% !TeX program = lualatex
\documentclass[11pt,leqno]{article}
\usepackage{amsmath,amssymb}
\usepackage{fontspec}
\usepackage{unicode-math}
\setmainfont{STIXTwoText}[Extension=.otf,UprightFont=*-Regular,ItalicFont=*-Italic,BoldFont=*-Bold,BoldItalicFont=*-BoldItalic]
\setmathfont{STIXTwoMath-Regular.otf}
\usepackage[main=english]{babel}
\usepackage{tikz-cd}
\usetikzlibrary{babel}
\usepackage[a4paper,inner=27mm,outer=23mm,top=25mm,bottom=27mm]{geometry}
\usepackage{xcolor}
\usepackage[hidelinks,unicode]{hyperref}
\setlength{\parindent}{0pt}
\setlength{\parskip}{0.55\baselineskip}
\setlength{\emergencystretch}{3em}
\newcommand{\Page}[1]{\ifvmode\noindent\else\unskip\ \fi{\footnotesize\textcolor{gray}{[p.\,#1]}}\ \ignorespaces}
\title{Hodge Theory, III}
\author{Pierre Deligne}
\date{}
\begin{document}
% ---- 005-012 ----
\maketitle
\Page{5}

\begin{center}
{\Large HODGE THEORY, III\par}
\medskip
{\large \textit{by} \textsc{Pierre} DELIGNE\par}
\end{center}

\section*{\centering\normalfont CONTENTS}

\begingroup
\small
\noindent\textsc{Introduction}\dotfill 6\par
\medskip
\noindent\textsc{Terminology and notation}\dotfill 7\par
\medskip
\noindent 5. \textbf{Cohomological descent}\dotfill 8\par
\noindent\hspace*{1.5em}1. Simplicial topological spaces\dotfill 8\par
\noindent\hspace*{1.5em}2. Cohomology of simplicial topological spaces\dotfill 12\par
\noindent\hspace*{1.5em}3. Cohomological descent\dotfill 14\par
\medskip
\noindent 6. \textbf{Examples of simplicial topological spaces}\dotfill 16\par
\noindent\hspace*{1.5em}1. Classifying spaces\dotfill 16\par
\noindent\hspace*{1.5em}2. Construction of hypercoverings\dotfill 19\par
\noindent\hspace*{1.5em}3. Relative cohomology\dotfill 21\par
\noindent\hspace*{1.5em}4. Multisimplicial spaces\dotfill 22\par
\medskip
\noindent 7. \textbf{Complements to § 1}\dotfill 23\par
\noindent\hspace*{1.5em}1. Filtered derived category\dotfill 23\par
\noindent\hspace*{1.5em}2. Complements to the lemma on two filtrations\dotfill 25\par
\medskip
\noindent 8. \textbf{Hodge theory of algebraic spaces}\dotfill 28\par
\noindent\hspace*{1.5em}1. Hodge complexes\dotfill 28\par
\noindent\hspace*{1.5em}2. Separated algebraic spaces\dotfill 38\par
\noindent\hspace*{1.5em}3. Hodge theory of simplicial schemes\dotfill 41\par
\medskip
\noindent 9. \textbf{Examples and applications}\dotfill 45\par
\noindent\hspace*{1.5em}1. Cohomology of groups and classifying spaces\dotfill 45\par
\noindent\hspace*{1.5em}2. Hodge theory of smooth hypersurfaces, following Griffiths\dotfill 48\par
\noindent\hspace*{1.5em}3. Construction of complexes of differential operators of the 1\textsuperscript{st} order\dotfill 50\par
\medskip
\noindent 10. \textbf{\boldmath Hodge theory at level \(\leq 1\)}\dotfill 53\par
\noindent\hspace*{1.5em}1. 1-motives\dotfill 53\par
\noindent\hspace*{1.5em}2. 1-motives and biextensions\dotfill 60\par
\noindent\hspace*{1.5em}3. Algebraic interpretation of mixed \(\mathrm{H}^{1}\) : the case of curves\dotfill 69\par
\noindent\hspace*{1.5em}4. Translation of a theorem of Picard\dotfill 75\par
\medskip
\noindent\textsc{Bibliography}\dotfill 77\par
\endgroup

\Page{6}

\section*{\centering Introduction}

This paper follows [4] and [5], cited as I and II. The numbering of its sections continues that of II (which contains §§ 1 to 4). To refer to a result of II, we simply give its number. For example : (3.2.17).

In I, we presented the yoga underlying II and the present paper III. However, II and III are logically independent of I. They do not contain all the results announced in I.

In II, we presented the Hodge theory of nonsingular algebraic varieties (not necessarily complete). Here we treat the case of arbitrary singularities. From § 7 onward, essential use will be made of the results of §§ 1 to 3 of II.

In the case of a \textit{complete} singular algebraic variety \(\mathrm{X}\), the basic idea is to use resolution of singularities to “replace” \(\mathrm{X}\) by a simplicial system of smooth projective schemes \(\mathrm{X}_{\bullet}\)
\[
\begin{tikzcd}[column sep=2.6em]
\ldotp\,\ldotp\,\ldotp\,\ldotp\,\ldotp\ \mathrm{X}_2
  \arrow[r, shift left=3.2ex]
  \arrow[r]
  \arrow[r, shift right=3.2ex]
&
\mathrm{X}_1
  \arrow[l, shift right=1.6ex, shorten <=1.1em]
  \arrow[l, shift left=1.6ex, shorten <=1.1em]
  \arrow[r, shift left=1.6ex]
  \arrow[r, shift right=1.6ex]
&
\mathrm{X}_0
  \arrow[l, shorten <=1.1em]
\end{tikzcd}
\]
(or, as we shall say, by a smooth projective simplicial scheme). The results of §§ 5 and 6 allow us to find such simplicial schemes \(\mathrm{X}_{\bullet}\) having, in a suitable sense, the same cohomology as \(\mathrm{X}\). We can then “express”, \textit{via} a spectral sequence, the cohomology of \(\mathrm{X}\) in terms of that of the \(\mathrm{X}_n\). In accordance with the general principles of I, we can then use classical Hodge theory on each of the \(\mathrm{X}_n\), and deduce from it a mixed Hodge structure on the cohomology of \(\mathrm{X}\).

In the case of an arbitrary algebraic variety \(\mathrm{X}\), we “replace” \(\mathrm{X}\) by a smooth simplicial scheme \(\mathrm{X}_{\bullet}\), compactified by a smooth projective simplicial scheme \(\overline{\mathrm{X}}_n\). We also arrange for \(\overline{\mathrm{X}}_n-\mathrm{X}_n\) to be a divisor with normal crossings, a union of smooth divisors \(\mathrm{D}_{n,i}\). We then “express” \textit{via} a spectral sequence the cohomology of \(\mathrm{X}\) in terms of the cohomology of the \(p\)-by-\(p\) intersections of the \(\mathrm{D}_{n,i}\) (\(n\geq 0,\ p\geq 0\)), and deduce a mixed Hodge structure on \(\mathrm{H}^{\bullet}(\mathrm{X})\).

§§ 5 and 6 constitute a summary (without proofs) of the theory of “cohomological descent”. This theory, in the more general setting of topoi, is expounded by B. Saint-Donat in \textit{S.G.A.}4, V \textit{bis}.

In § 7 n\textsuperscript{o} 1, we recast some of the results obtained earlier in the language of filtered derived categories. In n\textsuperscript{o} 2, we give a proof of the lemma on two filtrations that is simpler than the one in § 1, n\textsuperscript{o} 3.

The forbidding n\textsuperscript{o} 1 of § 8 is the heart of this work. The essential result is Theorem (8.1.15). In its proof, in order to apply the lemma on two
\Page{7}
filtrations, we use the fact that every morphism of mixed Hodge structures is \textit{strictly} compatible with the filtrations \(\mathrm{W}\) and \(\mathrm{F}\) (2.3.5). In n\textsuperscript{o} 2, we define the mixed Hodge structure on \(\mathrm{H}^n(\mathrm{X},\mathbf{Z})\) (\(\mathrm{X}\) a separated scheme). We show that the Hodge numbers \(h^{pq}\) of \(\mathrm{H}^n(\mathrm{X},\mathbf{Z})\) can be nonzero only for \((p,q)\in[0,n]\times[0,n]\). For \(n>\dim\mathrm{X}\), we have a more precise result (8.2.4). If \(\mathrm{X}\) is smooth (resp. complete), they can be nonzero only if, in addition, \(p+q\geq n\) (resp. \(p+q\leq n\)). If \(\mathrm{X}\) is a complete subscheme with complement \(\mathrm{U}\) of a complete smooth scheme \(\mathrm{Z}\) of dimension \(n\), we can, according to N. Katz, interpret this “duality” between the complete and smooth cases as arising from “Alexander duality”
\[
\ldots\ \mathrm{H}^i(\mathrm{Z},\mathbf{Q})
\to \mathrm{H}^i(\mathrm{X},\mathbf{Q})
\to \bigl(\mathrm{H}^{2n-i-1}(\mathrm{U},\mathbf{Q})\bigr)^*
\to \mathrm{H}^{i+1}(\mathrm{Z},\mathbf{Q})\ \ldots
\]

From the functoriality properties of the theory we have constructed, we deduce the useful corollaries (8.2.5) to (8.2.8), which are stated in terms independent of any Hodge theory.

In § 8, n\textsuperscript{o} 3, we endow the cohomology of every simplicial scheme \(\mathrm{X}_{\bullet}\) with a mixed Hodge structure. This generality is not illusory.

\textit{a)} Relative cohomology spaces can be interpreted as the cohomology of suitable simplicial schemes.

\textit{b)} Let \(\mathrm{B}_{\mathrm{G}}\) be the classifying space of the Lie group underlying an algebraic group \(\mathrm{G}\). The cohomology of \(\mathrm{B}_{\mathrm{G}}\) can be interpreted as that of a suitable simplicial scheme.

In § 9, n\textsuperscript{o} 1, after computing the mixed Hodge structure on the cohomology of \(\mathrm{B}_{\mathrm{G}}\) (for \(\mathrm{G}\) linear), we deduce that of \(\mathrm{G}\). As a corollary, we find that if a linear group \(\mathrm{G}\) acts on a complete nonsingular variety \(\mathrm{X}\), then the map
\[
\mathrm{H}^{\bullet}(\mathrm{X},\mathbf{Q})
\to \mathrm{H}^{\bullet}(\mathrm{X},\mathbf{Q})
\otimes \mathrm{H}^{\bullet}(\mathrm{G},\mathbf{Q})
\]
factors through \(\mathrm{H}^{\bullet}(\mathrm{X},\mathbf{Q})\otimes\mathrm{H}^{0}(\mathrm{G},\mathbf{Q})\subset\mathrm{H}^{\bullet}(\mathrm{X},\mathbf{Q})\otimes\mathrm{H}^{\bullet}(\mathrm{G},\mathbf{Q})\).

In § 10, we interpret in terms of abelian schemes the mixed Hodge structures purely of type \(\{(-1,-1),(-1,0),(0,-1),(0,0)\}\) (III.0.5), and we consider the \(\mathrm{H}^{1}\) of curves in detail.

The theory developed so far is an absolute theory (we do not study the functors \(\mathrm{R}f_*\) in it), and concerns only constant coefficients. I conjecture that if \(\mathscr{H}\) is a polarizable variation of Hodge structures (in the sense of Griffiths [6]) on a scheme \(\mathrm{X}\), then the cohomology of \(\mathrm{X}\) with coefficients in the local system \(\mathscr{H}\) is endowed with a natural mixed Hodge structure. I can prove this only when \(\mathrm{X}\) is complete.

\section*{\centering Terminology and notation}

\textbf{(III.0.1)} Let \(u:\mathrm{X}\to\mathrm{Y}\) be a continuous map between topological spaces. We shall say that \(u\) is \textit{proper} if it is proper in the sense of Bourbaki (i.e. universally closed) and is moreover separated (i.e. if the diagonal of \(\mathrm{X}\times_{\mathrm{Y}}\mathrm{X}\) is closed).

\Page{8}

\textbf{(III.0.2)} Following Gabriel and Zisman, we shall say \textit{simplicial} where one formerly said semi-simplicial.

\textbf{(III.0.3)} We shall denote by \(\mathrm{A}\) a noetherian subring of \(\mathbf{R}\) such that \(\mathrm{A}\otimes\mathbf{Q}\) is a field. The useful cases are \(\mathrm{A}=\mathbf{Z}\), \(\mathbf{Q}\) or \(\mathbf{R}\).

\textbf{(III.0.4)} An \(\mathrm{A}\)-\textit{mixed Hodge structure} consists of a finitely generated \(\mathrm{A}\)-module \(\mathrm{H}_{\mathrm{A}}\), a finite increasing filtration \(\mathrm{W}\) on the \(\mathrm{A}\otimes\mathbf{Q}\)-vector space \(\mathrm{H}_{\mathrm{A}\otimes\mathbf{Q}}=\mathrm{H}_{\mathrm{A}}\otimes\mathbf{Q}\), and a finite decreasing filtration \(\mathrm{F}\) on the \(\mathbf{C}\)-vector space \(\mathrm{H}_{\mathbf{C}}=\mathrm{H}_{\mathrm{A}}\otimes_{\mathrm{A}}\mathbf{C}\). We require that the \((\operatorname{Gr}^{\mathrm{W}}_n(\mathrm{H}_{\mathrm{A}\otimes\mathbf{Q}}),\operatorname{Gr}^{\mathrm{W}}_n(\mathrm{F}))\) be \(\mathrm{A}\otimes\mathbf{Q}\)-Hodge structures. For \(\mathrm{A}=\mathbf{Z}\) (resp. \(\mathbf{Q}\)) we recover (2.3.1) (resp. 2.3.8); the results of (2.3) carry over unchanged.

\textbf{(III.0.5)} Let \(\mathscr{E}\) be a subset of \(\mathbf{Z}\times\mathbf{Z}\). A mixed Hodge structure \(\mathrm{H}\) is of \textit{type} \(\mathscr{E}\) if its Hodge numbers \(h^{pq}\) vanish for \((p,q)\notin\mathscr{E}\).

\textbf{(III.0.6)} We shall henceforth denote by \(\Omega_{\mathrm{X}}^{\bullet}(\log\mathrm{D})\) what we denoted by \(\Omega_{\mathrm{X}}^*\langle\mathrm{D}\rangle\) (3.1.2).

\textbf{(III.0.7)} Unless explicitly stated otherwise, \textit{scheme} means “scheme of finite type over \(\mathbf{C}\)”, and a \textit{sheaf} on a scheme \(\mathrm{X}\) is a sheaf on the topological space underlying \(\mathrm{X}_{\mathrm{an}}\).

\section*{\centering 5. COHOMOLOGICAL DESCENT}

\section*{5.1. Simplicial topological spaces.}

This subsection begins with a review, for which we may refer to the Strasbourg homotopy seminar, 63/64.

\textbf{(5.1.1)} We shall use the following notation.

\noindent \(\mathscr{A}^{0}=\) opposite category of a category \(\mathscr{A}\).\par
\noindent \(\mathit{Hom}(\mathscr{A},\mathscr{B})=\) category of functors from \(\mathscr{A}\) to \(\mathscr{B}\).\par

Let \(n\) and \(k\) denote integers \(\geq-1\).

\noindent \(\Delta_n=\) the finite totally ordered set \([0,n]\).\par
\noindent \(\delta_i:\Delta_n\to\Delta_{n+1}=\) the order-preserving injection such that \(i\notin\delta_i(\Delta_n)\) (\(0\leq i\leq n+1\)).\par
\noindent \(s_i:\Delta_{n+1}\to\Delta_n=\) the order-preserving surjection such that \(s_i(i)=s_i(i+1)\) (\(0\leq i\leq n\)).\par
\noindent \(\varepsilon:\Delta_{-1}\to\Delta_n=\) the unique map from \(\Delta_{-1}\) to \(\Delta_n\).\par
\noindent \((\Delta^+)=\) the category whose objects are the \(\Delta_n\) (\(n\geq-1\)), and whose arrows are the order-preserving maps between the \(\Delta_n\).\par
\noindent \((\Delta)=\) the full subcategory of \((\Delta^+)\) with objects the \(\Delta_n\) (\(n\geq0\)).\par
\noindent \((\Delta^+)_k=\) the full subcategory of \((\Delta^+)\) with objects the \(\Delta_n\) (\(k\geq n\)).\par
\noindent \((\Delta)_k=\) the full subcategory of \((\Delta^+)\) with objects the \(\Delta_n\) (\(k\geq n\geq0\)).\par

\Page{9}

For any category \(\mathscr{C}\), a \textit{simplicial object} (resp. \textit{\(k\)-truncated simplicial object}) of \(\mathscr{C}\) is an object of \(\mathit{Hom}((\Delta)^0,\mathscr{C})\) (resp. \(\mathit{Hom}((\Delta)_k^0,\mathscr{C})\)). Similarly, a \textit{cosimplicial object} (resp. …) of \(\mathscr{C}\) is an object of \(\mathit{Hom}((\Delta),\mathscr{C})\) (resp. …).

The \textit{skeleton} functor \(\operatorname{sq}_k\) is, for \(k\geq-1\), the restriction functor
\[
\operatorname{sq}_k:\mathit{Hom}((\Delta)^0,\mathscr{C})
\to\mathit{Hom}((\Delta)_k^0,\mathscr{C}).
\]
The \textit{coskeleton} functor \(\operatorname{cosq}_k\) is the right adjoint to \(\operatorname{sq}_k\)
\[
\operatorname{cosq}_k:\mathit{Hom}((\Delta)_k^0,\mathscr{C})
\to\mathit{Hom}((\Delta)^0,\mathscr{C}).
\]
Let \(\mathrm{Y}\) be a simplicial object of \(\mathscr{C}\). We also use the name \textit{skeleton} for the functor
\[
\operatorname{sq}_k:\mathit{Hom}((\Delta)^0,\mathscr{C})/\mathrm{Y}
\to\mathit{Hom}((\Delta)_k^0,\mathscr{C})/\operatorname{sq}_k(\mathrm{Y}).
\]
Its right adjoint is denoted by \(\operatorname{cosq}_k^{\mathrm{Y}}\) (coskeleton relative to \(\mathrm{Y}\))
\[
\operatorname{cosq}_k^{\mathrm{Y}}:
\mathit{Hom}((\Delta)_k^0,\mathscr{C})/\operatorname{sq}_k\mathrm{Y}
\to\mathit{Hom}((\Delta)^0,\mathscr{C})/\mathrm{Y}.
\]

The coskeleton functors are defined if finite projective limits exist in \(\mathscr{C}\). The functors \(\operatorname{cosq}_k^{\mathrm{Y}}\) exist whenever fiber products exist. We have
\[
\operatorname{cosq}_k^{\mathrm{Y}}(\mathrm{X})
=\operatorname{cosq}_k(\mathrm{X})
\times_{\operatorname{cosq}_k\operatorname{sq}_k(\mathrm{Y})}\mathrm{Y}.
\]

\textbf{(5.1.2)} If \(\mathrm{X}_{\bullet}:(\Delta)^0\to\mathscr{C}\) is a simplicial object of \(\mathscr{C}\), we set \(\mathrm{X}_n=\mathrm{X}_{\bullet}(\Delta_n)\), \(\delta_i=\mathrm{X}_{\bullet}(\delta_i:\Delta_n\to\Delta_{n+1}):\mathrm{X}_{n+1}\to\mathrm{X}_n\) and \(s_i=\mathrm{X}_{\bullet}(s_i:\Delta_{n+1}\to\Delta_n):\mathrm{X}_n\to\mathrm{X}_{n+1}\).
\[
\begin{tikzcd}[column sep=2.8em]
\ldotp\,\ldotp\,\ldotp\,\ldotp\,\ldotp\ \mathrm{X}_2
  \arrow[r, shift left=4.4ex, "\delta_0"{near end}]
  \arrow[r, "\delta_1"{near end}]
  \arrow[r, shift right=4.4ex, "\delta_2"{near end}]
&
\mathrm{X}_1
  \arrow[l, shift right=2.2ex, shorten <=1.2em, "s_0"'{near end}]
  \arrow[l, shift left=2.2ex, shorten <=1.2em, "s_1"'{near end}]
  \arrow[r, shift left=2.2ex, "\delta_0"{near end}]
  \arrow[r, shift right=2.2ex, "\delta_1"{near end}]
&
\mathrm{X}_0
  \arrow[l, shorten <=1.2em, "s_0"'{near end}]
\end{tikzcd}
\tag*{\(\mathrm{X}_{\bullet}:\)}
\]

\textbf{(5.1.3)} Let \(\mathrm{S}\in\operatorname{Ob}\mathscr{C}\). The \textit{constant simplicial object} \(\mathrm{S}_{\bullet}\) is the simplicial object for which \(\mathrm{S}_n=\mathrm{S}\) and \(\delta_i=s_i=\operatorname{Id}_{\mathrm{S}}\). A simplicial (resp. \(k\)-truncated simplicial) object of \(\mathscr{C}\), \textit{augmented to} \(\mathrm{S}\), is a morphism \(a:\mathrm{X}_{\bullet}\to\mathrm{S}_{\bullet}\) (resp. \(a:\mathrm{X}_{\bullet}\to\operatorname{sq}_k(\mathrm{S}_{\bullet})\)). We shall identify the simplicial (resp. \(k\)-truncated simplicial) objects of \(\mathscr{C}\), augmented to \(\mathrm{S}\), \(a:\mathrm{X}_{\bullet}\to\mathrm{S}_{\bullet}\) (resp. …) with the objects \(\mathrm{X}_{\bullet}^{+}\) of \(\mathit{Hom}((\Delta^+)^0,\mathscr{C})\) (resp. \(\mathit{Hom}((\Delta^+)_k^0,\mathscr{C})\)), such that \(\mathrm{X}_{\bullet}^{+}(\Delta_{-1})=\mathrm{S}\). We have \(a_n=\mathrm{X}_{\bullet}^{+}(\varepsilon:\Delta_{-1}\to\Delta_n)\). These objects will be denoted by notation of the form \(a:\mathrm{X}_{\bullet}\to\mathrm{S}\). Relative coskeletons will be used mainly in this setting, and denoted by \(\operatorname{cosq}_k^{\mathrm{S}}\) or simply \(\operatorname{cosq}_k\).

\textbf{(5.1.4)} The coskeleton of the augmented 0-truncated simplicial object \(a_0:\mathrm{X}\to\mathrm{S}\) is the simplicial object augmented to \(\mathrm{S}\) whose components are the cartesian powers of \(\mathrm{X}\) in \(\mathscr{C}/\mathrm{S}\)
\[
\operatorname{cosq}_0(\mathrm{X}\to\mathrm{S})
=(((\mathrm{X}/\mathrm{S})^{\Delta_n})_{n\geq0}\to\mathrm{S}).
\]

\textbf{(5.1.5)} Let \(u:\mathrm{X}\to\mathrm{Y}\) be a continuous map between topological spaces, \(\mathrm{F}\) a sheaf on \(\mathrm{X}\), and \(\mathrm{G}\) a sheaf on \(\mathrm{Y}\). The set \(\operatorname{Hom}_u(\mathrm{G},\mathrm{F})\) of \textit{\(u\)-morphisms} from \(\mathrm{G}\) to \(\mathrm{F}\) is the set \(\operatorname{Hom}(u^*\mathrm{G},\mathrm{F})\sim\operatorname{Hom}(\mathrm{G},u_*\mathrm{F})\).

\Page{10}

Giving a \(u\)-morphism \(f\) from \(\mathrm{G}\) to \(\mathrm{F}\) amounts to giving, for every pair of open sets \(\mathrm{U}\subset\mathrm{X}\) and \(\mathrm{V}\subset\mathrm{Y}\) such that \(u(\mathrm{U})\subset\mathrm{V}\), a map \(f_{\mathrm{UV}}:\mathrm{G}(\mathrm{V})\to\mathrm{F}(\mathrm{U})\); these maps must satisfy the condition

\noindent (*) For \(\mathrm{U}'\subset\mathrm{U}\subset\mathrm{X}\), \(\mathrm{V}'\subset\mathrm{V}\subset\mathrm{Y}\), \(u(\mathrm{U})\subset\mathrm{V}\) and \(u(\mathrm{U}')\subset\mathrm{V}'\), the diagram
\[
\begin{tikzcd}[column sep=2em, row sep=3em]
\mathrm{G}(\mathrm{V}) \arrow[r] \arrow[d]
  & \mathrm{G}(\mathrm{V}') \arrow[d] \\
\mathrm{F}(\mathrm{U}) \arrow[r]
  & \mathrm{F}(\mathrm{U}')
\end{tikzcd}
\]
is commutative.

\textbf{(5.1.6)} A \textit{simplicial topological space} is a simplicial object of the category whose objects are topological spaces, and whose morphisms are continuous maps.

A \textit{sheaf} \(\mathrm{F}^{\bullet}\) on a simplicial topological space \(\mathrm{X}_{\bullet}\) consists of

\noindent\textit{a)} a family \(\mathrm{F}^n\) of sheaves on the \(\mathrm{X}_n\);\par
\noindent\textit{b)} for \(f:\Delta_n\to\Delta_m\), an \(\mathrm{X}_{\bullet}(f)\)-morphism \(\mathrm{F}^{\bullet}(f)\) from \(\mathrm{F}^n\) to \(\mathrm{F}^m\).\par

We also require that \(\mathrm{F}^{\bullet}(f\circ g)=\mathrm{F}^{\bullet}(f)\circ\mathrm{F}^{\bullet}(g)\).

A morphism \(u\) from \(\mathrm{F}^{\bullet}\) to \(\mathrm{G}^{\bullet}\) is a family of morphisms \(u^n:\mathrm{F}^n\to\mathrm{G}^n\) such that for \(f\in\operatorname{Hom}_{(\Delta)}(\Delta_n,\Delta_m)\), we have \(u^m\mathrm{F}^{\bullet}(f)=\mathrm{G}^{\bullet}(f)\,u^n\).

\textbf{(5.1.7)} For \(\mathrm{U}\) an open subset of \(\mathrm{X}_n\), let \(\mathrm{F}^{\bullet}(\mathrm{U})=\mathrm{F}^n(\mathrm{U})\). For \(f:\Delta_n\to\Delta_m\), \(\mathrm{U}\subset\mathrm{X}_n\) and \(\mathrm{V}\) in \(\mathrm{X}_m\) such that \(\mathrm{X}_{\bullet}(f)(\mathrm{V})\subset\mathrm{U}\), let \(\mathrm{F}^{\bullet}(f,\mathrm{V},\mathrm{U}):\mathrm{F}^{\bullet}(\mathrm{U})\to\mathrm{F}^{\bullet}(\mathrm{V})\) be the map induced by \(\mathrm{F}^{\bullet}(f)\). One checks immediately that the system of sets \(\mathrm{F}^{\bullet}(\mathrm{U})\) (indexed by \(n\) and \(\mathrm{U}\subset\mathrm{X}_n\)) and maps \(\mathrm{F}^{\bullet}(f,\mathrm{V},\mathrm{U})\) determines \(\mathrm{F}^{\bullet}\) uniquely (cf. (5.1.5)). For such a system to arise from a sheaf on \(\mathrm{X}_{\bullet}\), it is necessary and sufficient that :

\noindent\textit{a)} \(\mathrm{F}^{\bullet}(fg,\mathrm{U},\mathrm{W})=\mathrm{F}^{\bullet}(f,\mathrm{U},\mathrm{V})\,\mathrm{F}^{\bullet}(g,\mathrm{V},\mathrm{W})\) whenever this makes sense;\par
\noindent\textit{b)} for every \(n\), the \(\mathrm{F}^{\bullet}(\mathrm{U})\) for \(\mathrm{U}\subset\mathrm{X}_n\) form a sheaf on \(\mathrm{X}_n\).\par

\textbf{(5.1.8)} This shows that sheaves on \(\mathrm{X}_{\bullet}\) can be interpreted as sheaves on a suitable site, and in particular form a topos \((\mathrm{X}_{\bullet})^{\sim}\). For sheaves on \(\mathrm{X}_{\bullet}\), we shall freely use the terminology in use for sheaves on a site. Thus, a sheaf of abelian groups (resp. a sheaf of rings, …) \(\mathrm{F}^{\bullet}\) is a system of sheaves of abelian groups (resp. sheaves of rings, …) \(\mathrm{F}^n\) on the \(\mathrm{X}_n\), and morphisms \(\mathrm{F}^{\bullet}(f)\).

\textbf{(5.1.9)} \textit{Examples.} (I) Let \(\mathrm{X}_{\bullet}\) be a simplicial analytic space. The structure sheaves \(\mathscr{O}_{\mathrm{X}_n}\) form a sheaf of rings on \(\mathrm{X}_{\bullet}\).

(II) Let \(\mathrm{X}_{\bullet}\) be a simplicial analytic space augmented to \(\mathrm{S}\). The sheaves \(\Omega^1_{\mathrm{X}_n/\mathrm{S}}\) form a sheaf of \(\mathscr{O}\)-modules on \(\mathrm{X}_{\bullet}\). Its \(i\)-th exterior power is the sheaf
\Page{11}
of \(\mathscr{O}\)-modules \((\Omega^i_{\mathrm{X}_n/\mathrm{S}})_{n\geq0}\), denoted by \(\Omega^i_{\mathrm{X}_{\bullet}/\mathrm{S}}\). The de Rham complexes \((\Omega^{\bullet}_{\mathrm{X}_n/\mathrm{S}})_{n\geq0}\) form a complex of sheaves on \(\mathrm{X}_{\bullet}\).

(III) Let \(\mathrm{F}^{\bullet}\) be a sheaf of abelian groups on \(\mathrm{X}_{\bullet}\). The canonical flasque resolutions of Godement \(\mathscr{C}^{\bullet}(\mathrm{F}^n)\) form a complex of sheaves on \(\mathrm{X}_{\bullet}\), and a resolution of \(\mathrm{F}^{\bullet}\).

(IV) Let \(\mathrm{S}\) be a topological space. The sheaves \(\mathrm{F}^{\bullet}\) on the constant simplicial space \(\mathrm{S}_{\bullet}\) are identified with cosimplicial sheaves on \(\mathrm{S}\). In particular, an abelian sheaf \(\mathrm{F}^{\bullet}\) on \(\mathrm{S}_{\bullet}\) defines a differential complex \((\mathrm{F}^n,d=\sum\limits_i(-1)^i\delta_i)\). A complex of abelian sheaves \(\mathrm{K}\) on \(\mathrm{S}_{\bullet}\) defines a double complex \((\mathrm{K}^{n,m})\) again denoted by \(\mathrm{K}\) (\(m=\) cosimplicial degree) with associated simple complex \(\boldsymbol{s}\mathrm{K}\) :
\[
(\boldsymbol{s}\mathrm{K})^n
=\bigoplus_{p+q=n}\mathrm{K}^{pq}
\tag*{\textbf{(5.1.9.1)}}
\]
and
\[
d(x^{pq})=d_{\mathrm{K}}(x^{pq})
+(-1)^p\sum_i(-1)^i\delta_i x^{pq}.
\tag*{\textbf{(5.1.9.2)}}
\]
We shall denote by \(\mathrm{L}\) the second filtration on \(\boldsymbol{s}\mathrm{K}\) :
\[
\mathrm{L}^r(\boldsymbol{s}\mathrm{K})
=\bigoplus_{q\geq r}\mathrm{K}^{pq}.
\tag*{\textbf{(5.1.9.3)}}
\]

\textbf{(5.1.10)} Let \(u:\mathrm{X}_{\bullet}\to\mathrm{Y}_{\bullet}\) be a morphism of simplicial topological spaces, with components \(u_n:\mathrm{X}_n\to\mathrm{Y}_n\). If \(\mathrm{G}\) (resp. \(\mathrm{F}\)) is a sheaf on \(\mathrm{Y}_{\bullet}\) (resp. on \(\mathrm{X}_{\bullet}\)), then \((u_n^*\mathrm{G}^n)_{n\geq0}\) (resp. \((u_{n*}\mathrm{F}^n)_{n\geq0}\)) is a sheaf on \(\mathrm{X}_{\bullet}\) (resp. on \(\mathrm{Y}_{\bullet}\)); we denote it by \(u^*\mathrm{G}\) (resp. \(u_*\mathrm{F}\)). The functors \(u_*\) and \(u^*\) are adjoint; they are the direct and inverse image morphisms of a morphism of topoi \(u:(\mathrm{X}_{\bullet})^{\sim}\to(\mathrm{Y}_{\bullet})^{\sim}\).

\textbf{(5.1.11)} Let \(a:\mathrm{X}_{\bullet}\to\mathrm{S}\) be a simplicial topological space augmented to \(\mathrm{S}\). If \(\mathrm{F}\) is a sheaf on \(\mathrm{S}\), \(a^*\mathrm{F}=(a_n^*\mathrm{F})_{n\geq0}\) “is” a sheaf on \(\mathrm{X}_{\bullet}\). The functor \(a^*\) has a right adjoint
\[
a_*:\mathrm{F}^{\bullet}\mapsto
\operatorname{Ker}\bigl(
\begin{tikzcd}[ampersand replacement=\&, cramped, column sep=1.7em, baseline=-0.5ex]
a_{0*}\mathrm{F}^0
  \arrow[r, shift left=0.5ex, "\delta_0^*"]
  \arrow[r, shift right=0.5ex, "\delta_1^*"']
\& a_{1*}\mathrm{F}^1
\end{tikzcd}
\bigr).
\tag*{\textbf{(5.1.11.1)}}
\]
The functors \(a_*\) and \(a^*\) are the direct and inverse image functors of a morphism of topoi \(a:(\mathrm{X}_{\bullet})^{\sim}\to(\mathrm{S})^{\sim}\).

\textbf{(5.1.12)} Let \(\mathrm{S}_{\bullet}\) be the constant simplicial space associated with \(\mathrm{S}\) and \(a_{\bullet}:\mathrm{X}_{\bullet}\to\mathrm{S}_{\bullet}\) the morphism defined by \(a\). For \(\mathrm{F}^{\bullet}\) an abelian sheaf on \(\mathrm{X}_{\bullet}\), we identify \(a_{\bullet *}(\mathrm{F}^{\bullet})\) with a cosimplicial sheaf on \(\mathrm{S}\) ((5.1.9) (IV)). For \(\mathrm{K}\) a complex of abelian sheaves on \(\mathrm{X}_{\bullet}\), if \(\mathrm{K}^{pq}\) is the degree \(p\) component of the restriction of \(\mathrm{K}\) to \(\mathrm{X}_q\), the components of the simple complex associated with the double complex defined by \(a_{\bullet *}\mathrm{K}\) are
\[
(\boldsymbol{s}a_{\bullet *}\mathrm{K})^n
=\bigoplus_{p+q=n}a_{q*}\mathrm{K}^{pq}.
\tag*{\textbf{(5.1.12.1)}}
\]
\Page{12}
The spectral sequence defined by the filtration \(\mathrm{L}\) ((5.1.9.3)) on \(\boldsymbol{s}a_{\bullet *}\mathrm{K}\) is written
\[
\mathrm{E}_1^{pq}
=H^q(a_{p*}(\mathrm{K}\mid\mathrm{X}^{p}))
\Rightarrow H^{p+q}(\boldsymbol{s}a_{\bullet *}\mathrm{K}).
\tag*{\textbf{(5.1.12.2)}}
\]

\textbf{(5.1.13)} Set \(\mathrm{S}=\mathrm{P}^{t}\) (a topological space consisting of a single point). We find that, for \(\mathrm{F}^{\bullet}\) a sheaf on \(\mathrm{X}_{\bullet}\), the \(\Gamma(\mathrm{X}_n,\mathrm{F}^n)\) are the components of a cosimplicial set \(\Gamma^{\bullet}(\mathrm{X}_{\bullet},\mathrm{F}^{\bullet})\). The functor \(\Gamma\) (global sections) is the functor
\[
\Gamma:\mathrm{F}^{\bullet}\mapsto
\operatorname{Ker}\bigl(
\begin{tikzcd}[ampersand replacement=\&, cramped, column sep=1.4em, baseline=-0.5ex]
\Gamma(\mathrm{X}_0,\mathrm{F}^0)
  \arrow[r, shift left=0.5ex]
  \arrow[r, shift right=0.5ex]
\& \Gamma(\mathrm{X}_1,\mathrm{F}^1)
\end{tikzcd}
\bigr).
\tag*{\textbf{(5.1.13.1)}}
\]
If \(\mathrm{K}\) is a complex of abelian sheaves, we denote by \(\Gamma^{\bullet}(\mathrm{X}_{\bullet},\mathrm{K})\) the differential complex whose components are the cosimplicial abelian groups \(\Gamma^{\bullet}(\mathrm{X}_{\bullet},\mathrm{K}^n)\) and by \(\boldsymbol{s}\Gamma^{\bullet}(\mathrm{X}_{\bullet},\mathrm{K})\) the associated simple complex.

\section*{5.2. Cohomology of simplicial topological spaces.}

\textbf{(5.2.1)} To each simplicial topological space \(\mathrm{X}_{\bullet}\) is attached an ordinary topological space \(\lvert\mathrm{X}_{\bullet}\rvert\), called its geometric realization (see below). In the special case where the \(\mathrm{X}_n\) are discrete, we recover the usual notion of the geometric realization of a simplicial set.

In [14], G. Segal defines the cohomology of \(\mathrm{X}_{\bullet}\) with values in an abelian group \(\mathrm{A}\) to be \(\mathrm{H}^{\bullet}(\lvert\mathrm{X}_{\bullet}\rvert,\mathrm{A})\). Under suitable hypotheses, the filtration of \(\lvert\mathrm{X}_{\bullet}\rvert\) by successive skeleta yields a spectral sequence.
\[
\mathrm{E}_1^{pq}
=\mathrm{H}^q(\mathrm{X}_p,\mathrm{A})
\Rightarrow\mathrm{H}^{p+q}(\mathrm{X}_{\bullet},\mathrm{A}).
\tag*{\textbf{(5.2.1.1)}}
\]

One of the advantages of this definition is that it applies equally well, for example, to the definition of groups \(\mathrm{K}(\mathrm{X}_{\bullet})\). We shall adopt another, better suited to sheaf-theoretic techniques.

\textit{Geometric realization.} — For \(n\) an integer \(\geq0\), we denote by \(\lvert\Delta_n\rvert\) the simplex in \(\mathbf{R}^{\Delta_n}\) with the set of basis vectors as its vertices. We identify \(\Delta_n\) with the set of vertices of \(\lvert\Delta_n\rvert\). Every function \(f:\Delta_n\to\Delta_m\) extends by linearity to \(\lvert f\rvert:\lvert\Delta_n\rvert\to\lvert\Delta_m\rvert\).

Let \(\mathrm{X}_{\bullet}\) be a simplicial topological space. Let \(\mathrm{Y}=\coprod\limits_{n\geq0}(\mathrm{X}_n\times\lvert\Delta_n\rvert)\). Let \(\mathrm{R}\) be the finest equivalence relation on \(\mathrm{Y}\) such that, for all \(f:\Delta_n\to\Delta_m\) in \((\Delta)\), \(x_m\in\mathrm{X}_m\) and \(a\in\lvert\Delta_n\rvert\), we have
\[
(x_m,\lvert f\rvert(a))\equiv(f(x_m),a)\qquad(\operatorname{mod}\mathrm{R}).
\]
The \textit{geometric realization} of \(\mathrm{X}_{\bullet}\) is by definition
\[
\lvert\mathrm{X}_{\bullet}\rvert=\mathrm{Y}/\mathrm{R}.
\]

\textit{Definition} \textbf{(5.2.2)}. — \textit{Let} \(\mathrm{X}_{\bullet}\) \textit{be a simplicial topological space. The functors} cohomology with coefficients in the abelian sheaf \(\mathrm{F}^{\bullet}\) on \(\mathrm{X}_{\bullet}\), \(\mathrm{H}^i(\mathrm{X}_{\bullet},\mathrm{F}^{\bullet})\), \textit{are the derived functors of the functor} \(\Gamma\) (5.1.13.1).

This definition is equivalent to the following one, which is sometimes easier to use.

% ---- 013-020 ----
\Page{13}
\textbf{(5.2.3)} Let \(\mathrm{F}^{\bullet}\) be an abelian sheaf on \(\mathrm{X}_{\bullet}\). One shows (for example using (5.1.9) (III)) that \(\mathrm{F}^{\bullet}\) always admits right resolutions \(\mathrm{K}\) such that
\[
\mathrm{H}^{r}(\mathrm{X}_{q},\mathrm{K}^{p,q})=0
\qquad \text{for}\quad p,q\geq 0
\quad\text{and}\quad r>0.
\]
If \(\mathrm{K}\) is such a resolution, we have canonically
\[
\mathrm{H}^{n}(\mathrm{X}_{\bullet},\mathrm{F}^{\bullet})
\simeq
\mathrm{H}^{n}(\boldsymbol{s}\Gamma^{\bullet}(\mathrm{X}_{\bullet},\mathrm{K})).
\tag*{\textbf{(5.2.3.1)}}
\]
It is easy to verify directly that the right-hand side is independent (up to unique isomorphism) of the choice of \(\mathrm{K}\); for what follows, one could \textit{define} \(\mathrm{H}^{\bullet}(\mathrm{X}_{\bullet},\mathrm{F}^{\bullet})\) by (5.2.3.1). One also verifies that the spectral sequence of the complex \(\boldsymbol{s}\Gamma^{\bullet}(\mathrm{X}_{\bullet},\mathrm{K})\), filtered by \(\mathrm{L}\) ((5.1.9.3) and (5.1.12.2) for \(\mathrm{S}=\mathrm{P}^{t}\))
\[
\mathrm{E}_{1}^{pq}=\mathrm{H}^{q}(\mathrm{X}_{p},\mathrm{F}^{p})
\Rightarrow
\mathrm{H}^{p+q}(\mathrm{X}_{\bullet},\mathrm{F}^{\bullet})
\tag*{\textbf{(5.2.3.2)}}
\]
is independent (up to unique isomorphism) of the choice of \(\mathrm{K}\) (cf. (5.2.1.1)).

\textbf{(5.2.4)} It is appropriate to make the preceding construction more precise by passing to derived categories, and to give a relative variant of it.

Let \(a:\mathrm{X}_{\bullet}\to\mathrm{S}\), let \(\mathrm{S}_{\bullet}\) be the constant simplicial space defined by \(\mathrm{S}\), let \(a_{\bullet}:\mathrm{X}_{\bullet}\to\mathrm{S}_{\bullet}\) be induced by \(a\) and let \(\varepsilon:\mathrm{S}_{\bullet}\to\mathrm{S}\) be the augmentation morphism. We have
\[
a_*=\varepsilon_*a_{\bullet *}:
(\mathrm{X}_{\bullet})^{\sim}\to(\mathrm{S})^{\sim}.
\tag*{\textbf{(5.2.4.1)}}
\]

Passing to derived functors in this formula gives
\[
\mathrm{R}a_*=\mathrm{R}\varepsilon_*\,\mathrm{R}a_{\bullet *}:
\mathrm{D}^{+}(\mathrm{X}_{\bullet})\to\mathrm{D}^{+}(\mathrm{S})
\tag*{\textbf{(5.2.4.2)}}
\]
(\(\mathrm{D}^{+}(\mathrm{X}_{\bullet})=\)the bounded-below derived category of the category of abelian sheaves on \(\mathrm{X}_{\bullet}\)). Let us compute \(\mathrm{R}a_{\bullet *}\) and \(\mathrm{R}\varepsilon_*\).

\textbf{(5.2.5)} Let \(u:\mathrm{X}_{\bullet}\to\mathrm{Y}_{\bullet}\) be a morphism of simplicial topological spaces. The functor \(\mathrm{R}u_*:\mathrm{D}^{+}(\mathrm{X}_{\bullet})\to\mathrm{D}^{+}(\mathrm{Y}_{\bullet})\) is computed “component by component” : if \(\mathrm{K}\) is a complex of abelian sheaves on \(\mathrm{X}_{\bullet}\), to compute \(\mathrm{R}u_*\mathrm{K}\), one takes a resolution \(\mathrm{K}\xrightarrow{\approx}\mathrm{K}'\) such that the components \(\mathrm{F}\) of \(\mathrm{K}'\) satisfy \(\mathrm{R}^{i}_{u_{p*}}(\mathrm{F}^{p})=0\) for \(i>0\), \(p\geq 0\) (cf. (5.1.9) (III)); one then has \(\mathrm{R}u_*\mathrm{K}\approx u_*\mathrm{K}'\).

\textbf{(5.2.6)} The functor \(\boldsymbol{s}:\mathrm{C}^{+}(\mathrm{S}_{\bullet})\to\mathrm{C}^{+}(\mathrm{S})\) takes acyclic complexes to acyclic complexes. It therefore trivially induces the derived functor
\[
\boldsymbol{s}:\mathrm{D}^{+}(\mathrm{S}_{\bullet})\to\mathrm{D}^{+}(\mathrm{S}).
\]

For \(\mathrm{F}\) an injective sheaf on \(\mathrm{S}_{\bullet}\), the differential complex \(\boldsymbol{s}\mathrm{F}\) is a resolution of \(\varepsilon_*\mathrm{F}\). This yields an isomorphism \(\mathrm{R}\varepsilon_*\xrightarrow{\sim}\boldsymbol{s}\) and (5.2.4.2) becomes
\[
\mathrm{R}a_*=\boldsymbol{s}\mathrm{R}a_{\bullet *}.
\tag*{\textbf{(5.2.6.1)}}
\]

\Page{14}
\textbf{(5.2.7)} Combined with (5.2.5) and specialized to the case \(\mathrm{S}=\mathrm{P}^{t}\), this formula proves (5.2.3.1). In concrete terms, it means that to compute \(\mathrm{R}a_*\mathrm{K}\), one can proceed in two steps :

\textit{a)} one takes a resolution \(\mathrm{K}\xrightarrow{\approx}\mathrm{K}'\), such that the components \(\mathrm{F}\) of \(\mathrm{K}'\) satisfy \(\mathrm{R}^{i}a_{p*}(\mathrm{F}^{p})=0\) for \(i>0\). The complex \(a_{\bullet *}\mathrm{K}'\in\mathrm{D}^{+}(\mathrm{S}_{\bullet})\) (derived category of the category of cosimplicial abelian sheaves on \(\mathrm{S}\)) is identified with \(\mathrm{R}a_{\bullet *}\mathrm{K}\).

\textit{b)} \(\mathrm{R}a_*\mathrm{K}\) is the simple complex \(\boldsymbol{s}a_{\bullet *}\mathrm{K}'\) associated with the double complex \(a_{\bullet *}\mathrm{K}'\).

The spectral sequence (5.2.3.2) generalizes to a spectral sequence deduced from (5.1.12.2)
\[
\mathrm{E}_{1}^{pq}=\mathrm{R}^{q}a_{p*}(\mathrm{K}\mid\mathrm{X}_{p})
\Rightarrow \mathrm{R}^{p+q}a_*\mathrm{K}.
\tag*{\textbf{(5.2.7.1)}}
\]

\section*{5.3. Cohomological descent.}

\textbf{(5.3.1)} Let \(a:\mathrm{X}_{\bullet}\to\mathrm{S}\) be an augmented simplicial topological space. For every sheaf \(\mathrm{F}\) on \(\mathrm{S}\), there is an adjunction morphism
\[
\varphi:\mathrm{F}\to a_*a^*\mathrm{F}.
\]
Deriving this morphism gives a morphism of functors from \(\mathrm{D}^{+}(\mathrm{S})\) to \(\mathrm{D}^{+}(\mathrm{S})\) :
\[
\varphi:\mathrm{Id}\to\mathrm{R}a_*a^*.
\tag*{\textbf{(5.3.1.1)}}
\]

\textit{Definition} \textbf{(5.3.2)}. — \textit{We say that \(a\) is of \textup{cohomological descent} if for every abelian sheaf \(\mathrm{F}\) on \(\mathrm{S}\), we have}
\[
\mathrm{F}\xrightarrow{\sim}
\operatorname{Ker}(a_{0*}a_{0}^{*}\mathrm{F}\to a_{1*}a_{1}^{*}\mathrm{F})
\]
\[
\makebox[\linewidth]{\textit{and}\hfill
\(\mathrm{R}^{i}a_*a^*\mathrm{F}=0
\qquad\text{\textit{for}}\quad i>0.\)\hfill}
\]

Equivalently, one may require that (5.3.1.1) be an isomorphism.

\textbf{(5.3.3)} If \(a\) is of cohomological descent, then, for \(\mathrm{K}\in\operatorname{Ob}\mathrm{D}^{+}(\mathrm{S})\), the canonical map
\[
\mathrm{R}\Gamma(\mathrm{S},\mathrm{K})
\to\mathrm{R}\Gamma(\mathrm{S},\mathrm{R}a_*a^*\mathrm{K})
\sim\mathrm{R}\Gamma(\mathrm{X}_{\bullet},a^*\mathrm{K})
\tag*{\textbf{(5.3.3.1)}}
\]
is an isomorphism. In particular, for \(\mathrm{F}\) an abelian sheaf on \(\mathrm{S}\), we have a spectral sequence (5.2.3.2)
\[
\mathrm{E}_{1}^{pq}=\mathrm{H}^{q}(\mathrm{X}_{p},a_{p}^{*}\mathrm{F})
\Rightarrow\mathrm{H}^{p+q}(\mathrm{S},\mathrm{F}).
\tag*{\textbf{(5.3.3.2)}}
\]

For a complex, there is again, in hypercohomology, a spectral sequence
\[
\mathrm{E}_{1}^{pq}=\mathbf{H}^{q}(\mathrm{X}_{p},a_{p}^{*}\mathrm{K})
\Rightarrow\mathbf{H}^{p+q}(\mathrm{S},\mathrm{K}).
\tag*{\textbf{(5.3.3.3)}}
\]

In both cases, the \(\mathrm{E}_{1}^{\bullet q}\) (\(q\) fixed) form a simplicial group, and
\[
d_{1}=\sum_i(-1)^i\delta_i:
\mathrm{E}_{1}^{p,q}\to\mathrm{E}_{1}^{p+1,q}.
\]

\Page{15}
\textit{Definition} \textbf{(5.3.4)}. — \textit{A continuous map \(a:\mathrm{X}\to\mathrm{S}\) is of \textup{cohomological descent} if the augmentation morphism of \(\operatorname{cosq}(\mathrm{X}\to\mathrm{S})\) :}
\[
((\mathrm{X}/\mathrm{S})^{\Delta_n})_{n\geq 0}\to\mathrm{S}
\]
\textit{is of cohomological descent. We say that \(a\) is of \textup{universal cohomological descent} if for every \(u:\mathrm{S}'\to\mathrm{S}\), the continuous map \(a':\mathrm{X}\times_{\mathrm{S}}\mathrm{S}'\to\mathrm{S}'\) is of cohomological descent.}

\textbf{(5.3.5)} The fundamental results, proved in \textit{S.G.A.}4, V \textit{bis}, are the following.

(I) The continuous maps of universal cohomological descent form a Grothendieck topology on the category of topological spaces.

It is called the \textit{topology of universal cohomological descent}.

(II) A proper (III.0.1) and surjective map is of universal cohomological descent.

(III) A map \(a:\mathrm{X}\to\mathrm{S}\) admitting sections locally on \(\mathrm{S}\) is of universal cohomological descent.

(IV) Let \(a:\mathrm{X}_{\bullet}\to\mathrm{S}\) be a \(k\)-truncated augmented simplicial space \((-1\leq k\leq\infty)\). For \(k\geq n\geq-1\), let \(\varphi_n:\operatorname{cosq}\mathrm{X}_{\bullet}\to\operatorname{cosq}\operatorname{sq}_n\mathrm{X}_{\bullet}\) be the obvious arrow. We say that \(\mathrm{X}_{\bullet}\) is a \textit{\(k\)-truncated hypercovering of \(\mathrm{S}\), for the topology of universal cohomological descent}, if the maps
\[
(\varphi_n)_{n+1}:\mathrm{X}_{n+1}
\to(\operatorname{cosq}\operatorname{sq}_n\mathrm{X}_{\bullet})_{n+1}
\qquad(-1\leq n\leq k-1)
\tag*{\textbf{(5.3.5.1)}}
\]
are of universal cohomological descent. If \(\mathrm{X}_{\bullet}\) is such a hypercovering, then the simplicial space augmented over \(\mathrm{S}\), \(\operatorname{cosq}(\mathrm{X}_{\bullet})\), is of cohomological descent.

(V) Let \(a\) be a morphism of simplicial topological spaces augmented over \(\mathrm{S}\)
\[
\begin{tikzcd}[column sep=2.5em,row sep=2.5em]
\mathrm{X}_{\bullet}\arrow[r,"a"]\arrow[d,"x"]
  & \mathrm{Y}_{\bullet}\arrow[d,"y"]\\
\mathrm{S}\arrow[r,equal] & \mathrm{S}
\end{tikzcd}
\]
We say that \(a\) is a \textit{hypercovering} for the topology of universal cohomological descent if the obvious maps \(\mathrm{X}_n\to(\operatorname{cosq}_{n-1}^{\mathrm{Y}_{\bullet}}\operatorname{sq}_{n-1}\mathrm{X}_{\bullet})_n\) are of universal cohomological descent. If \(a\) is such a hypercovering, then, for every \(\mathrm{K}\in\operatorname{Ob}\mathrm{D}^{+}(\mathrm{S})\)
\[
a^*:\mathrm{R}y_*y^*\mathrm{K}\xrightarrow{\sim}\mathrm{R}x_*x^*\mathrm{K}.
\]

\textbf{(5.3.6)} For \(k=\infty\), (IV) asserts that \(a:\mathrm{X}_{\bullet}\to\mathrm{S}\) is of cohomological descent if the \((\varphi_n)_{n+1}\) are of universal cohomological descent. For \(n=-1\) or \(0\), these maps are
\[
\begin{aligned}
n=-1&:\quad \mathrm{X}_0\xrightarrow{a}\mathrm{S}\\
n=0&:\quad \mathrm{X}_1\xrightarrow{(\delta_0,\delta_1)}
\mathrm{X}_0\times_{\mathrm{S}}\mathrm{X}_0.
\end{aligned}
\]
\Page{16}
For \(n=1\), \((\operatorname{cosq}\operatorname{sq}_1(\mathrm{X}_{\bullet}))_1\) is the subspace of \(\mathrm{X}_1\times_{\mathrm{S}}\mathrm{X}_1\times_{\mathrm{S}}\mathrm{X}_1\) consisting of the triples \((x,y,z)\) such that \(\delta_0x=\delta_0y\), \(\delta_1x=\delta_0z\) and \(\delta_1y=\delta_1z\). The map \((\varphi_1)_2\) is \(x\mapsto(\delta_0x,\delta_1x,\delta_2x)\).

For \(k=0\), (IV) is the definition (5.3.4).

\textit{Example} \textbf{(5.3.7)}. — Let \(\mathcal{U}=(\mathrm{U}_i)_{i\in\mathrm{I}}\) be an open covering, or a locally finite closed covering of \(\mathrm{S}\). Let \(\mathrm{X}=\coprod\limits_{i\in\mathrm{I}}\mathrm{U}_i\). Then, \(a:\mathrm{X}\to\mathrm{S}\) is of cohomological descent. The spectral sequence (5.3.3.2) for \(\mathrm{X}_{\bullet}=\operatorname{cosq}(\mathrm{X}\to\mathrm{S})\) is none other than the Leray spectral sequence of the covering \(\mathcal{U}\).

\textbf{(5.3.8)} Let \(a:\mathrm{X}_{\bullet}\to\mathrm{S}\) be as in (5.3.5) (IV). We say that \(\mathrm{X}_{\bullet}\) is a \textit{\(k\)-truncated proper hypercovering} of \(\mathrm{S}\) if the arrows (5.3.5.1) are proper and surjective. For \(k=\infty\), we speak simply of a proper hypercovering.

\section*{6. EXAMPLES OF SIMPLICIAL TOPOLOGICAL SPACES}

\section*{6.1. Classifying spaces.}

\textbf{(6.1.1)} Let \(u:\mathrm{X}\to\mathrm{S}\) be a continuous map. For every sheaf \(\mathrm{F}\) on \(\mathrm{S}\), the sheaf \(u^*\mathrm{F}\) is equipped with a “descent datum” relative to \(u\) : there is an isomorphism between the two inverse images of \(u^*\mathrm{F}\) on \(\mathrm{X}\times_{\mathrm{S}}\mathrm{X}\), and this isomorphism satisfies a cocycle condition. If, locally on \(\mathrm{S}\), \(u\) admits a section, this construction defines an equivalence between the category of sheaves on \(\mathrm{S}\) and that of sheaves on \(\mathrm{X}\), equipped with a descent datum.

Take for \(\mathrm{X}\) a (left) principal homogeneous space with group \(\mathrm{G}\) over a space \(\mathrm{S}\). The \(\mathrm{G}\)-equivariant sheaves on \(\mathrm{X}\) are precisely the sheaves equipped with a descent datum : every equivariant sheaf on \(\mathrm{X}\) is in one and only one way (as an equivariant sheaf) the inverse image of a sheaf on \(\mathrm{S}=\mathrm{X}/\mathrm{G}\).

\textbf{(6.1.2)} If a topological group \(\mathrm{G}\) acts on a space \(\mathrm{X}\), then \(\mathrm{G}\) acts on \(\mathrm{G}^{\Delta_n}\times\mathrm{X}\) by
\[
g\mathbin{.}(g_0,\ldots,g_n,x)
=(g_0g^{-1},\ldots,g_ng^{-1},gx).
\]
We denote by \([\mathrm{X}/\mathrm{G}]_{\bullet}\) the simplicial space
\[
[\mathrm{X}/\mathrm{G}]_{\bullet}
=((\mathrm{G}^{\Delta_n}\times\mathrm{X})/\mathrm{G})_{n\geq 0}.
\tag*{\textbf{(6.1.2.1)}}
\]

\textit{a)} If \(\mathrm{X}\) is a principal homogeneous space with group \(\mathrm{G}\) over \(\mathrm{S}=\mathrm{X}/\mathrm{G}\), then the map
\[
\mathrm{G}^{\Delta_n}\times\mathrm{X}\to\mathrm{X}^{\Delta_n}:
(g_0,\ldots,g_n,x)\mapsto(g_0x,\ldots,g_nx)
\]
identifies \([\mathrm{X}/\mathrm{G}]_n\) with the iterated fiber product \((\mathrm{X}/\mathrm{S})^{\Delta_n}\) :
\[
\operatorname{cosq}(\mathrm{X}\to\mathrm{S})
=([\mathrm{X}/\mathrm{G}]_{\bullet}\to\mathrm{S}).
\]
\Page{17}
In particular, we have (5.3.5 (III))
\[
\mathrm{H}^{\bullet}([\mathrm{X}/\mathrm{G}]_{\bullet})
\simeq\mathrm{H}^{\bullet}(\mathrm{X}/\mathrm{G})
\qquad\text{(for a principal homogeneous space).}
\tag*{\textbf{(6.1.2.2)}}
\]

\textit{b)} For every \(n\), \(\mathrm{G}^{\Delta_n}\times\mathrm{X}\) is a principal homogeneous space with group \(\mathrm{G}\) over \([\mathrm{X}/\mathrm{G}]_n\). For every equivariant sheaf \(\mathrm{F}\) on \(\mathrm{X}\), \(\operatorname{pr}_2^*\mathrm{F}\) is an equivariant sheaf on \(\mathrm{G}^{\Delta_n}\times\mathrm{X}\); by (6.1.1), the latter is the inverse image of \(\mathrm{F}^n\) on \([\mathrm{X}/\mathrm{G}]_n\). Every equivariant sheaf \(\mathrm{F}\) on \(\mathrm{X}\) thus defines a sheaf on \([\mathrm{X}/\mathrm{G}]_{\bullet}\). It is easy to verify that one thus obtains an equivalence of the category of equivariant sheaves on \(\mathrm{X}\) with the category of sheaves \(\mathrm{F}^{\bullet}\) on \([\mathrm{X}/\mathrm{G}]_{\bullet}\) satisfying

\((\ast)\) For every \(f:\Delta_n\to\Delta_m\), the structural morphism \(f^*\mathrm{F}^n\to\mathrm{F}^m\) is an isomorphism.

\textit{c)} The construction \textit{b)} is natural in \((\mathrm{G},\mathrm{X},\mathrm{F})\). We set
\[
\mathrm{H}^{\bullet}(\mathrm{X},\mathrm{G};\mathrm{F})
=\mathrm{H}^{\bullet}([\mathrm{X}/\mathrm{G}]_{\bullet},\mathrm{F}^{\bullet})
\tag*{\textbf{(6.1.2.3)}}
\]
(mixed cohomology of \(\mathrm{X}\), \(\mathrm{G}\) with coefficients in \(\mathrm{F}\)). Under the hypotheses of \textit{a)}, if \(\mathrm{F}\) is the inverse image of \(\mathrm{F}^{-1}\) on \(\mathrm{S}=\mathrm{X}/\mathrm{G}\), we have
\[
\mathrm{H}^{\bullet}(\mathrm{X},\mathrm{G};\mathrm{F})
=\mathrm{H}^{\bullet}(\mathrm{X}/\mathrm{G},\mathrm{F}^{-1})
\qquad\text{(for a principal homogeneous space).}
\tag*{\textbf{(6.1.2.4)}}
\]
This generalizes (6.1.2.2) (which is the case \(\mathrm{F}=\mathbf{Z}\)).

\textbf{(6.1.3)} Let \(\mathrm{P}^{t}\) be the topological space consisting of a single point. The \textit{simplicial classifying space} of \(\mathrm{G}\), denoted by \(\mathrm{B}_{\bullet\mathrm{G}}\), is the simplicial space
\[
\mathrm{B}_{\bullet\mathrm{G}}=[\mathrm{P}^{t}/\mathrm{G}]_{\bullet}.
\]

Let \(\mathrm{P}\) be a principal homogeneous space with group \(\mathrm{G}\) over \(\mathrm{S}\). The obvious morphism
\[
\operatorname{cosq}(\mathrm{P}\to\mathrm{S})
=[\mathrm{P}/\mathrm{G}]_{\bullet}
\to[\mathrm{P}^{t}/\mathrm{G}]_{\bullet}
=\mathrm{B}_{\bullet\mathrm{G}}
\]
defines a composite morphism
\[
[\mathrm{P}]:
\mathrm{H}^{\bullet}(\mathrm{B}_{\bullet\mathrm{G}})
\longrightarrow
\mathrm{H}^{\bullet}([\mathrm{P}/\mathrm{G}]_{\bullet})
\xleftarrow[\text{\textbf{6.1.2.2}}]{\sim}
\mathrm{H}^{\bullet}(\mathrm{S}).
\tag*{\textbf{(6.1.3.1)}}
\]

We shall see below that in good cases \(\mathrm{H}^{\bullet}(\mathrm{B}_{\bullet\mathrm{G}})=\mathrm{H}^{\bullet}(\mathrm{B}_{\mathrm{G}})\), and that the image of \([\mathrm{P}]\) consists of the \textit{characteristic classes} of \(\mathrm{P}\).

\textbf{(6.1.4)} Let \(\mathrm{G}\) be a Lie group, \(\mathrm{B}_{\mathrm{G}}\) a classifying space for \(\mathrm{G}\) and \(a:\mathrm{U}_{\mathrm{G}}\to\mathrm{B}_{\mathrm{G}}\) the universal principal homogeneous \(\mathrm{G}\)-space. Let \(\mathrm{X}\) be a \(\mathrm{G}\)-space; \(\mathrm{X}\times\mathrm{U}_{\mathrm{G}}\) is a principal homogeneous \(\mathrm{G}\)-space over \(\mathrm{X}\times\mathrm{U}_{\mathrm{G}}/\mathrm{G}\), so that for every equivariant sheaf \(\mathrm{F}\) on \(\mathrm{X}\), \(\operatorname{pr}_1^*\mathrm{F}\) is the inverse image of a sheaf \(\mathrm{F}^{\mathrm{G}}\) on \(\mathrm{X}\times\mathrm{U}_{\mathrm{G}}/\mathrm{G}\). Since \(\mathrm{U}_{\mathrm{G}}\) is contractible and by (6.1.2.2), we have
\[
\mathrm{H}^{\bullet}(\mathrm{X},\mathrm{G};\mathrm{F})
\xrightarrow{\sim}
\mathrm{H}^{\bullet}(\mathrm{X}\times\mathrm{U}_{\mathrm{G}},\mathrm{G};
\operatorname{pr}_1^*\mathrm{F})
\xleftarrow{\sim}
\mathrm{H}^{\bullet}(\mathrm{X}\times\mathrm{U}_{\mathrm{G}}/\mathrm{G},\mathrm{F}^{\mathrm{G}}).
\tag*{\textbf{(6.1.4.1)}}
\]

In particular, for \(\mathrm{X}=\mathrm{P}^{t}\)
\[
\mathrm{H}^{\bullet}(\mathrm{B}_{\bullet\mathrm{G}})
=\mathrm{H}^{\bullet}(\mathrm{B}_{\mathrm{G}}).
\tag*{\textbf{(6.1.4.2)}}
\]

\Page{18}
One verifies that the isomorphism (6.1.4.2) is the special case of (6.1.3.1) for \(\mathrm{S}=\mathrm{B}_{\mathrm{G}}\), \(\mathrm{P}=\mathrm{U}_{\mathrm{G}}\).

\textbf{(6.1.5)} The spectral sequence (5.2.3.2) for \(\mathrm{B}_{\bullet\mathrm{G}}\)
\[
\mathrm{E}_{1}^{pq}=\mathrm{H}^{q}(\mathrm{G}^{\Delta_p}/\mathrm{G})
\Rightarrow
\mathrm{H}^{p+q}(\mathrm{B}_{\bullet\mathrm{G}})
=\mathrm{H}^{p+q}(\mathrm{B}_{\mathrm{G}})
\]
is essentially the \textit{Eilenberg-Moore spectral sequence}. Let us briefly recall how it can be used to relate the rational cohomology of \(\mathrm{G}\) and of \(\mathrm{B}_{\mathrm{G}}\), when \(\mathrm{G}\) is connected.

\textit{a)} The algebra \(\mathrm{H}^{\bullet}(\mathrm{G},\mathbf{Q})\) is a connected graded Hopf algebra of finite dimension over \(\mathbf{Q}\). If \(\mathrm{P}^{\bullet}(\mathrm{G})\) is the graded module of its primitive elements, we therefore have
\[
\mathrm{H}^{\bullet}(\mathrm{G},\mathbf{Q})
=\mathop{\textstyle\bigwedge}\mathrm{P}^{\bullet}(\mathrm{G}),
\]
and the generators of \(\mathrm{P}^{\bullet}(\mathrm{G})\) have odd degrees.

\textit{b)} The simplicial algebra \((\mathrm{E}_{1}^{p\bullet})_{p\geq 0}\) is
\[
\mathrm{E}_{1}^{p\bullet}
=\mathop{\textstyle\bigwedge}
(\mathrm{P}^{\bullet}(\mathrm{G})^{\Delta_p}/\mathrm{P}^{\bullet}(\mathrm{G}));
\]
it is the exterior algebra of the suspension of the constant cosimplicial module \(\mathrm{P}^{\bullet}(\mathrm{G})\); hence (Quillen [12]) \(\mathrm{E}_{2}^{p\bullet}=\operatorname{Sym}^{p}(\mathrm{P}^{\bullet}(\mathrm{G}))\). The terms \(\mathrm{E}_{2}^{pq}\) are nonzero only when \(p+q\) is even; thus \(\mathrm{E}_{2}^{pq}=\mathrm{E}_{\infty}^{pq}\) and, for a suitable filtration, we have canonically
\[
\operatorname{Gr}\mathrm{H}^{\bullet}(\mathrm{B}_{\mathrm{G}},\mathbf{Q})
\simeq\operatorname{Sym}^{\bullet}(\mathrm{P}^{\bullet}(\mathrm{G})[-1])
\]
and noncanonically
\[
\mathrm{H}^{\bullet}(\mathrm{B}_{\mathrm{G}},\mathbf{Q})
\simeq\operatorname{Sym}^{\bullet}(\mathrm{P}^{\bullet}(\mathrm{G})[-1]).
\]

\textbf{(6.1.6)} Take for \(\mathrm{G}\) a (complex) linear algebraic group.

\textit{a)} If \(\mathrm{T}\) is a maximal torus of \(\mathrm{G}\), with Weyl group \(\mathrm{W}\), we have
\[
\mathrm{H}^{\bullet}(\mathrm{B}_{\mathrm{G}},\mathbf{Q})
\xrightarrow{\sim}
\mathrm{H}^{\bullet}(\mathrm{B}_{\mathrm{T}},\mathbf{Q})^{\mathrm{W}}.
\tag*{\textbf{(6.1.6.1)}}
\]

\textit{b)} For a torus \(\mathrm{T}\) with character group \(\mathrm{X}(\mathrm{T})\), we have
\[
\mathrm{H}^{\bullet}(\mathrm{T},\mathbf{Z})
\simeq\overset{\bullet}{\mathop{\textstyle\bigwedge}}\mathrm{X}(\mathrm{T})
\tag*{\textbf{(6.1.6.2)}}
\]
(isomorphism of graded Hopf algebras).

We shall use \textit{a)} only in the following weaker form :

\(a')\) The map \(\mathrm{H}^{\bullet}(\mathrm{B}_{\mathrm{G}},\mathbf{Q})\to\mathrm{H}^{\bullet}(\mathrm{B}_{\mathrm{T}},\mathbf{Q})\) is injective \textit{(splitting principle)}.

For completeness, let us recall a proof of \(a')\).

If \(\mathrm{B}\) is a Borel subgroup of \(\mathrm{G}\), the bundle \(\mathrm{U}_{\mathrm{G}}/\mathrm{B}\) over \(\mathrm{B}_{\mathrm{G}}\) is a bundle of flag varieties. By ([3] \(2.1+2.6.3\)) as explained more clearly in P. A. Griffiths, Periods III, \textit{Publ. Math. I.H.E.S.}, \textbf{38}, prop. (3.1), or by [1], the Leray spectral sequence of \(\mathrm{U}_{\mathrm{G}}/\mathrm{B}\to\mathrm{B}_{\mathrm{G}}\) degenerates in rational cohomology. We therefore have \(\mathrm{H}^{\bullet}(\mathrm{B}_{\mathrm{G}},\mathbf{Q})\hookrightarrow\mathrm{H}^{\bullet}(\mathrm{U}_{\mathrm{G}}/\mathrm{B},\mathbf{Q})\). We conclude by noting that \(\mathrm{U}_{\mathrm{G}}/\mathrm{B}\sim\mathrm{B}_{\mathrm{B}}\sim\mathrm{B}_{\mathrm{T}}\).

\Page{19}
\section*{6.2. Construction of hypercoverings.}

\textbf{(6.2.1)} Let \(\mathrm{X}_{\bullet}\) be a simplicial set. Denote by \(\mathrm{D}(\Delta_n,\Delta_m)\) the set of surjective order-preserving maps from \(\Delta_n\) to \(\Delta_m\) (degeneracy operators), and set
\[
\mathrm{N}(\mathrm{X}_n)
=\mathrm{X}_n-\bigcup_{s\in\mathrm{D}(\Delta_n,\Delta_{n-1})}
s(\mathrm{X}_{n-1}).
\tag*{\textbf{(6.2.1.1)}}
\]

Recall that for every \(n\), the map
\[
\coprod s:
\coprod_{m\leq n}\coprod_{s\in\mathrm{D}(\Delta_n,\Delta_m)}
\mathrm{N}(\mathrm{X}_m)\to\mathrm{X}_n
\tag*{\textbf{(6.2.1.2)}}
\]
is bijective.

\textit{Definition} \textbf{(6.2.2)}. — \textit{A simplicial topological space is said to be \textup{\(s\)-split} if the maps \textup{(6.2.1.2)} are homeomorphisms.}

Let \(\mathrm{X}_{\bullet}\) be a \(k\)-truncated simplicial set. For \(n\leq k\), we again define \(\mathrm{N}(\mathrm{X}_n)\) by (6.2.1.1), and (6.2.1.2) is a bijection. A \(k\)-truncated simplicial topological space is said to be \textit{\(s\)-split} if (6.2.1.2) is a homeomorphism for \(n\leq k\).

\textbf{(6.2.3)} For \(\mathrm{X}\) an \((n+1)\)-truncated \(s\)-split simplicial topological space, augmented over \(\mathrm{S}\), let \(\alpha(\mathrm{X})\) be the triple consisting of \(\operatorname{sq}_n(\mathrm{X})\), \(\mathrm{N}\mathrm{X}_{n+1}\) and the obvious map from \(\mathrm{N}\mathrm{X}_{n+1}\) to \((\operatorname{cosq}\operatorname{sq}_n\mathrm{X})_{n+1}\).

This triple \((\mathrm{Y},\mathrm{N},\beta)\) satisfies

\((\ast)\) \(\mathrm{Y}\) is an \(n\)-truncated \(s\)-split simplicial topological space augmented over \(\mathrm{S}\), and \(\beta\) is a continuous map from \(\mathrm{N}\) to \((\operatorname{cosq}\mathrm{Y})_{n+1}\).

\textit{Proposition} \textbf{(6.2.4)} (\textit{S.G.A.}4, V \textit{bis} (5.1.3)). — \textit{Let \((\mathrm{Y},\mathrm{N},\beta)\) satisfy \((\ast)\) above.}

(i) \textit{Up to unique isomorphism, there exists one and only one topological space \(\mathrm{X}\), \((n+1)\)-truncated, \(s\)-split and augmented over \(\mathrm{S}\), with \(\alpha(\mathrm{X})\simeq(\mathrm{Y},\mathrm{N},\beta)\).}

(ii) \textit{Giving \(f:\mathrm{X}\to\mathrm{Z}\) is equivalent to giving :}

a) \textit{a morphism \(f':\mathrm{Y}\to\operatorname{sq}_n(\mathrm{Z})\),}

b) \textit{a morphism \(f'':\mathrm{N}\to\mathrm{Z}_{n+1}\) making the diagram commutative}
\[
\begin{tikzcd}[column sep=3em,row sep=3em]
\mathrm{N}\arrow[r,"\beta"]\arrow[d,"f''"']
  & (\operatorname{cosq}\mathrm{Y})_{n+1}\arrow[d,"f'"']\\
\mathrm{Z}_{n+1}\arrow[r]
  & (\operatorname{cosq}\operatorname{sq}_n\mathrm{Z})_{n+1}
\end{tikzcd}
\]

This proposition (6.2.4) also applies to simplicial objects in categories \(\mathcal{C}\) other than that of topological spaces. It suffices that \(\mathcal{C}\) satisfy
\Page{20}

\textbf{(6.2.4.1)} Finite inverse limits exist. Finite sums exist, are disjoint and universal.

\textbf{(6.2.5)} This proposition makes it possible to construct proper hypercoverings of \(\mathrm{S}\) by induction as follows.

0) We take \(f_0:\mathrm{X}_0\to\mathrm{S}\), proper and surjective. \(\{\mathrm{X}_0\}\) is a \(0\)-truncated proper hypercovering of \(\mathrm{S}\) (5.3.8), and is \(s\)-split.

1) We take \(f_1:\mathrm{N}_1\to\operatorname{cosq}(\{\mathrm{X}_0\})_1\), i.e. \(f_1:\mathrm{N}_1\to\mathrm{X}_0\times_{\mathrm{S}}\mathrm{X}_0\). Applying (6.2.4), we associate to \(f_1\) the \(1\)-truncated augmented \(s\)-split simplicial topological space
\[
\begin{tikzcd}[column sep=2em]
{}_{1}\mathrm{X}_{\bullet}:\ \mathrm{N}_1\amalg\mathrm{X}_0
  \arrow[r,shift left=1.1ex]\arrow[r,shift right=1.1ex]
  & \mathrm{X}_0\arrow[l]\arrow[r] & \mathrm{S}.
\end{tikzcd}
\]
We assume \(f_1\) is chosen so that
\[
f'_1:\mathrm{N}_1\amalg\mathrm{X}_0
\to\mathrm{X}_0\times_{\mathrm{S}}\mathrm{X}_0
\]
is proper and surjective (for example \(f_1\) proper and surjective). Then, \({}_{1}\mathrm{X}_{\bullet}\) is a \(1\)-truncated \(s\)-split proper hypercovering of \(\mathrm{S}\).

\noindent\(\ldotp\;\ldotp\;\ldotp\;\ldotp\;\ldotp\;\ldotp\;\ldotp\;\ldotp\;\ldotp\;\ldotp\;\ldotp\;\ldotp\;\ldotp\)

\(k+1)\) Suppose that a \(k\)-truncated \(s\)-split proper hypercovering \({}_{k}\mathrm{X}_{\bullet}\to\mathrm{S}\) has already been constructed. We take \(f_{k+1}:\mathrm{N}_{k+1}\to(\operatorname{cosq}({}_{k}\mathrm{X}_{\bullet}))_{k+1}\), and, applying (6.2.4), we associate to \(f_{k+1}\) a \(k+1\)-truncated augmented \(s\)-split semi-simplicial topological space \({}_{k+1}\mathrm{X}_{\bullet}\). We assume that \(f_{k+1}\) is such that
\[
f'_{k+1}:{}_{k+1}\mathrm{X}_{k+1}
\to\operatorname{cosq}({}_{k}\mathrm{X}_{\bullet})_{k+1}
\]
is proper and surjective (for example \(f_{k+1}\) proper and surjective). Then, \({}_{k+1}\mathrm{X}_{\bullet}\) is a \(k+1\)-truncated \(s\)-split proper hypercovering of \(\mathrm{S}\).

\noindent\(\ldotp\;\ldotp\;\ldotp\;\ldotp\;\ldotp\;\ldotp\;\ldotp\;\ldotp\;\ldotp\;\ldotp\;\ldotp\;\ldotp\;\ldotp\)

\(\infty)\) The \({}_{k}\mathrm{X}_{\bullet}\) so constructed are the successive skeleta of an \(s\)-split proper hypercovering of \(\mathrm{S}\).

\textbf{(6.2.6)} A simplicial scheme \(\mathrm{X}_{\bullet}\) over \(\mathbf{C}\) will be said to be \textit{smooth} if the \(\mathrm{X}_n\) are smooth; it will be said to be \textit{proper} if the \(\mathrm{X}_n\) are compact. A \textit{divisor with normal crossings} \(\mathrm{D}_{\bullet}\) on \(\mathrm{X}_{\bullet}\), assumed smooth, is a family \(\mathrm{D}_n\subset\mathrm{X}_n\) of divisors with normal crossings (3.1.2) such that the \(\mathrm{U}_n=\mathrm{X}_n-\mathrm{D}_n\) form a simplicial subscheme \(\mathrm{U}_{\bullet}\) of \(\mathrm{X}_{\bullet}\). This definition is justified by the following lemma.

\textit{Lemma} \textbf{(6.2.7)}. — \textit{If \(\mathrm{D}_{\bullet}\) is a divisor with normal crossings on \(\mathrm{X}_{\bullet}\), then the logarithmic De Rham complexes \((\Omega_{\mathrm{X}_n}^{\bullet}(\log\mathrm{D}_n))_{n\geq 0}\), equipped with the weight filtration \textup{(3.1.5)}, form a filtered complex on \(\mathrm{X}_{\bullet}\).}

This follows from (3.1.3 (ii)). The complex \((\Omega_{\mathrm{X}_n}^{\bullet}(\log\mathrm{D}_n))_{n\geq 0}\) will be denoted by \(\Omega_{\mathrm{X}_{\bullet}}^{\bullet}(\log\mathrm{D}_{\bullet})\).

\textbf{(6.2.8)} Using (6.2.5), one shows that for every separated scheme \(\mathrm{S}\) over \(\mathbf{C}\), there exist :

% ---- 021-028 ----
\Page{21}
\noindent
\(\alpha)\) a simplicial scheme \(\mathrm{X}_{\bullet}\) over \(\mathbf{C}\), proper and smooth, which may be taken to be \(s\)-split;

\noindent
\(\beta)\) a divisor with normal crossings \(\mathrm{D}_{\bullet}\) in \(\mathrm{X}_{\bullet}\); put \(\mathrm{U}_{\bullet}=\mathrm{X}_{\bullet}-\mathrm{D}_{\bullet}\);

\noindent
\(\gamma)\) an augmentation \(a:\mathrm{U}_{\bullet}\to\mathrm{S}\) making \(\mathrm{U}_{\bullet}^{\mathrm{an}}\) a proper hypercovering of \(\mathrm{S}^{\mathrm{an}}\).

Moreover, two such systems are dominated by a common third, and a morphism \(u:\mathrm{S}\to\mathrm{T}\) can be dominated by a morphism
\[
\begin{tikzcd}[column sep=1.3em,row sep=1.8em]
& \mathrm{X}_{\bullet} \arrow[rr] & & \mathrm{Y}_{\bullet} \\
\mathrm{U}_{\bullet} \arrow[ur,hook] \arrow[rr] \arrow[d,"a"']
& & \mathrm{V}_{\bullet} \arrow[ur,hook] \arrow[d,"b"'] & \\
\mathrm{S} \arrow[rr,"u"] & & \mathrm{T} &
\end{tikzcd}
\]
of such systems... (see \textit{S.G.A.}4, V \textit{bis}).

\section*{6.3. Relative cohomology.}

\textbf{(6.3.1)} The \textit{mapping cone} construction for morphisms of simplicial sets carries over unchanged to simplicial objects in any category \(\mathscr{C}\) having a final object \(e\) and finite sums. For \(u:\mathrm{Y}_{\bullet}\to\mathrm{X}_{\bullet}\), the cone \(\mathrm{C}(u)\) satisfies
\[
\mathrm{C}(u)_n
=\mathrm{X}_n\amalg\coprod_{i<n}\mathrm{Y}_i\amalg e.
\]

We shall take for \(\mathscr{C}\):

\noindent
\textit{a)} the category of topological spaces and continuous maps (final object: \(\mathrm{P}^{t}\));

\noindent
\textit{b)} the category of pairs \((\mathrm{X},\mathrm{F})\) consisting of a topological space \(\mathrm{X}\) and an abelian sheaf \(\mathrm{F}\) on \(\mathrm{X}\), an arrow \((u,f):(\mathrm{Y},\mathrm{F})\to(\mathrm{X},\mathrm{G})\) being a continuous map \(u:\mathrm{Y}\to\mathrm{X}\), together with a \(u\)-morphism (5.1.5) \(f:\mathrm{G}\to\mathrm{F}\) (final object: \((\mathrm{P}^{t},0)\)).

\textbf{(6.3.2)} Let \(u:\mathrm{Y}_{\bullet}\to\mathrm{X}_{\bullet}\) be a morphism of simplicial topological spaces, with cone \(\mathrm{C}(u)\), \(\mathrm{F}\) an abelian sheaf on \(\mathrm{X}_{\bullet}\), \(\mathrm{G}\) an abelian sheaf on \(\mathrm{Y}_{\bullet}\), and \(f:\mathrm{G}\to\mathrm{F}\) a \(u\)-morphism. The cone \(\mathrm{C}(f)\) of \(f\) is an abelian sheaf on \(\mathrm{C}(u)\), and we set
\[
\mathrm{H}^{n}(\mathrm{C}(u),\mathrm{C}(f))
=\mathrm{H}^{n}(\mathrm{X}_{\bullet}\operatorname{mod}\mathrm{Y}_{\bullet},
\mathrm{F}\operatorname{mod}\mathrm{G}).
\tag*{\textbf{(6.3.2.1)}}
\]
These are the \textit{relative cohomology groups}. One easily checks that they fit into a long exact sequence
\[
\ldots\to
\mathrm{H}^{i}(\mathrm{X}_{\bullet}\operatorname{mod}\mathrm{Y}_{\bullet},
\mathrm{F}\operatorname{mod}\mathrm{G})
\to\mathrm{H}^{i}(\mathrm{X}_{\bullet},\mathrm{F})
\to\mathrm{H}^{i}(\mathrm{Y}_{\bullet},\mathrm{G})
\to\ldots
\tag*{\textbf{(6.3.2.2)}}
\]

\textbf{(6.3.3)} More generally, let \(\mathrm{L}\) and \(\mathrm{K}\) be bounded below complexes of abelian sheaves on \(\mathrm{Y}_{\bullet}\) and \(\mathrm{X}_{\bullet}\), and let \(f:\mathrm{K}\to\mathrm{L}\) be a \(u\)-morphism. This gives a complex \(\mathrm{C}(f)\) on \(\mathrm{C}(u)\). In hypercohomology we again set
\[
\mathbf{H}^{n}(\mathrm{C}(u),\mathrm{C}(f))
=\mathbf{H}^{n}(\mathrm{X}_{\bullet}\operatorname{mod}\mathrm{Y}_{\bullet},
\mathrm{K}\operatorname{mod}\mathrm{L}).
\]
\Page{22}
These groups occur in an exact sequence analogous to (6.3.2.2), arising, in the appropriate derived category, from a distinguished triangle
\[
\begin{tikzcd}[column sep=0.5em,row sep=2em]
& \mathrm{R}\Gamma(\mathrm{Y}_{\bullet},\mathrm{K})
  \arrow[dl,"\partial"'] & \\
\mathrm{R}\Gamma(\mathrm{C}(u),\mathrm{C}(f)) \arrow[rr]
& & \mathrm{R}\Gamma(\mathrm{X}_{\bullet},\mathrm{L}) \arrow[ul]
\end{tikzcd}
\tag*{\textbf{(6.3.3.1)}}
\]

\textbf{(6.3.4)} The construction presented above is not the only possible one. It has the disadvantage that, even if one starts with genuine topological spaces \(\mathrm{X}\) and \(\mathrm{Y}\) (i.e. with constant simplicial spaces), one is led to consider nonconstant simplicial spaces.

\section*{6.4. Multisimplicial spaces.}

\textbf{(6.4.1)} Let \(r\) be an integer \(\geq 0\). An \textit{\(r\)-simplicial object} \(\mathrm{Z}_{\bullet}\) of a category \(\mathscr{C}\) is a contravariant functor from \((\Delta)^r\) to \(\mathscr{C}\). The \textit{diagonal} simplicial object \(\delta\mathrm{Z}_{\bullet}\) is the composite functor \((\Delta)\to(\Delta)^r\to\mathscr{C}\).

\textbf{(6.4.2)} The topos of \textit{sheaves} on an \(r\)-simplicial topological space is defined as in (5.1.6). Let \(\Gamma^{\bullet}\) be the functor
\[
\mathrm{F}\mapsto
(\Gamma(\mathrm{X}_{n_1\ldots n_r},\mathrm{F}^{n_1,\ldots,n_r}))
\]
from sheaves on \(\mathrm{X}_{\bullet}\) to \(r\)-cosimplicial sets. For small \(r\), it will often be denoted instead by \(\Gamma^{\bullet\cdots\bullet}\) (\(r\) dots). There is a coaugmentation \(\Gamma(\mathrm{X}_{\bullet},\mathrm{F}^{\bullet})\to\Gamma^{\bullet}(\mathrm{X}_{\bullet},\mathrm{F}^{\bullet})\). The functors of \textit{cohomology with values in an abelian sheaf} \(\mathrm{F}\) are the derived functors of the “global sections” functor \(\Gamma\). They can be computed by a procedure parallel to (5.2.3):
\[
\mathrm{R}\Gamma=\boldsymbol{s}\mathrm{R}\Gamma^{\bullet}:
\mathrm{D}^{+}(\mathrm{X}_{\bullet})\to\mathrm{D}^{+}((\mathrm{Ab})).
\tag*{\textbf{(6.4.2.1)}}
\]

A sheaf \(\mathrm{F}\) on an \(r\)-simplicial topological space \(\mathrm{X}_{\bullet}\) induces a sheaf \(\delta\mathrm{F}\) on the diagonal simplicial space \(\delta\mathrm{X}_{\bullet}\). It follows from the Cartier-Eilenberg-Zilber theorem that
\[
\mathrm{R}\Gamma(\mathrm{X}_{\bullet},\mathrm{K})
\xrightarrow{\sim}
\mathrm{R}\Gamma(\delta\mathrm{X}_{\bullet},\delta\mathrm{K}).
\tag*{\textbf{(6.4.2.2)}}
\]

\textbf{(6.4.3)} Let us restrict ourselves to the case \(r=2\). A bisimplicial object 1-augmented to a simplicial object \(\mathrm{S}_{\bullet}\) is a contravariant functor from \((\Delta^{+})\times(\Delta)\) to \(\mathscr{C}\), such that \(\mathrm{S}_{\bullet}\) is the composite functor
\[
(\Delta)\xrightarrow{(\Delta_{-1},\bullet)}
(\Delta^{+})\times(\Delta)\longrightarrow\mathscr{C}.
\]
To denote a bisimplicial object 1-augmented to \(\mathrm{S}_{\bullet}\), with underlying bisimplicial object \(\mathrm{X}_{\bullet\bullet}\), we shall use notation of the form
\[
a:\mathrm{X}_{\bullet\bullet}\to\mathrm{S}_{\bullet}.
\]
\Page{23}
For \(n\geq 0\), \(a:(\mathrm{X}_{\bullet n})\to\mathrm{S}_n\) is a simplicial object augmented to \(\mathrm{S}_n\). If \(\mathrm{F}\) is a sheaf on \(\mathrm{X}_{\bullet\bullet}\), the \((a_{\ast}(\mathrm{F}\mid\mathrm{X}_{\bullet n})\text{ on }\mathrm{S}_n)_{n\geq 0}\) form a sheaf on \(\mathrm{S}_{\bullet}\); this defines a morphism of topoi
\[
a:\mathrm{X}_{\bullet\bullet}^{\sim}\to\mathrm{S}_{\bullet}^{\sim}.
\]

It is explained in \textit{S.G.A.}4, V \textit{bis} that \(\mathrm{R}a_{\ast}\) is computed “component by component”:
\[
\mathrm{R}a_{\ast}\mathrm{K}\mid\mathrm{S}_n
=\mathrm{R}a_{\ast}(\mathrm{K}\mid\mathrm{X}_{\bullet n}).
\]
It follows that if, for each \(n\), \(a:(\mathrm{X}_{\bullet n})\to\mathrm{S}_n\) is of cohomological descent, then \(a:\mathrm{X}_{\bullet\bullet}\to\mathrm{S}_{\bullet}\) is of cohomological descent: for every complex of abelian sheaves \(\mathrm{K}\in\mathrm{D}^{+}(\mathrm{S}_{\bullet})\), we have
\[
\mathrm{K}\xrightarrow{\approx}\mathrm{R}a_{\ast}a^{\ast}\mathrm{K}.
\]

\textbf{(6.4.4)} In \textit{S.G.A.}4, V \textit{bis}, it is shown that for every separated simplicial scheme \(\mathrm{S}_{\bullet}\), there exist a bisimplicial scheme \(\mathrm{X}_{\bullet\bullet}\) 1-augmented to \(\mathrm{S}_{\bullet}\), and \(i:\mathrm{X}_{\bullet\bullet}\hookrightarrow\overline{\mathrm{X}}_{\bullet\bullet}\) such that:

\textit{a)} The \(\overline{\mathrm{X}}_{nm}\) are projective and smooth; \(\mathrm{X}_{nm}\) is the complement of a divisor with normal crossings \(\mathrm{D}_{nm}\) in \(\overline{\mathrm{X}}_{nm}\). We may assume that \(\mathrm{D}_{nm}\) is a union of smooth divisors.

\textit{b)} For \(n\geq 0\), \(\mathrm{X}_{\bullet n}\) is a proper hypercovering of \(\mathrm{S}_n\). We may assume it to be \(s\)-split.

The construction proceeds as in (6.2), but the induction is more complicated. The uniqueness assertions (6.2.8) remain valid, \textit{mutatis mutandis}.

\section*{7. SUPPLEMENTS TO § 1}

\section*{7.1. Filtered derived category.}

This subsection supplements § 1, n\textsuperscript{o} 4.

\textbf{(7.1.1)} Let \(\mathscr{A}\) be an abelian category. Set:

\noindent
\(\mathrm{F}\mathscr{A}\) (resp. \(\mathrm{F}_2\mathscr{A}\)) \(=\) the category of filtered (resp. bifiltered) objects with finite filtration(s) in \(\mathscr{A}\).

\noindent
\(\mathrm{K}^{+}\mathrm{F}\mathscr{A}\) (resp. \(\mathrm{K}^{+}\mathrm{F}_2\mathscr{A}\)) \(=\) the category of bounded below filtered (resp. bifiltered) complexes of objects of \(\mathscr{A}\), up to homotopy preserving the filtration (resp. the filtrations).

\noindent
\(\mathrm{D}^{+}\mathrm{F}\mathscr{A}\) (resp. \(\mathrm{D}^{+}\mathrm{F}_2\mathscr{A}\)) \(=\) the triangulated category obtained from \(\mathrm{K}^{+}\mathrm{F}\mathscr{A}\) (resp. \(\mathrm{K}^{+}\mathrm{F}_2\mathscr{A}\)) by inverting the filtered (resp. bifiltered) quasi-isomorphisms (1.3.6). This is the \textit{filtered derived category} (resp. \textit{bifiltered}).

\textbf{(7.1.2)} A filtered quasi-isomorphism \(u:(\mathrm{K},\mathrm{F})\to(\mathrm{K}',\mathrm{F}')\) induces an isomorphism of spectral sequences \(u:\mathrm{E}_{\bullet}^{\bullet\bullet}(\mathrm{K},\mathrm{F})\to\mathrm{E}_{\bullet}^{\bullet\bullet}(\mathrm{K}',\mathrm{F}')\). An object \(\mathrm{K}\) of \(\mathrm{D}^{+}\mathrm{F}\mathscr{A}\) therefore defines
\Page{24}
a spectral sequence \(\mathrm{E}_{\bullet}^{\bullet\bullet}(\mathrm{K})\). Similarly, an object \(\mathrm{L}\) of \(\mathrm{D}^{+}\mathrm{F}_2\mathscr{A}\) defines a collection of spectral sequences of the type considered in (1.4.9).

\textbf{(7.1.3)} Let \(\mathrm{T}\) be a left exact functor from \(\mathscr{A}\) to an abelian category \(\mathscr{B}\). We assume that every object of \(\mathscr{A}\) embeds into an injective object. The functor \(\mathrm{T}\) can then be “derived” to give functors
\[
\mathrm{RT}:\mathrm{D}^{+}(\mathscr{A})\to\mathrm{D}^{+}(\mathscr{B})
\qquad\text{(for reference)}
\tag*{\textbf{(7.1.3.1)}}
\]
\[
\mathrm{RT}:\mathrm{D}^{+}\mathrm{F}(\mathscr{A})
\to\mathrm{D}^{+}\mathrm{F}(\mathscr{B}),
\tag*{\textbf{(7.1.3.2)}}
\]
\[
\mathrm{RT}:\mathrm{D}^{+}\mathrm{F}_2(\mathscr{A})
\to\mathrm{D}^{+}\mathrm{F}_2(\mathscr{B}).
\tag*{\textbf{(7.1.3.3)}}
\]
They are computed as follows: if \(\mathrm{K}'\) is a \(\mathrm{T}\)-acyclic resolution (resp. a filtered resolution, resp. a bifiltered resolution) of \(\mathrm{K}\) ((1.4.5), (1.4.11)), then \(\mathrm{RT}(\mathrm{K})=\mathrm{T}(\mathrm{K}')\).

The hypercohomology spectral sequence (for \(\mathrm{T}\)) of \(\mathrm{K}\in\operatorname{Ob}\mathrm{D}^{+}\mathrm{F}(\mathscr{A})\) is the spectral sequence of \(\mathrm{RTK}\in\operatorname{Ob}\mathrm{D}^{+}\mathrm{F}(\mathscr{B})\) (cf. (1.4.5)).

\textbf{(7.1.4)} We shall need more precise results for the functors \(\mathrm{R}a_{\ast}\), where \(a:\mathrm{X}_{\bullet}\to\mathrm{S}\) is an augmentation of a simplicial topological space. The case \(\mathrm{S}=\mathrm{P}^{t}\), \(\mathrm{R}a_{\ast}=\mathrm{R}\Gamma\) would suffice for us.

Let us return to the notation of (5.1.12), and recall formula (5.2.6.1):
\[
\mathrm{R}a_{\ast}=\boldsymbol{s}\mathrm{R}a_{\bullet\ast}.
\]
For every complex \(\mathrm{K}\in\operatorname{Ob}\mathrm{C}^{+}(\mathrm{S}_{\bullet})\), the simple complex \(\boldsymbol{s}\mathrm{K}\) is equipped with a natural filtration \(\mathrm{L}\) (5.1.9.3). A quasi-isomorphism \(u:\mathrm{K}'\xrightarrow{\approx}\mathrm{K}''\) induces a filtered quasi-isomorphism \(u:(\boldsymbol{s}\mathrm{K}',\mathrm{L})\to(\boldsymbol{s}\mathrm{K}'',\mathrm{L})\). Thus, \(\boldsymbol{s}\) factors through
\[
\boldsymbol{s}:\mathrm{D}^{+}(\mathrm{S}_{\bullet})
\to\mathrm{D}^{+}\mathrm{F}(\mathrm{S})
\]
and \(\mathrm{R}a_{\ast}\) factors through
\[
\mathrm{R}a_{\ast}:\mathrm{D}^{+}(\mathrm{X}_{\bullet})
\to\mathrm{D}^{+}\mathrm{F}(\mathrm{S}).
\]
The spectral sequence of the filtered complex \((\mathrm{R}a_{\ast}\mathrm{K},\mathrm{L})\in\mathrm{D}^{+}\mathrm{F}(\mathrm{S})\) is precisely (5.2.7.1).

\textbf{(7.1.5)} If \(\mathrm{K}\) is filtered (resp. bifiltered), then \(\mathrm{R}a_{\bullet\ast}\mathrm{K}\) is filtered (bifiltered): we have
\[
\mathrm{R}a_{\bullet\ast}:
\mathrm{D}^{+}\mathrm{F}(\mathrm{X}_{\bullet})
\to\mathrm{D}^{+}\mathrm{F}(\mathrm{S}_{\bullet})
\tag*{\textbf{(7.1.5.1)}}
\]
\[
\mathrm{R}a_{\bullet\ast}:
\mathrm{D}^{+}\mathrm{F}_2(\mathrm{X}_{\bullet})
\to\mathrm{D}^{+}\mathrm{F}_2(\mathrm{S}_{\bullet}).
\tag*{\textbf{(7.1.5.2)}}
\]

\textbf{(7.1.6)} Let \(\mathrm{K}\) be a complex of cosimplicial sheaves on \(\mathrm{S}\), equipped with an increasing filtration \(\mathrm{W}\). The \textit{diagonal filtration} \(\delta(\mathrm{W},\mathrm{L})\) of \(\mathrm{W}\) and \(\mathrm{L}\) is the following increasing filtration of \(\boldsymbol{s}\mathrm{K}\)
\begin{align*}
\delta(\mathrm{W},\mathrm{L})_n(\boldsymbol{s}\mathrm{K})
&=\bigoplus_{p,q}\mathrm{W}_{n+p}(\mathrm{K}^{p,q})
\tag*{\textbf{(7.1.6.1)}}\\
&=\sum_p\boldsymbol{s}(\mathrm{W}_{n+p}(\mathrm{K}))
\cap\mathrm{L}^{p}(\boldsymbol{s}\mathrm{K}).
\end{align*}

\Page{25}
We have
\[
\operatorname{Gr}_{n}^{\delta(\mathrm{W},\mathrm{L})}
(\boldsymbol{s}\mathrm{K})
\simeq
\bigoplus_p\operatorname{Gr}_{n+p}^{\mathrm{W}}
(\mathrm{K}^{\bullet p})[-p].
\tag*{\textbf{(7.1.6.2)}}
\]

The functor \((\mathrm{K},\mathrm{W})\mapsto
(\boldsymbol{s}\mathrm{K},\delta(\mathrm{W},\mathrm{L}))\) carries filtered quasi-isomorphisms to filtered quasi-isomorphisms and defines
\[
(\boldsymbol{s},\delta):
(\mathrm{K},\mathrm{W})\mapsto
(\boldsymbol{s}\mathrm{K},\delta(\mathrm{W},\mathrm{L})):
\mathrm{D}^{+}\mathrm{F}(\mathrm{S}_{\bullet})
\to\mathrm{D}^{+}\mathrm{F}(\mathrm{S})
\tag*{\textbf{(7.1.6.3)}}
\]
whence, by composition with (7.1.5.1),
\[
(\mathrm{R}\Gamma,\delta(\, ,\mathrm{L})):
(\mathrm{K},\mathrm{W})\mapsto
(\mathrm{R}\Gamma\mathrm{K},\delta(\mathrm{W},\mathrm{L})):
\mathrm{D}^{+}\mathrm{F}(\mathrm{X}_{\bullet})
\to\mathrm{D}^{+}\mathrm{F}(\mathrm{S}).
\tag*{\textbf{(7.1.6.4)}}
\]
From (7.1.6.2) we deduce that
\[
\operatorname{Gr}_{n}^{\delta(\mathrm{W},\mathrm{L})}
(\mathrm{R}a_{\ast}\mathrm{K})
=\bigoplus_p\mathrm{R}a_{p\ast}
(\operatorname{Gr}_{n+p}^{\mathrm{W}}\mathrm{K})[-p].
\tag*{\textbf{(7.1.6.5)}}
\]

\textbf{(7.1.7)} If \((\mathrm{K},\mathrm{W},\mathrm{F})\) is a bifiltered complex of cosimplicial sheaves, then \(\boldsymbol{s}\mathrm{K}\) is equipped with the three filtrations \(\mathrm{W}\), \(\mathrm{F}\), and \(\mathrm{L}\) and defines various bifiltered complexes.

For example, when \(\mathrm{W}\) is increasing, the functor \((\mathrm{K},\mathrm{W},\mathrm{F})\mapsto
(\mathrm{K},\delta(\mathrm{W},\mathrm{L}),\mathrm{F})\) carries a bifiltered quasi-isomorphism to a bifiltered quasi-isomorphism, hence defines
\[
(\mathrm{K},\mathrm{W},\mathrm{F})\mapsto
(\mathrm{K},\delta(\mathrm{W},\mathrm{L}),\mathrm{F}):
\mathrm{D}^{+}\mathrm{F}_2(\mathrm{S}_{\bullet})
\to\mathrm{D}^{+}\mathrm{F}_2(\mathrm{S}).
\tag*{\textbf{(7.1.7.1)}}
\]
By composition with (7.1.5.2), we deduce
\[
(\mathrm{K},\mathrm{W},\mathrm{F})\mapsto
(\mathrm{R}\Gamma\mathrm{K},\delta(\mathrm{W},\mathrm{L}),\mathrm{F}):
\mathrm{D}^{+}\mathrm{F}_2(\mathrm{X}_{\bullet})
\to\mathrm{D}^{+}\mathrm{F}_2(\mathrm{S}),
\tag*{\textbf{(7.1.7.2)}}
\]
and we have
\[
\operatorname{Gr}_{n}^{\delta(\mathrm{W},\mathrm{L})}
(\mathrm{R}\Gamma\mathrm{K},\mathrm{F})
=\bigoplus_p\mathrm{R}a_{p\ast}
(\operatorname{Gr}_{n+p}^{\mathrm{W}}\mathrm{K},\mathrm{F})[-p]
\tag*{\textbf{(7.1.7.3)}}
\]
in \(\mathrm{D}^{+}\mathrm{F}(\mathrm{S})\).

\section*{7.2. Supplements to the lemma on two filtrations.}

In this subsection, we give a new proof of the lemma on two filtrations (1.3.16) and some supplements.

\textbf{(7.2.1)} Let \((\mathrm{K},\mathrm{F},\mathrm{W})\) be a bounded below bifiltered complex of objects of an abelian category \(\mathscr{A}\). The filtration \(\mathrm{F}\) is assumed to be biregular.

We say that \((\mathrm{K},\mathrm{F},\mathrm{W})\) is \(\mathrm{F}\)-\textit{splittable} if the filtered complex \((\mathrm{K},\mathrm{W})\) can be represented as a sum of filtered complexes \((\mathrm{K}_n,\mathrm{W}_n)_{n\in\mathbf{Z}}\)
\[
(\mathrm{K},\mathrm{W})=\bigoplus_n(\mathrm{K}_n,\mathrm{W}_n)
\]
with
\[
\mathrm{F}^{n}\mathrm{K}=\bigoplus_{n'\geq n}\mathrm{K}_{n'}.
\]

\Page{26}
Let \(r_0\) be an integer \(\geq 0\), or \(+\infty\). The following condition was considered in (1.3.16), (1.3.17):

\textbf{(7.2.2; \(r_0\))} For every nonnegative integer \(r<r_0\), the differentials \(d_r\) of the graded complex \(\mathrm{E}_r(\mathrm{K},\mathrm{W})\) are strictly compatible with the recurrent filtration defined by \(\mathrm{F}\).

\textbf{(7.2.3)} It is clear that if \((\mathrm{K},\mathrm{F},\mathrm{W})\) is \(\mathrm{F}\)-splittable, condition (7.2.2; \(\infty\)) is satisfied. Conversely, it seems that when condition (7.2.2; \(\infty\)) is satisfied, everything behaves as if the functor \(\operatorname{Gr}_{\mathrm{F}}\) were exact. For example, we shall show that the \(\operatorname{Gr}_{\mathrm{F}}\) of the spectral sequence \(\mathrm{E}(\mathrm{K},\mathrm{W})\) can then be identified with the spectral sequence \(\mathrm{E}(\operatorname{Gr}_{\mathrm{F}}\mathrm{K},\mathrm{W})\), and that the spectral sequence \(\mathrm{E}(\mathrm{K},\mathrm{F})\) degenerates \((\mathrm{E}_1=\mathrm{E}_{\infty})\).

\textbf{(7.2.4)} It follows immediately from definition (1.3.8) that the first direct filtration of \(\mathrm{E}_r(\mathrm{K},\mathrm{W})\) is the filtration \(\mathrm{F}_d\) of \(\mathrm{E}_r(\mathrm{K},\mathrm{W})\) by the images
\[
\mathrm{F}_d^p(\mathrm{E}_r(\mathrm{K},\mathrm{W}))
=\operatorname{Im}(\mathrm{E}_r(\mathrm{F}^p\mathrm{K},\mathrm{W})
\to\mathrm{E}_r(\mathrm{K},\mathrm{W})).
\]
Dually, the second direct filtration (1.3.9) of \(\mathrm{E}_r(\mathrm{K},\mathrm{W})\) is the filtration \(\mathrm{F}_{d^{*}}\) of \(\mathrm{E}_r(\mathrm{K},\mathrm{W})\) by the kernels
\[
\mathrm{F}_{d^{*}}^p(\mathrm{E}_r(\mathrm{K},\mathrm{W}))
=\operatorname{Ker}(\mathrm{E}_r(\mathrm{K},\mathrm{W})
\to\mathrm{E}_r(\mathrm{K}/\mathrm{F}^p\mathrm{K},\mathrm{W})).
\]
The recurrent filtration \(\mathrm{F}_{\mathrm{rec}}\) of \(\mathrm{E}_r(\mathrm{K},\mathrm{W})\) lies between these two filtrations (1.3.13 (iii)).

\textit{Proposition} \textbf{(7.2.5)}. — \textit{Suppose that \((\mathrm{K},\mathrm{F},\mathrm{W})\) satisfies \((7.2.2;\,r_0)\). Then}

\noindent
(i) \textit{\((\mathrm{F}^a\mathrm{K}/\mathrm{F}^b\mathrm{K},\mathrm{F},\mathrm{W})\) also satisfies \((7.2.2;\,r_0)\).}

\noindent
(ii) \textit{For \(r\leq r_0\), the sequence}
\[
0\to\mathrm{E}_r(\mathrm{F}^p\mathrm{K},\mathrm{W})
\to\mathrm{E}_r(\mathrm{K},\mathrm{W})
\to\mathrm{E}_r(\mathrm{K}/\mathrm{F}^p\mathrm{K},\mathrm{W})\to 0
\]
\textit{is exact; for \(r=r_0+1\), the sequence}
\[
\mathrm{E}_r(\mathrm{F}^p\mathrm{K},\mathrm{W})
\to\mathrm{E}_r(\mathrm{K},\mathrm{W})
\to\mathrm{E}_r(\mathrm{K}/\mathrm{F}^p\mathrm{K},\mathrm{W})
\]
\textit{is exact. In particular, for \(r\leq r_0+1\), the two direct filtrations and the recurrent filtration of \(\mathrm{E}_r(\mathrm{K},\mathrm{W})\) coincide.}

Let us fix an integer \(p\) and prove by induction on \(r\) the assertion

\((\ast)_r\) If \(r\leq r_0\), \(\mathrm{E}_r(\mathrm{F}^p\mathrm{K},\mathrm{W})\) embeds into \(\mathrm{E}_r(\mathrm{K},\mathrm{W})\); if \(r\leq r_0+1\), its image is \(\mathrm{F}_{\mathrm{rec}}^p(\mathrm{E}_r(\mathrm{K},\mathrm{W}))\).

The assertion \((\ast)_0\) is always true. Assume \((\ast)_r\) and prove \((\ast)_{r+1}\). We may assume that \(0\leq r\leq r_0\). We have a diagram
\[
\begin{tikzcd}[column sep=2em,row sep=3em]
\mathrm{E}_r(\mathrm{K},\mathrm{W}) \arrow[r,"d_r"]
& \mathrm{E}_r(\mathrm{K},\mathrm{W}) \arrow[r,"d_r"]
& \mathrm{E}_r(\mathrm{K},\mathrm{W}) \\
\mathrm{E}_r(\mathrm{F}^p\mathrm{K},\mathrm{W}) \arrow[u,hook] \arrow[r]
& \mathrm{E}_r(\mathrm{F}^p\mathrm{K},\mathrm{W}) \arrow[u,hook] \arrow[r]
& \mathrm{E}_r(\mathrm{F}^p\mathrm{K},\mathrm{W}), \arrow[u,hook,"i"']
\end{tikzcd}
\]
\Page{27}
and the image of the vertical inclusions is \(\mathrm{F}_{\mathrm{rec}}^p(\mathrm{E}_r(\mathrm{K},\mathrm{W}))\). Since \(i\) is injective
\begin{align*}
\mathrm{F}_{\mathrm{rec}}^p(\mathrm{E}_{r+1}(\mathrm{K},\mathrm{W}))
&=\operatorname{Im}(\operatorname{Ker}(
\mathrm{F}_{\mathrm{rec}}^p\mathrm{E}_r(\mathrm{K},\mathrm{W})
\xrightarrow{d_r}\mathrm{E}_r(\mathrm{K},\mathrm{W}))
\to\mathrm{E}_{r+1}(\mathrm{K},\mathrm{W}))\\
&=\operatorname{Im}(\operatorname{Ker}(
\mathrm{E}_r(\mathrm{F}^p\mathrm{K},\mathrm{W})
\xrightarrow{d_r}\mathrm{E}_r(\mathrm{F}^p\mathrm{K},\mathrm{W}))
\to\mathrm{E}_{r+1}(\mathrm{K},\mathrm{W}))\\
&=\operatorname{Im}(
\mathrm{E}_{r+1}(\mathrm{F}^p\mathrm{K},\mathrm{W})
\to\mathrm{E}_{r+1}(\mathrm{K},\mathrm{W})).
\end{align*}
If \(r<r_0\), \(d_r\) is strictly compatible with \(\mathrm{F}_{\mathrm{rec}}\), hence
\[
d_r\mathrm{E}_r(\mathrm{K},\mathrm{W})
\cap\mathrm{E}_r(\mathrm{F}^p\mathrm{K},\mathrm{W})
=d_r\mathrm{E}_r(\mathrm{F}^p\mathrm{K},\mathrm{W}),
\]
and \(\mathrm{E}_{r+1}(\mathrm{F}^p\mathrm{K},\mathrm{W})\) therefore embeds into \(\mathrm{E}_{r+1}(\mathrm{K},\mathrm{W})\). This proves \((\ast)_{r+1}\).

The assertions \((\ast)_r\), together with the dual assertions, imply (ii).

It follows from (ii) that for \(r<r_0\) and \(a\leq b\), \(\mathrm{E}_r(\mathrm{F}^b\mathrm{K},\mathrm{W})\) embeds into \(\mathrm{E}_r(\mathrm{F}^a\mathrm{K},\mathrm{W})\) and that the differential \(d_r\) of \(\mathrm{E}_r(\mathrm{F}^a\mathrm{K},\mathrm{W})\) is strictly compatible with the first direct filtration of \(\mathrm{E}_r(\mathrm{F}^a\mathrm{K},\mathrm{W})\). We deduce by induction on \(r<r_0\) that the first direct filtration of \(\mathrm{E}_r(\mathrm{F}^a\mathrm{K},\mathrm{W})\) coincides with the recurrent filtration, and (i) follows.

\textbf{(7.2.6)} If condition (7.2.2; \(r_0\)) is satisfied, and \(r\leq r_0+1\), we shall denote by \(\mathrm{F}\) the filtration \(\mathrm{F}_d=\mathrm{F}_{d^{*}}=\mathrm{F}_{\mathrm{rec}}\) of \(\mathrm{E}_r(\mathrm{K},\mathrm{W})\). The exact sequences (7.2.5) (ii)
\[
\begin{tikzcd}[column sep=0.6em,row sep=1.8em]
0
& \mathrm{E}_r(\mathrm{K}/\mathrm{F}^p\mathrm{K},\mathrm{W}) \arrow[l]
& & \mathrm{E}_r(\mathrm{K}/\mathrm{F}^{p+1}\mathrm{K},\mathrm{W}) \arrow[ll]
& & 0 \\
& & \mathrm{E}_r(\mathrm{K},\mathrm{W}) \arrow[ul] \arrow[ur]
& & \mathrm{E}_r(\operatorname{Gr}_{\mathrm{F}}^p\mathrm{K},\mathrm{W})
    \arrow[ul] \arrow[ur]
& \\
0 \arrow[r]
& \mathrm{E}_r(\mathrm{F}^{p+1}\mathrm{K},\mathrm{W}) \arrow[ur] \arrow[rr]
& & \mathrm{E}_r(\mathrm{F}^p\mathrm{K},\mathrm{W}) \arrow[ul] \arrow[ur]
& & 0 \arrow[ul]
\end{tikzcd}
\]
define, for \(r\leq r_0\), a (self-dual) isomorphism compatible with the differentials \(d_r\)
\[
\operatorname{Gr}_{\mathrm{F}}^p(\mathrm{E}_r(\mathrm{K},\mathrm{W}))
\simeq\mathrm{E}_r(\operatorname{Gr}_{\mathrm{F}}^p\mathrm{K},\mathrm{W})
\qquad(r\leq r_0).
\tag*{\textbf{(7.2.6.1)}}
\]

\textbf{(7.2.7)} If condition (7.2.2; \(\infty\)) is satisfied, the sequences
\[
0\to\mathrm{E}_r(\mathrm{F}^p\mathrm{K},\mathrm{W})
\to\mathrm{E}_r(\mathrm{K},\mathrm{W})
\to\mathrm{E}_r(\mathrm{K}/\mathrm{F}^p\mathrm{K},\mathrm{W})\to 0
\]
are exact for every \(r\). If the filtration \(\mathrm{W}\) is biregular, the sequence
\[
0\to\mathrm{E}_{\infty}(\mathrm{F}^p\mathrm{K},\mathrm{W})
\to\mathrm{E}_{\infty}(\mathrm{K},\mathrm{W})
\to\mathrm{E}_{\infty}(\mathrm{K}/\mathrm{F}^p\mathrm{K},\mathrm{W})\to 0
\]
is therefore exact. We may rewrite this sequence as
\[
0\to\operatorname{Gr}_{\mathrm{W}}(\mathrm{H}(\mathrm{F}^p\mathrm{K}))
\to\operatorname{Gr}_{\mathrm{W}}(\mathrm{H}(\mathrm{K}))
\to\operatorname{Gr}_{\mathrm{W}}(\mathrm{H}(\mathrm{K}/\mathrm{F}^p\mathrm{K}))
\to 0.
\]

\textit{Lemma} \textbf{(7.2.7.1)}. — \textit{Let \(\mathrm{L}\) be a complex equipped with a biregular filtration \(\mathrm{G}\). For \(\operatorname{Gr}_{\mathrm{G}}(\mathrm{L})\) to be acyclic, it is necessary and sufficient that \(\mathrm{L}\) be acyclic and that the differentials of \(\mathrm{L}\) be strictly compatible with the filtration \(\mathrm{G}\).}

This is a special case of (1.3.2).

\Page{28}
Let us apply this lemma, taking for \(\mathrm{L}\) the complex
\[
0\to\mathrm{H}(\mathrm{F}^p\mathrm{K})
\to\mathrm{H}(\mathrm{K})
\to\mathrm{H}(\mathrm{K}/\mathrm{F}^p\mathrm{K})\to 0.
\]
We find that:

\noindent
\textit{a)} this sequence is exact, i.e. the differentials of \(\mathrm{K}\) are strictly compatible with the filtration \(\mathrm{F}\);

\noindent
\textit{b)} the filtration of \(\mathrm{H}(\mathrm{F}^p\mathrm{K})\), obtained from the filtration \(\mathrm{W}\) of \(\mathrm{F}^p\mathrm{K}\), coincides with the filtration induced by the filtration \(\mathrm{W}\) of \(\mathrm{H}(\mathrm{K})\);

\noindent
\textit{c)} an analogous statement holds for \(\mathrm{K}/\mathrm{F}^p\mathrm{K}\).

In conclusion:

\textit{Proposition} \textbf{(7.2.8)}. — \textit{If \((\mathrm{K},\mathrm{F},\mathrm{W})\) satisfies \((7.2.2;\,\infty)\), the spectral sequence \(\mathrm{E}(\mathrm{K},\mathrm{F})\) degenerates \((\mathrm{E}_1=\mathrm{E}_{\infty})\). We also have an isomorphism of spectral sequences}
\[
\operatorname{Gr}_{\mathrm{F}}^p(\mathrm{E}_{\bullet}(\mathrm{K},\mathrm{W}))
\simeq\mathrm{E}_{\bullet}(\operatorname{Gr}_{\mathrm{F}}^p\mathrm{K},\mathrm{W}).
\]

\section*{8. HODGE THEORY OF ALGEBRAIC SPACES}

\section*{8.1. Hodge complexes.}

\textbf{(8.1.1)} Let \(\mathrm{A}\) be as in (III.0.3). An \(\mathrm{A}\)-\textit{Hodge complex} \(\mathrm{K}\) \textit{of weight \(n\)} consists of:

\noindent
\(\alpha)\) a complex \(\mathrm{K}_{\mathrm{A}}\in\operatorname{Ob}\mathrm{D}^{+}(\mathrm{A})\), such that \(\mathrm{H}^{k}(\mathrm{K}_{\mathrm{A}})\) is a finitely generated \(\mathrm{A}\)-module for every \(k\);

\noindent
\(\beta)\) a filtration \(\mathrm{F}\) on \(\mathrm{K}_{\mathrm{A}}\otimes\mathbf{C}\), i.e. a filtered complex \((\mathrm{K}_{\mathbf{C}},\mathrm{F})\in\operatorname{Ob}\mathrm{D}^{+}\mathrm{F}(\mathbf{C})\), and an isomorphism \(\alpha:\mathrm{K}_{\mathbf{C}}\simeq\mathrm{K}_{\mathrm{A}}\otimes\mathbf{C}\) in \(\mathrm{D}^{+}(\mathbf{C})\).

The following axioms must be satisfied:

\noindent
(CH 1) the differential \(d\) of \(\mathrm{K}_{\mathbf{C}}\) is strictly compatible with the filtration \(\mathrm{F}\); in other words, the spectral sequence defined by \((\mathrm{K}_{\mathbf{C}},\mathrm{F})\) degenerates at \(\mathrm{E}_1\) \((\mathrm{E}_1=\mathrm{E}_{\infty})\) (1.3.2);

\noindent
(CH 2) for every \(k\), the filtration \(\mathrm{F}\) on \(\mathrm{H}^{k}(\mathrm{K}_{\mathbf{C}})\simeq\mathrm{H}^{k}(\mathrm{K}_{\mathrm{A}})\otimes\mathbf{C}\) defines an \(\mathrm{A}\)-Hodge structure of weight \(n+k\) on \(\mathrm{H}^{k}(\mathrm{K}_{\mathrm{A}})\): the filtration \(\mathrm{F}\) is \((n+k)\)-opposed to its complex conjugate (this makes sense since \(\mathrm{A}\subset\mathbf{R}\)).

\textbf{(8.1.2)} Let \(\mathrm{X}\) be a topological space. An \(\mathrm{A}\)-\textit{cohomological Hodge complex of weight \(n\)} on \(\mathrm{X}\), \(\mathrm{K}\), consists of:

\noindent
\(\alpha)\) a complex \(\mathrm{K}_{\mathrm{A}}\in\operatorname{Ob}\mathrm{D}^{+}(\mathrm{X},\mathrm{A})\),

\noindent
\(\beta)\) a filtered complex \((\mathrm{K}_{\mathbf{C}},\mathrm{F})\in\operatorname{Ob}\mathrm{D}^{+}\mathrm{F}(\mathrm{X},\mathbf{C})\),

\noindent
\(\gamma)\) an isomorphism \(\alpha:\mathrm{K}_{\mathbf{C}}\simeq\mathrm{K}_{\mathrm{A}}\otimes\mathbf{C}\) in \(\mathrm{D}^{+}(\mathrm{X},\mathbf{C})\).

The following axiom must be satisfied:

\noindent
(CHC) The triple \((\mathrm{R}\Gamma(\mathrm{K}_{\mathrm{A}}),\mathrm{R}\Gamma((\mathrm{K}_{\mathbf{C}},\mathrm{F})),\mathrm{R}\Gamma(\alpha))\) is a Hodge complex of weight \(n\).

% ---- 029-036 ----
\Page{29}
When \(\mathrm{A}=\mathbf{Z}\), we shall simply speak of a Hodge complex and a cohomological Hodge complex.

The following statement reformulates part of classical Hodge theory (cf. (2.2.1)).

\textit{Scholium} \textbf{(8.1.3)}. — \textit{Let \(\mathrm{X}\) be a compact Kähler manifold. Let \(\mathrm{K}_{\mathbf{Z}}\) be the complex \(\mathbf{Z}[0]\) consisting of the constant sheaf \(\mathbf{Z}\), \(\mathrm{K}_{\mathbf{C}}\) the holomorphic de Rham complex \(\Omega_{\mathrm{X}}^\bullet\), and \(\mathrm{F}\) its stupid filtration. We have \(\alpha:\mathrm{K}_{\mathbf{Z}}\otimes\mathbf{C}\simeq\mathbf{C}[0]\xrightarrow{\approx}\Omega_{\mathrm{X}}^\bullet\) (Poincaré lemma), and \((\mathrm{K}_{\mathbf{Z}},(\mathrm{K}_{\mathbf{C}},\mathrm{F}),\alpha)\) is a cohomological Hodge complex of weight \(0\).}

\textit{Remark} \textbf{(8.1.4)}. — If \(\mathrm{K}\) is a Hodge complex (resp. a cohomological Hodge complex) of weight \(n\), then \(\mathrm{K}[m]\) is a Hodge complex (resp. a cohomological Hodge complex) of weight \(n+m\).

\textbf{(8.1.5)} An \(\mathrm{A}\)-\textit{mixed Hodge complex} \(\mathrm{K}\) consists of:

\(\alpha)\) a complex \(\mathrm{K}_{\mathrm{A}}\in\operatorname{Ob}\mathrm{D}^{+}(\mathrm{A})\) such that \(\mathrm{H}^{k}(\mathrm{K}_{\mathrm{A}})\) is a finitely generated \(\mathrm{A}\)-module for every \(k\);

\(\beta)\) a filtered complex \((\mathrm{K}_{\mathrm{A}\otimes\mathbf{Q}},\mathrm{W})\in\operatorname{Ob}\mathrm{D}^{+}\mathrm{F}(\mathrm{A}\otimes\mathbf{Q})\) (\(\mathrm{W}\) an \textit{increasing} filtration) and an isomorphism \(\mathrm{K}_{\mathrm{A}\otimes\mathbf{Q}}\simeq\mathrm{K}_{\mathrm{A}}\otimes\mathbf{Q}\) in \(\mathrm{D}^{+}(\mathrm{A}\otimes\mathbf{Q})\);

\(\gamma)\) a bifiltered complex \((\mathrm{K}_{\mathbf{C}},\mathrm{W},\mathrm{F})\in\operatorname{Ob}\mathrm{D}^{+}\mathrm{F}_{2}(\mathbf{C})\) (\(\mathrm{W}\) an \textit{increasing} filtration and \(\mathrm{F}\) a \textit{decreasing} filtration) and an isomorphism \(\alpha:(\mathrm{K}_{\mathbf{C}},\mathrm{W})\simeq(\mathrm{K}_{\mathrm{A}\otimes\mathbf{Q}},\mathrm{W})\otimes\mathbf{C}\) in \(\mathrm{D}^{+}\mathrm{F}(\mathbf{C})\).

The following axiom must be satisfied:

(CHM) For every \(n\), the system consisting of the complex \(\operatorname{Gr}_{n}^{\mathrm{W}}(\mathrm{K}_{\mathrm{A}\otimes\mathbf{Q}})\in\operatorname{Ob}\mathrm{D}^{+}(\mathrm{A}\otimes\mathbf{Q})\), the complex \(\operatorname{Gr}_{n}^{\mathrm{W}}(\mathrm{K}_{\mathbf{C}})\in\operatorname{Ob}\mathrm{D}^{+}\mathrm{F}(\mathbf{C})\) filtered by \(\mathrm{F}\), and the isomorphism
\[
\operatorname{Gr}_{n}^{\mathrm{W}}(\alpha):
\operatorname{Gr}_{n}^{\mathrm{W}}(\mathrm{K}_{\mathrm{A}\otimes\mathbf{Q}})\otimes\mathbf{C}
\simeq
\operatorname{Gr}_{n}^{\mathrm{W}}(\mathrm{K}_{\mathbf{C}}),
\]
is an \(\mathrm{A}\otimes\mathbf{Q}\)-Hodge complex of weight \(n\).

\textbf{(8.1.6)} An \(\mathrm{A}\)-\textit{cohomological mixed Hodge complex} \(\mathrm{K}\) on a topological space \(\mathrm{X}\) consists of:

\(\alpha)\) a complex \(\mathrm{K}_{\mathrm{A}}\in\operatorname{Ob}\mathrm{D}^{+}(\mathrm{X},\mathrm{A})\), such that the \(\mathbf{H}^{k}(\mathrm{X},\mathrm{K}_{\mathrm{A}})\) are finitely generated \(\mathrm{A}\)-modules, and such that \(\mathbf{H}^{\bullet}(\mathrm{X},\mathrm{K}_{\mathrm{A}})\otimes\mathbf{Q}\xrightarrow{\sim}\mathbf{H}^{\bullet}(\mathrm{X},\mathrm{K}_{\mathrm{A}}\otimes\mathbf{Q})\);

\(\beta)\) a filtered complex \((\mathrm{K}_{\mathrm{A}\otimes\mathbf{Q}},\mathrm{W})\in\operatorname{Ob}\mathrm{DF}^{+}(\mathrm{X},\mathrm{A}\otimes\mathbf{Q})\) (\(\mathrm{W}\) an \textit{increasing} filtration) and an isomorphism \(\mathrm{K}_{\mathrm{A}\otimes\mathbf{Q}}\simeq\mathrm{K}_{\mathrm{A}}\otimes\mathbf{Q}\) in \(\mathrm{D}^{+}(\mathrm{X},\mathrm{A}\otimes\mathbf{Q})\);

\(\gamma)\) a bifiltered complex \((\mathrm{K}_{\mathbf{C}},\mathrm{W},\mathrm{F})\in\operatorname{Ob}\mathrm{D}^{+}\mathrm{F}_{2}(\mathrm{X},\mathbf{C})\) (\(\mathrm{W}\) an \textit{increasing} filtration and \(\mathrm{F}\) a \textit{decreasing} filtration) and an isomorphism \(\alpha:(\mathrm{K}_{\mathbf{C}},\mathrm{W})\simeq(\mathrm{K}_{\mathrm{A}\otimes\mathbf{Q}},\mathrm{W})\otimes\mathbf{C}\) in \(\mathrm{D}^{+}\mathrm{F}(\mathrm{X},\mathbf{C})\).

The following axiom must be satisfied:

(CHMC) For every \(n\), the system consisting of the complex \(\operatorname{Gr}_{n}^{\mathrm{W}}(\mathrm{K}_{\mathrm{A}\otimes\mathbf{Q}})\in\operatorname{Ob}\mathrm{D}^{+}(\mathrm{X},\mathrm{A}\otimes\mathbf{C})\), the complex \(\operatorname{Gr}_{n}^{\mathrm{W}}(\mathrm{K}_{\mathbf{C}})\in\operatorname{Ob}\mathrm{D}^{+}\mathrm{F}(\mathrm{X},\mathbf{C})\) filtered by \(\mathrm{F}\), and the isomorphism
\[
\operatorname{Gr}_{n}^{\mathrm{W}}(\alpha):
\operatorname{Gr}_{n}^{\mathrm{W}}(\mathrm{K}_{\mathbf{C}})
\simeq
\operatorname{Gr}_{n}^{\mathrm{W}}(\mathrm{K}_{\mathrm{A}\otimes\mathbf{Q}})\otimes\mathbf{C}
\]
is an \(\mathrm{A}\otimes\mathbf{Q}\)-cohomological Hodge complex of weight \(n\).

\Page{30}
It is trivial to verify:

\textit{Proposition} \textbf{(8.1.7)}. — \textit{If \(\mathrm{K}=(\mathrm{K}_{\mathrm{A}},(\mathrm{K}_{\mathrm{A}\otimes\mathbf{Q}},\mathrm{W}),(\mathrm{K}_{\mathbf{C}},\mathrm{W},\mathrm{F}))\) is an \(\mathrm{A}\)-cohomological mixed Hodge complex, then \(\mathrm{R}\Gamma\mathrm{K}=(\mathrm{R}\Gamma\mathrm{K}_{\mathrm{A}},\mathrm{R}\Gamma(\mathrm{K}_{\mathrm{A}\otimes\mathbf{Q}},\mathrm{W}),\mathrm{R}\Gamma(\mathrm{K}_{\mathbf{C}},\mathrm{W},\mathrm{F}))\) is an \(\mathrm{A}\)-mixed Hodge complex.}

\textit{Scholium} \textbf{(8.1.8)}. — \textit{Let \(\mathrm{U}\) be the complement, in a proper and smooth scheme \(\mathrm{X}\), of a divisor with normal crossings \(\mathrm{D}\), and let \(j:\mathrm{U}\hookrightarrow\mathrm{X}\) be the inclusion morphism. Let \(\mathrm{W}\) be the canonical filtration of \(\mathrm{R}j_{\ast}\mathbf{Q}=\mathrm{R}j_{\ast}\mathbf{Z}\otimes\mathbf{Q}\): \(\mathrm{W}_{n}(\mathrm{R}j_{\ast}\mathbf{Q})=\tau_{\leq n}(\mathrm{R}j_{\ast}\mathbf{Q})\). We have (3.1.8)}
\[
(\mathrm{R}j_{\ast}\mathbf{Q})\otimes\mathbf{C}
\xrightarrow{\approx}\mathrm{R}j_{\ast}\mathbf{C}
\xrightarrow{\approx}j_{\ast}\Omega_{\mathrm{U}}^\bullet
\xleftarrow{\approx}\Omega_{\mathrm{X}}^\bullet(\log\mathrm{D}).
\]
\textit{Let \(\mathrm{W}\) be the weight filtration of \(\Omega_{\mathrm{X}}^\bullet(\log\mathrm{D})\) and \(\mathrm{F}\) its Hodge filtration (by the \(\sigma_{\geq p}\)). Then (3.1.8) gives an isomorphism \((\mathrm{R}j_{\ast}\mathbf{Q},\mathrm{W})\otimes\mathbf{C}\simeq(\Omega_{\mathrm{X}}^\bullet(\log\mathrm{D}),\mathrm{W})\) in \(\mathrm{D}^{+}\mathrm{F}(\mathrm{X},\mathbf{C})\) and, by (3.1.5) and (3.1.9)}
\[
(\mathrm{R}j_{\ast}\mathbf{Z},(\mathrm{R}j_{\ast}\mathbf{Q},\mathrm{W}),
(\Omega_{\mathrm{X}}^\bullet(\log\mathrm{D}),\mathrm{W},\mathrm{F}))
\]
\textit{is a cohomological mixed Hodge complex on \(\mathrm{X}\).}

The following result is an “abstract” version of subsection (3.2).

\textit{Scholium} \textbf{(8.1.9)}. — \textit{Let \(\mathrm{K}\) be an \(\mathrm{A}\)-mixed Hodge complex.}

(i) \textit{On the terms \({}_{\mathrm{W}}\mathrm{E}_{r}^{pq}\) of the spectral sequence of \((\mathrm{K}_{\mathbf{C}},\mathrm{W})\), the recurrent filtration and the two direct filtrations defined by \(\mathrm{F}\) coincide.}

(ii) \textit{The filtration \(\mathrm{W}[n]\) (1.1.2) of \(\mathrm{H}^{n}(\mathrm{K}_{\mathrm{A}\otimes\mathbf{Q}})\simeq\mathrm{H}^{n}(\mathrm{K}_{\mathrm{A}})\otimes\mathbf{Q}\) and the filtration \(\mathrm{F}\) of \(\mathrm{H}^{n}(\mathrm{K}_{\mathbf{C}})\simeq\mathrm{H}^{n}(\mathrm{K}_{\mathrm{A}})\otimes_{\mathrm{A}}\mathbf{C}\) define on \(\mathrm{H}^{n}(\mathrm{K}_{\mathrm{A}})\) an \(\mathrm{A}\)-mixed Hodge structure ((2.3.1) as extended by (III.0.4)).}

(iii) \textit{The morphisms \({}_{\mathrm{W}}d_{1}:{}_{\mathrm{W}}\mathrm{E}_{1}^{pq}\to{}_{\mathrm{W}}\mathrm{E}_{1}^{p+1,q}\) are compatible with the Hodge bigrading; in particular, they are strictly compatible with the Hodge filtration.}

(iv) \textit{The spectral sequence of \((\mathrm{K}_{\mathrm{A}\otimes\mathbf{Q}},\mathrm{W})\) degenerates at \(\mathrm{E}_{2}\): \({}_{\mathrm{W}}\mathrm{E}_{2}={}_{\mathrm{W}}\mathrm{E}_{\infty}\).}

(v) \textit{The spectral sequence of \((\mathrm{K}_{\mathbf{C}},\mathrm{F})\) degenerates at \(\mathrm{E}_{1}\): \({}_{\mathrm{F}}\mathrm{E}_{1}={}_{\mathrm{F}}\mathrm{E}_{\infty}\).}

(vi) \textit{The spectral sequence of the complex \(\operatorname{Gr}_{\mathrm{F}}^{p}(\mathrm{K}_{\mathbf{C}})\), equipped with the filtration induced by \(\mathrm{W}\), degenerates at \(\mathrm{E}_{2}\).}

The proof of (i), (ii), (iii), (iv) is parallel to that of (3.2.5). The analogue of lemmas (3.2.6) and (3.2.7) has been taken as an axiom (CH 1 and CH 2). The proofs of (3.2.8), (3.2.9) and (3.2.10) then carry over unchanged, and as in (3.2.10), we deduce (i), (ii), (iii), (iv). Assertions (v), (vi) then follow from (7.2.8).

\textbf{(8.1.10)} We shall abbreviate “differential graded” to \(\mathrm{DG}\) and “differential graded with degrees bounded below” to \(\mathrm{DG}^{+}\). A complex of \(\mathrm{DG}\) \(\mathrm{A}\)-modules can be viewed as a double complex of \(\mathrm{A}\)-modules; the \textit{first} degree is that of the complex, the \textit{second} is the one defined by the grading of the \(\mathrm{DG}\) modules. We denote by \(\mathrm{D}^{+}((\mathrm{A}\text{-mod }\mathrm{DG})^{+})\)
\Page{31}
the derived category of the category of bounded below complexes of \(\mathrm{DG}\) \(\mathrm{A}\)-modules, with degrees uniformly bounded below.

A \(\mathrm{DG}\) \(\mathrm{A}\)-\textit{mixed} \textit{Hodge complex} \(\mathrm{K}\) consists of:

\textit{a)} a complex \(\mathrm{K}_{\mathrm{A}}\in\operatorname{Ob}\mathrm{D}^{+}((\mathrm{A}\text{-mod }\mathrm{DG})^{+})\);

\textit{b)} a filtered complex \((\mathrm{K}_{\mathrm{A}\otimes\mathbf{Q}},\mathrm{W})\in\operatorname{Ob}\mathrm{D}^{+}\mathrm{F}((\mathrm{A}\otimes\mathbf{Q}\text{-mod }\mathrm{DG})^{+})\), and an isomorphism \(\mathrm{K}_{\mathrm{A}}\otimes\mathbf{Q}\simeq\mathrm{K}_{\mathrm{A}\otimes\mathbf{Q}}\) in \(\mathrm{D}^{+}\);

\textit{c)} a bifiltered complex \((\mathrm{K}_{\mathbf{C}},\mathrm{F},\mathrm{W})\in\operatorname{Ob}\mathrm{D}^{+}\mathrm{F}_{2}((\mathbf{C}\text{-mod }\mathrm{DG})^{+})\), and an isomorphism \(\mathrm{K}_{\mathrm{A}\otimes\mathbf{Q}}\otimes\mathbf{C}\simeq\mathrm{K}_{\mathbf{C}}\) in \(\mathrm{D}^{+}\mathrm{F}\).

We require that, for each \(n\), the component \((\mathrm{K}^{\bullet n},\mathrm{F},\mathrm{W})\) of second degree \(n\) of \(\mathrm{K}\) be an \(\mathrm{A}\)-mixed Hodge complex.

We denote by \(\mathrm{L}\) the filtration of \(\mathbf{s}\mathrm{K}\) by the second degree. We denote by \(\operatorname{Dec}_{1}\mathrm{W}\) the filtration of \(\mathrm{K}_{\mathrm{A}\otimes\mathbf{Q}}\) obtained from \(\mathrm{W}\) by décalage (1.3.3). The filtration it induces on \(\mathbf{s}\mathrm{K}_{\mathrm{A}\otimes\mathbf{Q}}\) (again denoted by \(\operatorname{Dec}_{1}\mathrm{W}\)) must not be confused with the décalage of the filtration induced on \(\mathbf{s}\mathrm{K}_{\mathrm{A}\otimes\mathbf{Q}}\) by \(\mathrm{W}\).

\textit{Variant} \textbf{(8.1.10.1)}. — We similarly define \(\mathrm{A}\)-\textit{cosimplicial mixed Hodge complexes} (resp. \textit{\(r\)-cosimplicial}) by replacing \(\mathrm{DG}\) everywhere by cosimplicial (resp. \(r\)-cosimplicial). The functor (cosimplicial \(\mathrm{A}\)-module) \(\to\) (\(\mathrm{DG}\) \(\mathrm{A}\)-module) takes cosimplicial \(\mathrm{A}\)-mixed Hodge complexes to \(\mathrm{DG}\) \(\mathrm{A}\)-mixed Hodge complexes. Likewise for \(r\)-cosimplicial.

\textbf{(8.1.11)} Let \(\mathrm{X}_{\bullet}\) be a simplicial topological space. An \(\mathrm{A}\)-\textit{cohomological mixed Hodge complex} \(\mathrm{K}\) on \(\mathrm{X}_{\bullet}\) consists of:

\(\alpha)\) a complex \(\mathrm{K}_{\mathrm{A}}\in\operatorname{Ob}\mathrm{D}^{+}(\mathrm{X}_{\bullet},\mathrm{A})\);

\(\beta)\) a filtered complex \((\mathrm{K}_{\mathrm{A}\otimes\mathbf{Q}},\mathrm{W})\in\operatorname{Ob}\mathrm{D}^{+}\mathrm{F}(\mathrm{X}_{\bullet},\mathrm{A}\otimes\mathbf{Q})\), and an isomorphism \(\mathrm{K}_{\mathrm{A}\otimes\mathbf{Q}}\simeq\mathrm{K}_{\mathrm{A}}\otimes\mathbf{Q}\) in \(\mathrm{D}^{+}(\mathrm{X}_{\bullet},\mathrm{A}\otimes\mathbf{Q})\);

\(\gamma)\) a bifiltered complex \((\mathrm{K}_{\mathbf{C}},\mathrm{W},\mathrm{F})\in\operatorname{Ob}\mathrm{D}^{+}\mathrm{F}_{2}(\mathrm{X}_{\bullet},\mathbf{C})\), and an isomorphism
\[
(\mathrm{K}_{\mathbf{C}},\mathrm{W})
\simeq
(\mathrm{K}_{\mathrm{A}\otimes\mathbf{Q}},\mathrm{W})\otimes\mathbf{C}
\]
in \(\mathrm{D}^{+}\mathrm{F}(\mathrm{X}_{\bullet},\mathbf{C})\).

The following axiom must be satisfied:

\((\mathrm{CHM}_{\bullet})\) The restriction of \(\mathrm{K}\) to each of the \(\mathrm{X}_{n}\) is an \(\mathrm{A}\)-cohomological mixed Hodge complex.

For \(\mathrm{A}=\mathbf{Z}\), we shall simply speak of a cohomological mixed Hodge complex.

\textit{Example} \textbf{(8.1.12)}. — Let \(\mathrm{X}_{\bullet}\) be a proper and smooth simplicial scheme (over \(\mathbf{C}\)) and let \(\mathrm{Y}_{\bullet}\) be a divisor with normal crossings in \(\mathrm{X}_{\bullet}\). We set \(\mathrm{U}_{\bullet}=\mathrm{X}_{\bullet}-\mathrm{Y}_{\bullet}\) and denote by \(j\) the inclusion of \(\mathrm{U}_{\bullet}\) in \(\mathrm{X}_{\bullet}\). The constructions (8.1.8) can be carried out “uniform-%
\Page{32}ly in \(\mathrm{X}_{n}\)” as (5.2.5) and (6.2.7) show. They therefore give a cohomological mixed Hodge complex
\[
(\mathrm{R}j_{\ast}\mathbf{Z},(\mathrm{R}j_{\ast}\mathbf{Q},\tau_{\leq}),
(\Omega_{\mathrm{X}_{\bullet}}^\bullet(\log\mathrm{Y}_{\bullet}),\mathrm{W},\mathrm{F}))
\]
on \(\mathrm{X}_{\bullet}\).

\textbf{(8.1.13)} Let \(\mathrm{K}\) be an \(\mathrm{A}\)-cohomological mixed Hodge complex on \(\mathrm{X}_{\bullet}\). Apply the functor \(\mathrm{R}\Gamma^\bullet\) to \(\mathrm{K}\). We obtain

\textit{a)} a cosimplicial complex \(\mathrm{R}\Gamma^\bullet\mathrm{K}_{\mathrm{A}}\in\operatorname{Ob}\mathrm{D}^{+}((\mathrm{A}\text{-cosimplicial modules}))\);

\textit{b)} a filtered cosimplicial complex \(\mathrm{R}\Gamma^\bullet(\mathrm{K}_{\mathrm{A}\otimes\mathbf{Q}},\mathrm{W})\in\operatorname{Ob}\mathrm{D}^{+}\mathrm{F}((\mathrm{A}\otimes\mathbf{Q}\text{-cosimplicial modules}))\);

\textit{c)} a bifiltered cosimplicial complex
\[
\mathrm{R}\Gamma^\bullet(\mathrm{K}_{\mathbf{C}},\mathrm{W},\mathrm{F})
\in\operatorname{Ob}\mathrm{D}^{+}\mathrm{F}_{2}((\mathbf{C}\text{-cosimplicial vector spaces}));
\]

\textit{d)} isomorphisms (8.1.10.\textit{b)}, \textit{c)}) between these objects.

It is clear that \(\mathrm{R}\Gamma^\bullet(\mathrm{K},\mathrm{W},\mathrm{F})\) described above is a cosimplicial \(\mathrm{A}\)-mixed Hodge complex. It defines a \(\mathrm{DG}\) \(\mathrm{A}\)-mixed Hodge complex (8.1.10.1).

\textbf{(8.1.14)} Let \(\mathrm{K}\) be a \(\mathrm{DG}\) \(\mathrm{A}\)-mixed Hodge complex, and consider the spectral sequence with abutment the cohomology of \(\mathbf{s}\mathrm{K}\) defined by the filtration \(\mathrm{L}\). By hypothesis, the \(\mathrm{E}_{1}\) terms of this spectral sequence are equipped with \(\mathrm{A}\)-mixed Hodge structures. The content of the theorem below is that the entire spectral sequence, including its abutment, is equipped with an \(\mathrm{A}\)-mixed Hodge structure. In other words, it is a spectral sequence in the abelian category of \(\mathrm{A}\)-mixed Hodge structures. The structure on the abutment corresponds to a natural “structure” of an \(\mathrm{A}\)-mixed Hodge complex on \(\mathbf{s}\mathrm{K}\).

\textit{Theorem} \textbf{(8.1.15)}. — \textit{Let \(\mathrm{K}\) be a \(\mathrm{DG}\) \(\mathrm{A}\)-mixed Hodge complex.}

(i) \textit{With the notation of (8.1.10), (8.1.10.1) and (7.1.6), \((\mathbf{s}\mathrm{K},\delta(\mathrm{W},\mathrm{L}),\mathrm{F})\) is an \(\mathrm{A}\)-mixed Hodge complex.}

(ii) \textit{We have}
\[
{}_{\delta(\mathrm{W},\mathrm{L})}\mathrm{E}_{1}^{ab}
(\mathbf{s}\mathrm{K}\otimes\mathbf{Q})
=
\sum_{\substack{a=\gamma-\beta\\b=\alpha+\beta}}
\mathrm{H}^{\alpha}
(\operatorname{Gr}_{\beta}^{\mathrm{W}}\mathrm{K}^{\bullet,\gamma});
\tag*{\textbf{(8.1.15.1)}}
\]
\textit{the complex \(({}_{\delta(\mathrm{W},\mathrm{L})}\mathrm{E}_{1}^{\bullet b},d_{1})\) is the simple complex associated to the following double complex of \(\mathrm{A}\otimes\mathbf{Q}\)-Hodge structures of weight \(b\)}
\[
\begin{tikzcd}[ampersand replacement=\&, column sep=1em, row sep=3em]
\mathrm{H}^{b-(\beta+1)}
(\operatorname{Gr}_{\beta+1}^{\mathrm{W}}\mathrm{K}^{\bullet,\gamma+1})
\arrow[r, "\partial"]
\&
\mathrm{H}^{b-\beta}
(\operatorname{Gr}_{\beta}^{\mathrm{W}}\mathrm{K}^{\bullet,\gamma+1})
\arrow[r, "\partial"]
\&
\mathrm{H}^{b-(\beta-1)}
(\operatorname{Gr}_{\beta-1}^{\mathrm{W}}\mathrm{K}^{\bullet,\gamma+1})
\\
\&
\mathrm{H}^{b-\beta}
(\operatorname{Gr}_{\beta}^{\mathrm{W}}\mathrm{K}^{\bullet,\gamma})
\arrow[u, "d''"]
\arrow[r, "\partial"]
\&
\mathrm{H}^{b-(\beta-1)}
(\operatorname{Gr}_{\beta-1}^{\mathrm{W}}\mathrm{K}^{\bullet,\gamma})
\arrow[u, "d''"]
\end{tikzcd}
\tag*{\textbf{(8.1.15.2)}}
\]

\Page{33}
(iii) \textit{On the terms \({}_{\mathrm{L}}\mathrm{E}_{r}\) of the spectral sequence defined by \((\mathbf{s}\mathrm{K}_{\mathrm{A}\otimes\mathbf{Q}},\mathrm{L})\), the recurrent filtration and the two direct filtrations defined by \(\operatorname{Dec}_{1}(\mathrm{W})\) coincide; likewise, on the \({}_{\mathrm{L}}\mathrm{E}_{r}(\mathbf{s}\mathrm{K}_{\mathbf{C}})\), the recurrent filtration and the two direct filtrations induced by \(\mathrm{F}\) coincide. We shall denote these filtrations by \(\mathrm{W}\) and \(\mathrm{F}\).}

(iv) \textit{For \(r\geq1\), \(({}_{\mathrm{L}}\mathrm{E}_{r},\mathrm{W},\mathrm{F})\) is an \(\mathrm{A}\)-mixed Hodge structure (III.0.4). The differentials \(d_{r}\) are morphisms of \(\mathrm{A}\)-mixed Hodge structures.}

(v) \textit{The filtration \(\mathrm{L}\) of \(\mathrm{H}^{\bullet}(\mathbf{s}\mathrm{K})\) is a filtration in the category of mixed Hodge structures, and}
\[
\operatorname{Gr}_{\mathrm{L}}^{p}
(\mathrm{H}^{p+q}(\mathbf{s}\mathrm{K}),
\delta(\mathrm{W},\mathrm{L})[p+q],\mathrm{F})
=
({}_{\mathrm{L}}\mathrm{E}_{\infty}^{pq},\mathrm{W},\mathrm{F}).
\tag*{\textbf{(8.1.15.3)}}
\]

\textit{Remark} \textbf{(8.1.16)}. — In (8.1.15), the filtration \(\operatorname{Dec}_{1}(\mathrm{W})\) (8.1.10) plays a greater role than \(\mathrm{W}\). We have
\[
\operatorname{Dec}(\delta(\mathrm{W},\mathrm{L}))
=
\operatorname{Dec}_{1}(\mathrm{W})
\quad\text{on }\mathbf{s}\mathrm{K},
\tag*{\textbf{(8.1.16.1)}}
\]
and (8.1.15.3) can be rewritten as
\[
\operatorname{Gr}_{\mathrm{L}}
(\mathrm{H}(\mathrm{K}),\operatorname{Dec}_{1}\mathrm{W},\mathrm{F})
=
({}_{\mathrm{L}}\mathrm{E}_{\infty},\mathrm{W},\mathrm{F}).
\tag*{\textbf{(8.1.16.2)}}
\]
In the same spirit, note that, with the notation of (8.1.8)
\[
\mathrm{R}\Gamma
(\mathrm{R}j_{\ast}\mathbf{Z},
(\mathrm{R}j_{\ast}\mathbf{Q},\operatorname{Dec}\mathrm{W}),
(\Omega_{\mathrm{X}}^\bullet(\log\mathrm{D}),
\operatorname{Dec}\mathrm{W},\mathrm{F}))
\]
depends only on \(\mathrm{U}\) and not on the chosen compactification \(\mathrm{X}\).

\textbf{(8.1.17)} Let us prove (8.1.15). Assertions (i) and (ii) follow from the fact that
\[
\operatorname{Gr}_{n}^{\delta(\mathrm{W},\mathrm{L})}
(\mathbf{s}\mathrm{K}\otimes\mathbf{Q})
=
\sum_{p-q=n}
\operatorname{Gr}_{p}^{\mathrm{W}}\operatorname{Gr}_{\mathrm{L}}^{q}
(\mathbf{s}\mathrm{K}\otimes\mathbf{Q})
=
\sum_{p-q=n}
(\operatorname{Gr}_{p}^{\mathrm{W}}\mathrm{K}^{\bullet q})[-q]
\]
(this isomorphism is compatible with \(\mathrm{F}\), and we apply (8.1.4)). The verification of (8.1.15.2) is left to the reader.

Consider the assertions:

\((a_{r})\) The morphisms \(d_{r}\) of \({}_{\mathrm{L}}\mathrm{E}_{r}(\mathbf{s}\mathrm{K}_{\mathrm{A}\otimes\mathbf{Q}})\) are strictly compatible with the recurrent filtration defined by \(\operatorname{Dec}_{1}(\mathrm{W})\).

\((b_{r})\) The morphisms \(d_{r}\) of \({}_{\mathrm{L}}\mathrm{E}_{r}(\mathbf{s}\mathrm{K}_{\mathbf{C}})\) are strictly compatible with the recurrent filtration defined by \(\mathrm{F}\).

\((c_{r})\) \({}_{\mathrm{L}}\mathrm{E}_{r}\), equipped with the filtrations considered in \((a_{r})\), \((b_{r})\), is an \(\mathrm{A}\)-mixed Hodge structure.

\((d_{r})\) The \(d_{r}\) are morphisms of \(\mathrm{A}\)-mixed Hodge structures.

We shall prove \((a_{r})\) (\(r\geq0\)), \((b_{r})\) (\(r\geq0\)), \((c_{r})\) (\(r\geq1\)), \((d_{r})\) (\(r\geq1\)) by simultaneous induction.

A) \textit{Proof of} \((a_{0})\): we have
\[
({}_{\mathrm{L}}\mathrm{E}_{0}^{p\bullet}
(\mathbf{s}\mathrm{K}_{\mathrm{A}\otimes\mathbf{Q}}),
\operatorname{Dec}_{1}\mathrm{W})
=
(\mathrm{K}_{\mathrm{A}\otimes\mathbf{Q}}^{\bullet p},
\operatorname{Dec}\mathrm{W}).
\]
By (8.1.9) (iv), the spectral sequence of \((\mathrm{K}_{\mathrm{A}\otimes\mathbf{Q}}^{\bullet p},\mathrm{W})\) degenerates at \(\mathrm{E}_{2}\). By (1.3.4), the spectral sequence of \((\mathrm{K}_{\mathrm{A}\otimes\mathbf{Q}}^{\bullet p},\operatorname{Dec}\mathrm{W})\) degenerates at \(\mathrm{E}_{1}\), and this proves \((a_{0})\) (1.3.2).

\Page{34}
B) \textit{Proof of} \((b_{0})\): we have
\[
({}_{\mathrm{L}}\mathrm{E}_{0}^{p\bullet}
(\mathbf{s}\mathrm{K}_{\mathbf{C}}),\mathrm{F})
=
(\mathrm{K}_{\mathbf{C}}^{\bullet p},\mathrm{F}).
\]
We apply (8.1.9) (v) and (1.3.2).

C) \textit{Proof of} \((c_{1})\): we must show that \(\mathrm{H}^{n}(\mathrm{K}_{\mathrm{A}}^{\bullet p})\), equipped with the filtrations induced by \(\operatorname{Dec}\mathrm{W}\) on \(\mathrm{H}^{n}(\mathrm{K}_{\mathrm{A}\otimes\mathbf{Q}}^{\bullet p})=\mathrm{H}^{n}(\mathrm{K}_{\mathrm{A}}^{\bullet p})\otimes\mathbf{Q}\) and \(\mathrm{F}\) on \(\mathrm{H}^{n}(\mathrm{K}_{\mathbf{C}}^{\bullet p})=\mathrm{H}^{n}(\mathrm{K}_{\mathrm{A}}^{\bullet p})\otimes_{\mathrm{A}}\mathbf{C}\), is an \(\mathrm{A}\)-mixed Hodge structure. \(\operatorname{Dec}\mathrm{W}\) induces on \(\mathrm{H}^{n}(\mathrm{K}_{\mathrm{A}\otimes\mathbf{Q}}^{\bullet p})\) the same filtration as \(\mathrm{W}[n]\), and we conclude by (8.1.9) (ii).

D) \textit{Proof of} \((c_{r})+(d_{r})\Rightarrow(a_{r})+(b_{r})+(c_{r+1})\): this is the key point of the proof, and a simple application of (2.3.5) (as extended by (III.0.4)).

E) \textit{Proof of} \((a_{r})_{r\leq r_{0}-2}+(b_{r})_{r\leq r_{0}-2}+(c_{r_{0}})\Rightarrow(d_{r_{0}})\): By (1.3.16), the recurrent filtration and the direct filtrations induced by \(\operatorname{Dec}_{1}\mathrm{W}\) (resp. \(\mathrm{F}\)) on \({}_{\mathrm{L}}\mathrm{E}_{r}\) coincide for \(r\leq r_{0}\), and we apply (1.3.13) (i).

This completes the proof by induction of \((a)\), \((b)\), \((c)\), \((d)\).

Assertion (iii) was proved in part E of the induction, and the assertions of (iv) are \((c_{r})\) and \((d_{r})\). We shall prove (8.1.15.3) in the form (8.1.16.2).

The equality
\[
\operatorname{Gr}_{\mathrm{L}}
(\mathrm{H}(\mathbf{s}\mathrm{K}_{\mathrm{A}\otimes\mathbf{Q}}),
\operatorname{Dec}_{1}\mathrm{W})
=
({}_{\mathrm{L}}\mathrm{E}_{\infty}
(\mathbf{s}\mathrm{K}_{\mathrm{A}\otimes\mathbf{Q}}),\mathrm{W})
\]
follows from \((a_{r})\) and (1.3.17). The equality
\[
\operatorname{Gr}_{\mathrm{L}}
(\mathrm{H}(\mathbf{s}\mathrm{K}_{\mathbf{C}}),\mathrm{F})
=
({}_{\mathrm{L}}\mathrm{E}_{\infty}
(\mathbf{s}\mathrm{K}_{\mathbf{C}}),\mathrm{F})
\]
follows from \((b_{r})\) and (1.3.17). We conclude by the following lemma, whose verification is left to the reader:

\textit{Lemma} \textbf{(8.1.18)}. — \textit{Let \(\mathrm{H}=(\mathrm{H}_{\mathrm{A}},\mathrm{W},\mathrm{F})\) be an \(\mathrm{A}\)-mixed Hodge structure. For a filtration \(\mathrm{L}\) of \(\mathrm{H}_{\mathrm{A}}\) to arise from a filtration of \(\mathrm{H}\), in the abelian category of \(\mathrm{A}\)-mixed Hodge structures, it is necessary and sufficient that for every \(n\), \((\operatorname{Gr}_{\mathrm{L}}^{n}(\mathrm{H}_{\mathrm{A}}),\operatorname{Gr}_{\mathrm{L}}^{n}(\mathrm{W}),\operatorname{Gr}_{\mathrm{L}}^{n}(\mathrm{F}))\) be an \(\mathrm{A}\)-mixed Hodge structure.}

\textbf{(8.1.19)} Let us apply (8.1.13) and (8.1.15) to example (8.1.12).

We obtain a mixed Hodge complex consisting of \textit{a)}, \textit{b)}, \textit{c)} below.

\textit{a)} The complex \(\mathrm{R}\Gamma\mathrm{R}j_{\ast}\mathbf{Z}\simeq\mathrm{R}\Gamma(\mathrm{U}_{\bullet},\mathbf{Z})\), whose cohomology groups are the \(\mathrm{H}^{n}(\mathrm{U}_{\bullet},\mathbf{Z})\).

\textit{b)} The filtration \(\delta(\mathrm{W},\mathrm{L})\) of \(\mathrm{R}\Gamma j_{\ast}\mathbf{Q}\simeq\mathrm{R}\Gamma(\mathrm{U}_{\bullet},\mathbf{Q})\). With the notation of (3.1.4), we have
\begin{align*}
\operatorname{Gr}_{n}^{\delta(\mathrm{W},\mathrm{L})}
\mathrm{R}\Gamma(\mathrm{U}_{\bullet},\mathbf{Q})
&\underset{\text{(7.1.6.5)}}{\simeq}
\bigoplus_m
\operatorname{Gr}_{n+m}^{\mathrm{W}}
\mathrm{R}\Gamma(\mathrm{U}_{m},\mathbf{Q})[-m]
\\
&\simeq
\bigoplus_m
\mathrm{R}\Gamma
(\widetilde{\mathrm{Y}}_{m}^{n+m},
\varepsilon_{\mathbf{Q}}^{n+m}[-n-m])[-m]
\\
&\simeq
\bigoplus_m
\mathrm{R}\Gamma
(\widetilde{\mathrm{Y}}_{m}^{n+m},
\varepsilon_{\mathbf{Q}}^{n+m})[-n-2m].
\end{align*}

\Page{35}
The \(\mathrm{E}_{1}\) term of the corresponding spectral sequence is a sum of groups of the form \(\mathrm{H}^{p}(\widetilde{\mathrm{Y}}_{q}^{r},\varepsilon_{\mathbf{Q}}^{r})\); the latter contributes to \(\mathrm{H}^{p+q+r}(\mathrm{U}_{\bullet})\); it corresponds to \(\operatorname{Gr}_{a}^{\delta(\mathrm{W},\mathrm{F})}\) for \(p+2r=(p+q+r)+a\):
\[
{}_{\delta(\mathrm{W},\mathrm{L})}\mathrm{E}_{1}^{-a,b}
=
\bigoplus_{\substack{p+2r=b\\q-r=-a}}
\mathrm{H}^{p}(\widetilde{\mathrm{Y}}_{q}^{r},\varepsilon_{\mathbf{Q}}^{r})
\Rightarrow
\mathrm{H}^{-a+b}(\mathrm{U}_{\bullet},\mathbf{Q}).
\tag*{\textbf{(8.1.19.1)}}
\]

\textit{c)} The bifiltration \((\delta(\mathrm{W},\mathrm{L}),\mathrm{F})\) of \(\mathrm{R}\Gamma\mathrm{R}j_{\ast}\mathbf{C}\simeq\mathbf{s}\mathrm{R}\Gamma^\bullet\Omega_{\mathrm{X}}^\bullet(\log\mathrm{D})\).

The filtration \(\mathrm{F}\) equips \({}_{\delta(\mathrm{W},\mathrm{L})}\mathrm{E}_{1}^{-a,b}\) with a Hodge structure of weight \(b\), and the differentials \(d_{1}\) of \({}_{\delta(\mathrm{W},\mathrm{L})}\mathrm{E}\) are morphisms of \(\mathbf{Q}\)-Hodge structures. The complex \({}_{\delta(\mathrm{W},\mathrm{L})}\mathrm{E}_{1}^{\bullet b}\) is the simple complex associated to the following double complex, whose components are Hodge structures of weight \(b\):

{\small
\[
\begin{tikzcd}[ampersand replacement=\&, column sep=1em, row sep=4em]
\mathrm{H}^{b-2(r+2)}
(\widetilde{\mathrm{Y}}_{q+1}^{r+1},\varepsilon_{\mathbf{Q}}^{r+1})(-r-1)
\arrow[r, "\mathrm{Gysin}"']
\&
\mathrm{H}^{b-2r}
(\widetilde{\mathrm{Y}}_{q+1}^{r},\varepsilon_{\mathbf{Q}}^{r})(-r)
\arrow[r, "\mathrm{Gysin}"']
\&
\mathrm{H}^{b-2(r-2)}
(\widetilde{\mathrm{Y}}_{q+1}^{r-1},\varepsilon_{\mathbf{Q}}^{r-1})(-r+1)
\\
\&
\mathrm{H}^{b-2r}
(\widetilde{\mathrm{Y}}_{q}^{r},\varepsilon_{\mathbf{Q}}^{r})(-r)
\arrow[u, "{\sum\limits_i(-1)^i\delta_i}"']
\arrow[r, "\mathrm{Gysin}"']
\&
\mathrm{H}^{b-2(r-2)}
(\widetilde{\mathrm{Y}}_{q}^{r-1},\varepsilon_{\mathbf{Q}}^{r-1})(-r+1)
\arrow[u, "{\sum\limits_i(-1)^i\delta_i}"']
\end{tikzcd}
\]
}

The indices of the components in \(\mathrm{E}_{1}^{ab}\) satisfy \(-a=q-r\).

In particular, the \(\mathrm{H}^{n}(\mathrm{U}_{\bullet},\mathbf{Z})\) are thus equipped with mixed Hodge structures (8.1.9).

\textit{Proposition} \textbf{(8.1.20)}. — \textit{Under the preceding hypotheses:}

(i) \textit{The mixed Hodge structure on \(\mathrm{H}^{n}(\mathrm{U}_{\bullet},\mathbf{Z})\) is functorial in the pair \((\mathrm{U}_{\bullet},\mathrm{X}_{\bullet})\).}

(ii) \textit{In rational cohomology, the spectral sequence (8.1.19.1) degenerates at \(\mathrm{E}_{2}\) \((\mathrm{E}_{2}=\mathrm{E}_{\infty})\) and abuts to the weight filtration of \(\mathrm{H}^{n}(\mathrm{U}_{\bullet},\mathbf{Q})\). The Hodge structure on \(\mathrm{E}_{2}\) induced by that on \(\mathrm{E}_{1}\) is also that on \(\operatorname{Gr}_{\mathrm{W}}\mathrm{H}^{n}(\mathrm{U}_{\bullet},\mathbf{Q})\).}

(iii) \textit{The Hodge numbers of \(\mathrm{H}^{n}(\mathrm{U}_{\bullet},\mathbf{Z})\) satisfy}
\[
h^{pq}\neq0
\quad\Rightarrow\quad
0\leq p\leq n\ \textit{and}\ 0\leq q\leq n.
\]

(iv) \textit{For \(\mathrm{Y}_{\bullet}=\varnothing\), the Hodge numbers of \(\mathrm{H}^{n}(\mathrm{X}_{\bullet},\mathbf{Z})\) satisfy}
\[
h^{pq}\neq0
\quad\Rightarrow\quad
0\leq p\leq n,\quad 0\leq q\leq n
\quad\textit{and}\quad p+q\leq n.
\]

(i) and (ii) follow from the general theory (8.1.9). To prove (iii), we examine (8.1.19.2), and observe that the Hodge structures \(\mathrm{H}^{p}(\widetilde{\mathrm{Y}}_{q}^{r},\varepsilon_{\mathbf{Q}}^{r})(-r)\) \((p+q+r=n)\) satisfy (iii). To prove (iv), we observe that the Hodge structures \(\mathrm{H}^{p}(\mathrm{X}_{q})\) \((p+q=n)\) satisfy (iv).

\textbf{(8.1.21)} We define an \(\mathrm{A}\)-\textit{cohomological mixed Hodge complex on an \(r\)-simplicial topological space} by adapting (8.1.11). If \(\mathrm{K}\) is such a complex, we see as in (8.1.13) that \(\mathrm{R}\Gamma^\bullet\mathrm{K}\) is an \(r\)-cosimplicial \(\mathrm{A}\)-mixed Hodge complex (8.1.10.1).

\Page{36}
Let \(\mathrm{X}_{\bullet}\) be a proper and smooth \(r\)-simplicial scheme. A divisor with normal crossings \(\mathrm{D}_{\bullet}\) in \(\mathrm{X}_{\bullet}\) is a system of divisors with normal crossings in the \(\mathrm{X}_{n}\) such that \(\mathrm{X}_{\bullet}-\mathrm{D}_{\bullet}=\mathrm{U}_{\bullet}\) is an \(r\)-simplicial subscheme of \(\mathrm{X}_{\bullet}\). Let \(j:\mathrm{U}_{\bullet}\hookrightarrow\mathrm{X}_{\bullet}\). As in (8.1.12), we define a cohomological mixed Hodge complex \(\mathrm{R}j_{\ast}\mathbf{Z}\) on \(\mathrm{X}_{\bullet}\).

As in (8.1.20), we could deduce from this a mixed Hodge structure on \(\mathrm{H}^{\bullet}(\mathrm{U}_{\bullet},\mathbf{Z})\), but, since
\[
\mathrm{H}^{\bullet}(\mathrm{U}^{\bullet},\mathbf{Z})
=
\mathrm{H}^{\bullet}(\delta\mathrm{U}_{\bullet},\mathbf{Z}),
\]
this would tell us nothing new (the preceding equality would be an identity between mixed Hodge structures).

\textbf{(8.1.22)} Set \(r=2\). For any bounded below bicosimplicial complex \(\mathrm{K}=(\mathrm{K}^{p n_{1} n_{2}})\) (\(p=\text{differential degree}\), \(n_i=i\)-th simplicial degree), we shall denote by \(\mathbf{s}_{1}\mathrm{K}\) the cosimplicial complex with simplicial degree \(n_{2}\), given by
\[
(\mathbf{s}_{1}\mathrm{K})^{\bullet n_{2}}
=
\mathbf{s}\mathrm{K}^{\bullet\bullet n_{2}}.
\]
We shall denote by \(\mathrm{L}_{1}\) the filtration of \(\mathbf{s}_{1}\mathrm{K}\) which, on each \(\mathbf{s}\mathrm{K}^{\bullet\bullet n_{2}}\), induces the filtration by the simplicial degree \(n_{1}\).

Let \(\mathrm{K}\) be a bicosimplicial \(\mathrm{A}\)-mixed Hodge complex. It follows from (8.1.15) that \((\mathbf{s}_{1}\mathrm{K},\delta(\mathrm{W},\mathrm{L}_{1}),\mathrm{F})\) is a cosimplicial mixed Hodge complex. Let \(\mathrm{L}_{2}\) be the filtration of \(\mathbf{s}\mathrm{K}=\mathbf{s}\mathbf{s}_{1}\mathrm{K}\) by the simplicial degree \(n_{2}\). By (8.1.14), (8.1.15), the spectral sequence defined by \(\mathrm{L}_{2}\) arises from a spectral sequence in the abelian category of \(\mathrm{A}\)-mixed Hodge structures
\[
{}_{\mathrm{L}_{2}}\mathrm{E}_{1}^{pq}
=
\mathrm{H}^{q}(\mathbf{s}\mathrm{K}^{\bullet\bullet p})
\Rightarrow
\mathrm{H}^{p+q}(\mathbf{s}\mathrm{K}).
\tag*{\textbf{(8.1.22.1)}}
\]
In this spectral sequence, the \(\mathrm{A}\)-mixed Hodge structure on \(\mathrm{E}_{1}^{pq}\) is the \(\mathrm{A}\)-mixed Hodge structure ((8.1.9) (ii)) on the cohomology of the \(\mathrm{A}\)-mixed Hodge complex \((\mathbf{s}\mathrm{K}^{\bullet\bullet p},\delta(\mathrm{W},\mathrm{L}_{1}),\mathrm{F})\). The \(\mathrm{A}\)-mixed Hodge structure on the abutment is the \(\mathrm{A}\)-mixed Hodge structure on the cohomology of the \(\mathrm{A}\)-mixed Hodge complex \((\mathbf{s}\mathrm{K},\delta(\delta(\mathrm{W},\mathrm{L}_{1}),\mathrm{L}_{2}),\mathrm{F})\). The filtration \(\delta(\delta(\mathrm{W},\mathrm{L}_{1}),\mathrm{L}_{2})\underset{\mathrm{dfn}}{=}\delta(\mathrm{W},\mathrm{L}_{1},\mathrm{L}_{2})\) is given by
\[
\delta(\mathrm{W},\mathrm{L}_{1},\mathrm{L}_{2})_{k}
(\mathrm{K}^{p n_{1} n_{2}})
=
\mathrm{W}_{k+n_{1}+n_{2}}\mathrm{K}^{p n_{1} n_{2}}.
\]

\textbf{(8.1.23)} In particular, an \(\mathrm{A}\)-cohomological mixed Hodge complex on a bisimplicial topological space \(\mathrm{X}_{\bullet\bullet}\) defines a spectral sequence of \(\mathrm{A}\)-mixed Hodge structures
\[
\mathrm{E}_{1}^{pq}
=
\mathrm{H}^{q}(\mathrm{X}_{\bullet p},\mathrm{K})
\Rightarrow
\mathrm{H}^{p+q}(\mathrm{X}_{\bullet\bullet},\mathrm{K}).
\tag*{\textbf{(8.1.23.1)}}
\]
In the special case considered in (8.1.21), for \(r=2\), we have \(\mathrm{A}=\mathbf{Z}\) and (8.1.23.1) reads
\[
\mathrm{E}_{1}^{pq}
=
\mathrm{H}^{q}(\mathrm{U}_{\bullet p},\mathbf{Z})
\Rightarrow
\mathrm{H}^{p+q}(\mathrm{U}_{\bullet\bullet},\mathbf{Z}).
\tag*{\textbf{(8.1.23.2)}}
\]

% ---- 037-044 ----
\Page{37}
\textbf{(8.1.24)} If \(\mathrm{K}'\) and \(\mathrm{K}''\) are two \(\mathrm{A}\)-mixed Hodge complexes, their tensor product \(\mathrm{K}'\otimes\mathrm{K}''\), defined by
\[
\begin{aligned}
&(\mathrm{K}'\otimes\mathrm{K}'')_{\mathrm{A}}
 =\mathrm{K}'_{\mathrm{A}}\overset{\mathrm{L}}{\otimes}\mathrm{K}''_{\mathrm{A}},\\
&((\mathrm{K}'\otimes\mathrm{K}'')_{\mathrm{A}\otimes\mathbf{Q}},\mathrm{W})
 =(\mathrm{K}'_{\mathrm{A}\otimes\mathbf{Q}},\mathrm{W}')
 \otimes(\mathrm{K}''_{\mathrm{A}\otimes\mathbf{Q}},\mathrm{W}''),\\
&((\mathrm{K}'\otimes\mathrm{K}'')_{\mathbf{C}},\mathrm{W},\mathrm{F})
 =(\mathrm{K}'_{\mathbf{C}},\mathrm{W}',\mathrm{F}')
 \otimes(\mathrm{K}''_{\mathbf{C}},\mathrm{W}'',\mathrm{F}''),
\end{aligned}
\]
is again an \(\mathrm{A}\)-mixed Hodge complex.

Let \(\mathrm{U}'\) (resp. \(\mathrm{U}''\)) be the complement of a divisor with normal crossings \(\mathrm{Y}'\) (resp. \(\mathrm{Y}''\)) in a proper and smooth scheme \(\mathrm{X}'\) (resp. \(\mathrm{X}''\)). Then, \(\mathrm{U}=\mathrm{U}'\times\mathrm{U}''\) is the complement of a divisor with normal crossings \(\mathrm{Y}\) in \(\mathrm{X}=\mathrm{X}'\times\mathrm{X}''\). Let \(j\) (resp. \(j',j''\)) be the inclusion of \(\mathrm{U}\) (resp. \(\mathrm{U}',\mathrm{U}''\)) in \(\mathrm{X}\) (resp. \(\mathrm{X}',\mathrm{X}''\)).

We have a quasi-isomorphism
\[
\mathrm{R}j'_*\mathbf{Z}\overset{\mathrm{L}}{\boxtimes}\mathrm{R}j''_*\mathbf{Z}
\underset{\mathrm{dfn}}{=}
\operatorname{pr}_1^*\mathrm{R}j'_*\mathbf{Z}
\overset{\mathrm{L}}{\otimes}
\operatorname{pr}_2^*\mathrm{R}j''_*\mathbf{Z}
\xrightarrow{\sim}\mathrm{R}j_*\mathbf{Z},
\]
a filtered quasi-isomorphism
\[
(\mathrm{R}j'_*\mathbf{Q},\tau_{\leq})
\overset{\mathrm{L}}{\boxtimes}
(\mathrm{R}j''_*\mathbf{Q},\tau_{\leq})
\xrightarrow{\sim}
(\mathrm{R}j_*\mathbf{Q},\tau_{\leq}),
\]
and a bifiltered morphism
\[
(\Omega_{\mathrm{X}'}^\bullet(\log\mathrm{Y}'),\mathrm{W},\mathrm{F})
\boxtimes
(\Omega_{\mathrm{X}''}^\bullet(\log\mathrm{Y}''),\mathrm{W},\mathrm{F})
\longrightarrow
(\Omega_{\mathrm{X}}^\bullet(\log\mathrm{Y}),\mathrm{W},\mathrm{F}).
\]

Let us apply the functor \(\mathrm{R}\Gamma\); by the Künneth formula, and by its analogue in the cohomology of coherent analytic sheaves, we obtain an isomorphism of mixed Hodge complexes
\[
\mathrm{R}\Gamma(\mathrm{U}',\mathbf{Z})
\overset{\mathrm{L}}{\otimes}
\mathrm{R}\Gamma(\mathrm{U}'',\mathbf{Z})
\xrightarrow{\sim}
\mathrm{R}\Gamma(\mathrm{U},\mathbf{Z}),
\]
whence an isomorphism of \(\mathbf{Q}\)-mixed Hodge structures
\[
\mathrm{H}^\bullet(\mathrm{U}',\mathbf{Q})
\otimes\mathrm{H}^\bullet(\mathrm{U}'',\mathbf{Q})
\xrightarrow{\sim}\mathrm{H}^\bullet(\mathrm{U},\mathbf{Q}).
\]

\textbf{(8.1.25)} Standard arguments, based on the Cartier-Eilenberg-Zilber theorem, yield the following simplicial variant of (8.1.24).

Let \(\mathrm{U}'_\bullet\) (resp. \(\mathrm{U}''_\bullet\)) be the complement of a divisor with normal crossings in a proper and smooth simplicial scheme \(\mathrm{X}'_\bullet\) (resp. \(\mathrm{X}''_\bullet\)). Let \(\mathrm{U}_\bullet=\mathrm{U}'_\bullet\times\mathrm{U}''_\bullet\) and \(\mathrm{X}_\bullet=\mathrm{X}'_\bullet\times\mathrm{X}''_\bullet\).

We have an isomorphism of mixed Hodge complexes
\[
\mathrm{R}\Gamma(\mathrm{U}'_\bullet,\mathbf{Z})
\overset{\mathrm{L}}{\otimes}
\mathrm{R}\Gamma(\mathrm{U}''_\bullet,\mathbf{Z})
\xrightarrow{\sim}
\mathrm{R}\Gamma(\mathrm{U}_\bullet,\mathbf{Z}),
\]
and in particular an isomorphism of \(\mathbf{Q}\)-mixed Hodge structures
\[
\mathrm{H}^\bullet(\mathrm{U}'_\bullet,\mathbf{Q})
\otimes\mathrm{H}^\bullet(\mathrm{U}''_\bullet,\mathbf{Q})
\xrightarrow{\sim}\mathrm{H}^\bullet(\mathrm{U}_\bullet,\mathbf{Q}).
\]

\textbf{(8.1.26)} Let \(\mathrm{U}'_{\bullet\bullet}\) (resp. \(\mathrm{U}''_{\bullet\bullet}\)) be the complement of a divisor with normal crossings in a proper and smooth bisimplicial scheme \(\mathrm{X}'_{\bullet\bullet}\) (resp. \(\mathrm{X}''_{\bullet\bullet}\)). Let \(\mathrm{U}_{\bullet\bullet}=\mathrm{U}'_{\bullet\bullet}\times\mathrm{U}''_{\bullet\bullet}\) and
\Page{38}
\(\mathrm{X}_{\bullet\bullet}=\mathrm{X}'_{\bullet\bullet}\times\mathrm{X}''_{\bullet\bullet}\). Standard arguments show that, in rational cohomology, the spectral sequence (8.1.23.2), for \(\mathrm{U}_{\bullet\bullet}\), is the tensor product of the spectral sequences (8.1.23.2) for \(\mathrm{U}'_{\bullet\bullet}\) and \(\mathrm{U}''_{\bullet\bullet}\).

\section*{8.2. Separated algebraic spaces.}

\textbf{(8.2.1)} Let \(\mathrm{X}\) be a scheme (or an algebraic space) of finite type over \(\mathbf{C}\), assumed to be separated. There is then a diagram
\[
\begin{tikzcd}[column sep=large,row sep=large]
\mathrm{Y}_\bullet \arrow[r,hook] \arrow[d,"\alpha"']
  & \overline{\mathrm{Y}}_\bullet\\
\mathrm{X} &
\end{tikzcd}
\tag*{\textbf{(8.2.1.1)}}
\]
in which \(\mathrm{Y}_\bullet\) is the complement of a divisor with normal crossings in a proper and smooth simplicial scheme \(\overline{\mathrm{Y}}_\bullet\), and is a proper hypercovering of \(\mathrm{X}\). We have (5.3.5 (II))
\[
\mathrm{H}^\bullet(\mathrm{X})
\xrightarrow{\sim}\mathrm{H}^\bullet(\mathrm{Y}_\bullet).
\tag*{\textbf{(8.2.1.2)}}
\]

In (8.1.19), we endowed \(\mathrm{H}^\bullet(\mathrm{Y}_\bullet)\) with a mixed Hodge structure. Let \(\mathcal{H}(\alpha)\) be the mixed Hodge structure on \(\mathrm{H}^\bullet(\mathrm{X})\) obtained from it by (8.2.1.2).

\textit{Proposition} \textbf{(8.2.2)}. — \textit{With the preceding notation, the mixed Hodge structure \(\mathcal{H}(\alpha)\) on \(\mathrm{H}^\bullet(\mathrm{X})\) is independent of the choice of \(\alpha\). It is called the} mixed Hodge structure on \(\mathrm{H}^\bullet(\mathrm{X},\mathbf{Z})\). \textit{For every morphism of schemes \(f:\mathrm{X}_1\to\mathrm{X}_2\), \(f^\bullet:\mathrm{H}^\bullet(\mathrm{X}_2)\to\mathrm{H}^\bullet(\mathrm{X}_1)\) is a morphism of mixed Hodge structures.}

The proof uses (6.2.8), and is parallel to that of (3.2.11 C)) (cf. also (8.3.3)).

\textit{Remark} \textbf{(8.2.3)}. — For \(\mathrm{X}\) smooth, one checks by taking \(\overline{\mathrm{Y}}_\bullet\) to be a constant simplicial scheme that the mixed Hodge structure (8.2.2) coincides with the one defined in (3.2.12).

\textit{Theorem} \textbf{(8.2.4)}. — \textit{The pairs \((p,q)\) for which the Hodge number \(h^{pq}\) of \(\mathrm{H}^n(\mathrm{X})\) is nonzero satisfy the following conditions}

\noindent\hspace*{\parindent}\makebox[2em][l]{(i)}
\((p,q)\in[0,n]\times[0,n]\).

\noindent\hspace*{\parindent}\makebox[2em][l]{(ii)}
\textit{If \(\dim\mathrm{X}=\mathrm{N}\), and \(n\geq\mathrm{N}\), then}
\[
(p,q)\in[n-\mathrm{N},\mathrm{N}]\times[n-\mathrm{N},\mathrm{N}].
\]

\noindent\hspace*{\parindent}\makebox[2em][l]{(iii)}
\textit{If \(\mathrm{X}\) is proper, then \(p+q\leq n\).}

\noindent\hspace*{\parindent}\makebox[2em][l]{(iv)}
\textit{If \(\mathrm{X}\) is smooth, then \(p+q\geq n\). The same conclusion holds for \(\mathrm{X}\) an equidimensional “rational homology manifold” of dimension \(\mathrm{N}\), \textup{i.e.} if for every \(x\in\mathrm{X}\), \(\mathrm{H}^i_{\{x\}}(\mathrm{X},\mathbf{Q})=0\) if \(i\neq2\mathrm{N}\), \(=\mathbf{Q}\) if \(i=2\mathrm{N}\).}

\Page{39}
\[
\begin{tikzcd}[
  baseline=(current bounding box.south),
  row sep=7mm,column sep=7mm,
  cells={nodes={inner sep=0pt,outer sep=0pt,text height=0pt,text depth=0pt}},
  arrows={start anchor=center,end anchor=center}
]
{} & {} & {} & {} & {}\\
{} \arrow[rr,no head,"\text{smooth}"'{pos=.7,yshift=-3mm}]
   \arrow[ddrr,no head]
 & {}
 & {} \arrow[dd,no head,"n"{pos=1,anchor=north,yshift=-2mm}]
 & {} & {}\\
{} & {} & {} & {} & {}\\
{} \arrow[uuu,"q"{pos=1,anchor=east,xshift=-2mm,yshift=-2mm}]
   \arrow[rrrr,"p"{pos=1,anchor=west,xshift=2mm}]
   \arrow[rrrr,phantom,"\text{proper}"{pos=.18,anchor=south,yshift=1mm}]
   \arrow[rrrr,phantom,"n\leq\mathrm{N}"'{pos=.22,anchor=north,yshift=-7mm}]
 & {} & {} & {} & {}
\end{tikzcd}
\qquad\qquad
\begin{tikzcd}[
  baseline=(current bounding box.south),
  row sep=5.5mm,column sep=5.5mm,
  cells={nodes={inner sep=0pt,outer sep=0pt,text height=0pt,text depth=0pt}},
  arrows={start anchor=center,end anchor=center}
]
{} & {} & {} & {} & {} & {} & {}\\
{} \arrow[dr,dashed,no head]
 & {} & {}
 & {} \arrow[d,dashed,no head]
 & {} & {} & {}\\
{} \arrow[r,dashed,no head,"\mathrm{N}"{pos=0,anchor=east,xshift=-2mm}]
 & {} \arrow[rr,no head,"\text{smooth}"'{pos=.7,yshift=-2mm}]
      \arrow[dd,no head]
      \arrow[ddrr,no head]
 & {}
 & {} \arrow[r,dashed,no head] \arrow[dd,no head]
 & {} & {} & {}\\
{} & {} & {} & {} & {} & {} & {}\\
{} & {} \arrow[rr,no head,"\text{proper}"{pos=.43,anchor=south,yshift=.7mm}]
 & {}
 & {} \arrow[dr,dashed,no head,"n"{pos=1,anchor=north,yshift=-2mm}]
      \arrow[d,dashed,no head,"\mathrm{N}"{pos=1,anchor=north,yshift=-2mm}]
 & {} & {} & {}\\
{} \arrow[uuuuu,"q"{pos=1,anchor=east,xshift=-2mm,yshift=-2mm}]
   \arrow[rrrrrr,"p"{pos=1,anchor=west,xshift=2mm}]
   \arrow[rrrrrr,phantom,"n\geq\mathrm{N}"'{pos=.61,anchor=north,yshift=-7mm}]
 & {} & {} & {} & {} & {} & {}
\end{tikzcd}
\]

Assertion (ii) will be proved in (8.3.10). Assertion (i) follows from (8.1.20 (iii)). For \(\mathrm{X}\) proper, one can find a diagram (8.2.1.1) in which \(\mathrm{Y}_\bullet=\overline{\mathrm{Y}}_\bullet\) and (iii) follows from (8.1.20 (iv)). For \(\mathrm{X}\) smooth, (iv) follows from (8.2.3) and (3.2.15). The general case follows: if \(p:\widetilde{\mathrm{X}}\to\mathrm{X}\) is a resolution of the singularities of \(\mathrm{X}\), \(p^*:\mathrm{H}^\bullet(\mathrm{X},\mathbf{Q})\to\mathrm{H}^\bullet(\widetilde{\mathrm{X}},\mathbf{Q})\) is injective, since the transpose under Poincaré duality \(p_!\) of \(p^*:\mathrm{H}_c^\bullet(\mathrm{X},\mathbf{Q})\to\mathrm{H}_c^\bullet(\widetilde{\mathrm{X}},\mathbf{Q})\) is a retraction of it.

\textit{Proposition} \textbf{(8.2.5)}. — \textit{Suppose \(\mathrm{X}\) is proper. If \(u:\mathrm{Y}\to\mathrm{X}\) is a proper surjective morphism, with \(\mathrm{Y}\) smooth, then the quotient of weight \(n\) of \(\mathrm{H}^n(\mathrm{X},\mathbf{Q})\) is the image of \(\mathrm{H}^n(\mathrm{X},\mathbf{Q})\) in \(\mathrm{H}^n(\mathrm{Y},\mathbf{Q})\). The sequence}
\[
\mathrm{H}^n(\mathrm{X},\mathbf{Q})
\longrightarrow
\mathrm{H}^n(\mathrm{Y},\mathbf{Q})
\xrightarrow[\operatorname{pr}_1^*-\operatorname{pr}_2^*]{}
\mathrm{H}^n(\mathrm{Y}\times_{\mathrm{X}}\mathrm{Y},\mathbf{Q})
\tag*{\textbf{(8.2.5.1)}}
\]
\textit{is exact.}

Let \(\mathrm{Y}_\bullet\) be as in (8.2.1.1), with \(\mathrm{Y}_0=\mathrm{Y}\) and \(\mathrm{Y}_\bullet=\overline{\mathrm{Y}}_\bullet\). By ((8.1.20 (ii))), we have an exact sequence
\[
\begin{tikzcd}[column sep=small]
0 \arrow[r]
 & \mathrm{H}^n(\mathrm{X},\mathbf{Q})/
   \mathrm{W}^{n-1}\mathrm{H}^n(\mathrm{X},\mathbf{Q}) \arrow[r]
 & \mathrm{H}^n(\mathrm{Y}_0,\mathbf{Q})
   \arrow[r,shift left=.4ex,"\delta_0"]
   \arrow[r,shift right=.4ex,"\delta_1"']
 & \mathrm{H}^n(\mathrm{Y}_1,\mathbf{Q})
\end{tikzcd}
\tag*{\textbf{(8.2.5.2)}}
\]
which implies (8.2.5) (conversely, this exact sequence follows easily from (8.2.5) and (8.2.7) below).

\textit{Proposition} \textbf{(8.2.6)}. — \textit{Consider morphisms of schemes}
\[
\begin{tikzcd}[column sep=small]
\mathrm{Y} \arrow[r,"f"]
 & \mathrm{X} \arrow[r,hook,"j"]
 & \overline{\mathrm{X}}.
\end{tikzcd}
\]
\textit{Assume that \(\mathrm{Y}\) is proper, that \(\mathrm{X}\) is smooth and that \(\overline{\mathrm{X}}\) is a smooth compactification of \(\mathrm{X}\). Then, \(\mathrm{H}^\bullet(\widetilde{\mathrm{X}},\mathbf{Q})\) and \(\mathrm{H}^\bullet(\mathrm{X},\mathbf{Q})\) have the same image in \(\mathrm{H}^\bullet(\mathrm{Y},\mathbf{Q})\).}

The proof is parallel to that of (3.2.18). It suffices to prove that, for all \(i\) and \(n\) \(\operatorname{Gr}_i^{\mathrm{W}}(f^*)\) and \(\operatorname{Gr}_i^{\mathrm{W}}((jf)^*)\) have the same image in \(\operatorname{Gr}_i^{\mathrm{W}}\mathrm{H}^n(\mathrm{Y},\mathbf{Q})\).

\textit{a)} For \(i>n\), these images are zero since \(\operatorname{Gr}_i^{\mathrm{W}}\mathrm{H}^n(\mathrm{Y},\mathbf{Q})=0\) (8.2.4 (iii)).

\textit{b)} For \(i<n\), these images are zero since \(\operatorname{Gr}_i^{\mathrm{W}}\mathrm{H}^n(\mathrm{X},\mathbf{Q})=0\) (8.2.4 (iv)).

\textit{c)} For \(i=n\), these images are equal since \(\operatorname{Gr}_n^{\mathrm{W}}(j^*)\) is surjective (3.2.17).

\Page{40}
\textit{Proposition} \textbf{(8.2.7)}. — \textit{Consider morphisms of schemes}
\[
\begin{tikzcd}[column sep=small]
\widetilde{\mathrm{X}} \arrow[r,"\pi"]
 & \mathrm{X} \arrow[r,"f"]
 & \mathrm{Y}.
\end{tikzcd}
\]
\textit{Assume that \(\mathrm{Y}\) is smooth, that \(\mathrm{X}\) is proper, that \(\widetilde{\mathrm{X}}\) is proper and smooth and that \(\pi\) is surjective. Then, the kernels of \(f^*\) and of \((f\pi)^*\) in \(\mathrm{H}^\bullet(\mathrm{Y},\mathbf{Q})\) are equal.}

Let us prove that \(\operatorname{Ker}\operatorname{Gr}_i^{\mathrm{W}}(f^*)=\operatorname{Ker}(\operatorname{Gr}_i^{\mathrm{W}}((f\pi)^*))\) in \(\operatorname{Gr}_i^{\mathrm{W}}(\mathrm{H}^n(\mathrm{Y},\mathbf{Q}))\).

\textit{a)} For \(i<n\), these kernels are zero since \(\operatorname{Gr}_i^{\mathrm{W}}\mathrm{H}^n(\mathrm{Y},\mathbf{Q})=0\) ((8.2.4 (iv))).

\textit{b)} For \(i>n\), these kernels are the whole of \(\operatorname{Gr}_i^{\mathrm{W}}(\mathrm{H}^n(\mathrm{Y},\mathbf{Q}))\) since \(\operatorname{Gr}_i^{\mathrm{W}}\mathrm{H}^n(\mathrm{X},\mathbf{Q})=0\) ((8.2.4 (iii))).

\textit{c)} For \(i=n\), these kernels are equal since \(\operatorname{Gr}_n^{\mathrm{W}}(\pi^*)\) is injective (8.2.5.2).

\textit{Corollary} \textbf{(8.2.8)}. — \textit{Consider morphisms of schemes}
\[
\begin{tikzcd}[column sep=small]
\widetilde{\mathrm{X}} \arrow[r,"\pi"]
 & \mathrm{X} \arrow[r,hook,"f"]
 & \mathrm{Y}.
\end{tikzcd}
\]
\textit{Assume that \(\mathrm{Y}\) is proper and smooth, that \(\mathrm{X}\) is a closed subscheme of \(\mathrm{Y}\) and that \(\widetilde{\mathrm{X}}\) is a resolution of the singularities of \(\mathrm{X}\). Let \(q=f\pi\) and \(\mathrm{U}=\mathrm{Y}-\mathrm{X}\). Then, the sequence}
\[
\mathrm{H}^\bullet(\widetilde{\mathrm{X}},\mathbf{Q})
\xrightarrow{q_!}\mathrm{H}^\bullet(\mathrm{Y},\mathbf{Q})
\longrightarrow\mathrm{H}^\bullet(\mathrm{U},\mathbf{Q})
\tag*{\textbf{(8.2.8.1)}}
\]
\textit{is exact.}

By (8.2.7), we have
\[
\operatorname{Ker}(f^*:\mathrm{H}^\bullet(\mathrm{Y},\mathbf{Q})
 \to\mathrm{H}^\bullet(\mathrm{X},\mathbf{Q}))
=
\operatorname{Ker}(q^*:\mathrm{H}^\bullet(\mathrm{Y},\mathbf{Q})
 \to\mathrm{H}^\bullet(\widetilde{\mathrm{X}},\mathbf{Q})).
\]
The long exact cohomology sequence of the pair \((\mathrm{Y},\mathrm{X})\) therefore gives an exact sequence
\[
\mathrm{H}_c^\bullet(\mathrm{U},\mathbf{Q})
\longrightarrow\mathrm{H}^\bullet(\mathrm{Y},\mathbf{Q})
\longrightarrow\mathrm{H}^\bullet(\widetilde{\mathrm{X}},\mathbf{Q}),
\tag*{\textbf{(8.2.8.2)}}
\]
and (8.2.8.1) is the transpose of (8.2.8.2) under Poincaré duality.

\textit{Remark} \textbf{(8.2.9)}. — Corollary (8.2.8) is used under more general hypotheses in ([8], (9.7), p. 161). As Lieberman pointed out, the justificati n given in \textit{loc. cit.} is insufficient.

We deduce from (8.1.25) that

\textit{Proposition} \textbf{(8.2.10)}. — \textit{The Künneth isomorphisms}
\[
\mathrm{H}^\bullet(\mathrm{X}\times\mathrm{Y},\mathbf{Q})
\simeq
\mathrm{H}^\bullet(\mathrm{X},\mathbf{Q})
\otimes\mathrm{H}^\bullet(\mathrm{Y},\mathbf{Q})
\]
\textit{are isomorphisms of mixed Hodge structures.}

\textit{Corollary} \textbf{(8.2.11)}. — \textit{The cup product}
\[
\mathrm{H}^\bullet(\mathrm{X})\otimes\mathrm{H}^\bullet(\mathrm{X})
\longrightarrow\mathrm{H}^\bullet(\mathrm{X})
\]
\textit{is a morphism of mixed Hodge structures.}

This follows from (8.2.10) and from (8.2.2) applied to the diagonal map \(\Delta:\mathrm{X}\to\mathrm{X}\times\mathrm{X}\).

\Page{41}
\section*{8.3. Hodge theory of simplicial schemes.}

\textbf{(8.3.1)} For \(\mathrm{X}^\bullet\) a separated simplicial scheme (or algebraic space), we shall have to consider diagrams of simplicial schemes
\[
\begin{tikzcd}[column sep=large,row sep=large]
\mathrm{Z}_\bullet \arrow[r,hook,"j"] \arrow[d,"p"']
 & \overline{\mathrm{Z}}_\bullet\\
\mathrm{X}_\bullet &
\end{tikzcd}
\tag*{\textbf{(8.3.1.1)}}
\]
satisfying:

\textit{a)} \(j\) is an open immersion with dense image; \(\overline{\mathrm{Z}}_\bullet\) is proper.

\textit{b)} The obvious morphisms \(\mathrm{Z}_n\to(\operatorname{cosq}_{n-1}^{\mathrm{X}_\bullet}\operatorname{sq}_{n-1}\mathrm{Z}_\bullet)_n\) are proper and surjective.

For \(\overline{\mathrm{Z}}_\bullet\) reduced, one can regard \(j\) as a simplicial object of the category \(\mathcal{C}\) of pairs \((\mathrm{Y},\overline{\mathrm{Y}})\) consisting of a proper reduced scheme \(\overline{\mathrm{Y}}\) and a dense (Zariski) open subset \(\mathrm{Y}\) of \(\overline{\mathrm{Y}}\). The category \(\mathcal{C}\) satisfies (6.2.4.1), which allows us to apply the analogue of (6.2.4) to it.

\textit{Proposition} \textbf{(8.3.2)}. — (i) \textit{For every separated simplicial scheme \(\mathrm{X}_\bullet\), there exists a diagram \textup{(8.3.1.1)} satisfying a), b) and such that \(\overline{\mathrm{Z}}_\bullet\) is proper and smooth and \(\mathrm{Z}_\bullet\) is the complement of a divisor with normal crossings.}

(ii) \textit{Let \(\mathrm{F}:\mathrm{X}_{\bullet1}\to\mathrm{X}_{\bullet2}\) be a morphism of separated simplicial schemes. Let \((\mathrm{X}_{\bullet i},\mathrm{Z}_{\bullet i},\overline{\mathrm{Z}}_{\bullet i})_{i=1,2}\) be two diagrams \textup{(8.3.1.1)} satisfying a) and b). There then exist \((\mathrm{X}_{\bullet1},\mathrm{Z}_\bullet,\overline{\mathrm{Z}}_\bullet)\) as in \textup{(i)} and a commutative diagram}
\[
\begin{tikzcd}[column sep=large,row sep=large]
\mathrm{Z}_\bullet
  \arrow[r,hook] \arrow[d] \arrow[drr]
 & \overline{\mathrm{Z}}_\bullet
   \arrow[d,crossing over] \arrow[drr]
 & &\\
\mathrm{Z}_{\bullet1}
  \arrow[r,hook,"j_1"'] \arrow[d,"p_1"']
 & \overline{\mathrm{Z}}_{\bullet1}
 & \mathrm{Z}_{\bullet2}
   \arrow[r,hook,"j_2"'] \arrow[d,"p_2"']
 & \overline{\mathrm{Z}}_{\bullet2}\\
\mathrm{X}_{\bullet1} \arrow[rr,"f"]
 & & \mathrm{X}_{\bullet2} &
\end{tikzcd}
\]

The proof, analogous to (6.2.5) but more complicated, is left to the reader.

\textbf{(8.3.3)} Consider a diagram (8.3.1.1) as in (8.3.2 (i)). By (5.3.5 (V)), the map \(\mathrm{H}^\bullet(\mathrm{X}_\bullet,\mathbf{Z})\to\mathrm{H}^\bullet(\mathrm{Z}_\bullet,\mathbf{Z})\) is an isomorphism. In (8.1.19) and (8.1.20), we endowed \(\mathrm{H}^\bullet(\mathrm{Z}_\bullet,\mathbf{Z})\) with a mixed Hodge structure, hence a mixed
\Page{42}
Hodge structure on \(\mathrm{H}^\bullet(\mathrm{X}_\bullet,\mathbf{Z})\). It is independent of the choice of the diagram (8.3.1.1): by (8.3.2 (ii)) for \(f=\mathrm{Id}\), two systems \((\mathrm{Z}_{\bullet i},\overline{\mathrm{Z}}_{\bullet i},p_i)_{i=1,2}\) are always dominated by a third \((\mathrm{Z}_\bullet,\overline{\mathrm{Z}}_\bullet,p)\), whence a commutative diagram of isomorphisms
\[
\begin{tikzcd}[column sep=large,row sep=huge]
 & \mathrm{H}^\bullet(\mathrm{X}_\bullet,\mathbf{Z})
   \arrow[dl] \arrow[d] \arrow[dr] &\\
\mathrm{H}^\bullet(\mathrm{Z}_{\bullet1},\mathbf{Z}) \arrow[r]
 & \mathrm{H}^\bullet(\mathrm{Z}_\bullet,\mathbf{Z})
 & \mathrm{H}^\bullet(\mathrm{Z}_{\bullet2},\mathbf{Z}) \arrow[l]
\end{tikzcd}
\]

\textit{Definition} \textbf{(8.3.4)}. — \textit{The} mixed Hodge structure \textit{on \(\mathrm{H}^\bullet(\mathrm{X}_\bullet,\mathbf{Z})\) is the one defined above.}

It follows from (8.3.2) (ii) that this mixed Hodge structure is functorial in \(\mathrm{X}_\bullet\).

\textit{Proposition} \textbf{(8.3.5)}. — \textit{Let \(\mathrm{X}_\bullet\) be a separated simplicial scheme. The spectral sequence}
\[
\mathrm{E}_1^{pq}=\mathrm{H}^q(\mathrm{X}_p,\mathbf{Z})
\Rightarrow\mathrm{H}^{p+q}(\mathrm{X}_\bullet,\mathbf{Z})
\]
\textit{is a spectral sequence of mixed Hodge structures.}

For a diagram as in (8.3.1), the maps \(p_n:\mathrm{Z}_n\to\mathrm{X}_n\) are proper and surjective. We can then apply (6.2.4) to the category of simplicial objects of the category \(\mathcal{C}\) (8.3.1) to obtain as in (6.2.5) the following lemma:

\textit{Lemma} \textbf{(8.3.6)}. — \textit{Let \(\mathrm{X}_\bullet\) be a separated simplicial scheme. There exist a bisimplicial scheme \(\mathrm{Z}_{\bullet\bullet}\) augmented over \(\mathrm{X}_\bullet\):}
\[
\varepsilon_{n,m}:\mathrm{Z}_{n,m}\to\mathrm{X}_n
\]
\textit{and \(i:\mathrm{Z}_{\bullet\bullet}\hookrightarrow\overline{\mathrm{Z}}_{\bullet\bullet}\) such that:}

\textit{a) \(\overline{\mathrm{Z}}_{\bullet\bullet}\) is a proper and smooth bisimplicial scheme, and \(\mathrm{Z}_{\bullet\bullet}\) is the complement of a divisor with normal crossings in \(\overline{\mathrm{Z}}_{\bullet\bullet}\);}

\textit{b) for every \(n\), the augmented simplicial scheme \(\mathrm{Z}_{n\bullet}\to\mathrm{X}_n\) is a proper hypercovering of \(\mathrm{X}_n\).}

In (8.1.2.1), we endowed \(\mathrm{H}^\bullet(\mathrm{Z}_{\bullet\bullet},\mathbf{Z})\) with a mixed Hodge structure. If \(\delta\mathrm{Z}_{\bullet\bullet}\) is the diagonal simplicial scheme of \(\mathrm{Z}_{\bullet\bullet}\), we have a commutative diagram of morphisms
\[
\begin{tikzcd}[column sep=large,row sep=huge]
\mathrm{H}^\bullet(\mathrm{X}_\bullet,\mathbf{Z})
  \arrow[r,"\gamma"] \arrow[dr,"\alpha"]
 & \mathrm{H}^\bullet(\mathrm{Z}_{\bullet\bullet},\mathbf{Z})
   \arrow[d,equal,"\beta"]\\
 & \mathrm{H}^\bullet(\delta\mathrm{Z}_{\bullet\bullet},\mathbf{Z})
\end{tikzcd}
\]
\Page{43}
in which \(\beta\) and \(\gamma\) are isomorphisms, and \(\alpha\) and \(\beta\) are morphisms of mixed Hodge structures; \(\gamma\) is therefore an isomorphism of mixed Hodge structures.

We have a morphism of spectral sequences
\[
\begin{tikzcd}[column sep=large,row sep=huge]
\mathrm{H}^q(\mathrm{X}_p,\mathbf{Z})
  \arrow[r,Rightarrow] \arrow[d]
 & \mathrm{H}^{p+q}(\mathrm{X}_\bullet,\mathbf{Z}) \arrow[d]
 & \text{\textbf{(5.2.3)}}\\
\mathrm{H}^q(\mathrm{Z}_{p\bullet},\mathbf{Z})
  \arrow[r,Rightarrow]
 & \mathrm{H}^{p+q}(\mathrm{Z}_{\bullet\bullet},\mathbf{Z})
 & \text{\textbf{(8.1.23.1)}}
\end{tikzcd}
\]
this morphism is an isomorphism compatible with the mixed Hodge structures, so that (8.3.5) follows from (8.1.23).

\textit{Example} \textbf{(8.3.7)}. — Let \(\mathrm{X}_\bullet\) be a simplicial scheme not necessarily separated. There then exist a separated simplicial scheme \(\mathrm{Y}_\bullet\) and \(u:\mathrm{Y}_\bullet\to\mathrm{X}_\bullet\) such that for \(n\geq0\), the map \(u'_n:\mathrm{Y}_n\to(\operatorname{cosq}_{n-1}^{\mathrm{X}_\bullet}\operatorname{sq}_{n-1}\mathrm{Y}_\bullet)_n\) is étale and surjective (for example the disjoint union of the open subsets of an open covering).

The map \(u^\bullet:\mathrm{H}^\bullet(\mathrm{X}_\bullet,\mathbf{Z})\to\mathrm{H}^\bullet(\mathrm{Y}_\bullet,\mathbf{Z})\) is an isomorphism ((5.3.5) (V) and (III)). By definition, the \textit{mixed Hodge structure} on \(\mathrm{H}^\bullet(\mathrm{X}_\bullet,\mathbf{Z})\) is obtained from that on \(\mathrm{H}^\bullet(\mathrm{Y}_\bullet,\mathbf{Z})\) \textit{via} \(u^\bullet\). It does not depend on the choice of \(\mathrm{Y}_\bullet\).

For \(\mathrm{X}\) a nonseparated scheme, the \textit{mixed Hodge structure} on \(\mathrm{H}^\bullet(\mathrm{X},\mathbf{Z})\) is that on \(\mathrm{H}^\bullet(\mathrm{X}_\bullet,\mathbf{Z})\), for \(\mathrm{X}_\bullet\) the constant simplicial scheme defined by \(\mathrm{X}\). It is plausible, but not proved, that (8.2.4) (iv) remains valid for \(\mathrm{X}\) smooth and nonseparated.

\textit{Example} \textbf{(8.3.8)}. — Let \(u:\mathrm{X}_\bullet\to\mathrm{Y}_\bullet\) be a morphism of simplicial schemes. The relative cohomology groups \(\mathrm{H}^n(\mathrm{Y}_\bullet\operatorname{mod}\mathrm{X}_\bullet,\mathbf{Z})\) are by definition (6.3.2) the cohomology groups of the simplicial scheme \(\mathrm{C}(u)\). They are therefore endowed with a mixed Hodge structure.

\textit{Proposition} \textbf{(8.3.9)}. — \textit{The long exact sequence of relative cohomology}
\[
\mathrm{H}^n(\mathrm{Y}_\bullet,\mathbf{Z})
\to\mathrm{H}^n(\mathrm{X}_\bullet,\mathbf{Z})
\xrightarrow{j}
\mathrm{H}^{n+1}(\mathrm{Y}_\bullet\operatorname{mod}\mathrm{X}_\bullet,\mathbf{Z})
\to\mathrm{H}^{n+1}(\mathrm{Y},\mathbf{Z})
\]
\textit{is a long exact sequence of mixed Hodge structures.}

We reduce to the case where \(\mathrm{X}_\bullet\) and \(\mathrm{Y}_\bullet\) are smooth, and where, moreover, there is a commutative diagram
\[
\begin{tikzcd}[column sep=large,row sep=large]
\mathrm{X}_\bullet \arrow[r,"u"] \arrow[d,hook]
 & \mathrm{Y}_\bullet \arrow[d,hook]\\
\overline{\mathrm{X}}_\bullet \arrow[r,"\bar{u}"]
 & \overline{\mathrm{Y}}_\bullet
\end{tikzcd}
\]
with \(\overline{\mathrm{X}}_\bullet\) and \(\overline{\mathrm{Y}}_\bullet\) proper and smooth, and \(\mathrm{X}_\bullet\) and \(\mathrm{Y}_\bullet\) complements of divisors with normal crossings \(\mathrm{D}_\bullet\) and \(\mathrm{E}_\bullet\) in \(\overline{\mathrm{X}}_\bullet\) and \(\overline{\mathrm{Y}}_\bullet\).

\Page{44}
There then exist bifiltered resolutions (with acyclic components in the associated graded) \(a:\Omega_{\overline{\mathrm{X}}_\bullet}^\bullet(\log\mathrm{D}_\bullet)\to\mathrm{K}\) and \(b:\Omega_{\overline{\mathrm{Y}}_\bullet}^\bullet(\log\mathrm{E}_\bullet)\to\mathrm{L}\) which fit into a commutative diagram of bifiltered complexes
\[
\begin{tikzcd}[column sep=large,row sep=huge]
\Omega_{\overline{\mathrm{X}}_\bullet}^\bullet(\log\mathrm{D}_\bullet)
  \arrow[d,"a"]
 & \Omega_{\overline{\mathrm{Y}}_\bullet}^\bullet(\log\mathrm{E}_\bullet)
   \arrow[l,"\bar{u}^*"'] \arrow[d,"b"]\\
\mathrm{K}
 & \mathrm{L} \arrow[l,"v"']
\end{tikzcd}
\]
(\(\bar{u}^*\) and \(v\) are \(\bar{u}\)-morphisms). The long exact sequence (8.3.9) is then identified with the exact sequence obtained from the distinguished triangle
\[
\begin{tikzcd}[column sep=small,row sep=large]
 & \mathrm{C}(\Gamma v) \arrow[dl] &\\
\Gamma\mathrm{L} \arrow[rr,"v"]
 & & \Gamma\mathrm{K} \arrow[ul]
\end{tikzcd}
\]
The complex \(\mathrm{C}(v)\) on \(\mathrm{C}(\bar{u})\) is a bifiltered resolution of \(\mathrm{C}(u^*:\Omega_{\overline{\mathrm{Y}}_\bullet}^\bullet(\log\mathrm{E}_\bullet)\to\Omega_{\overline{\mathrm{X}}_\bullet}^\bullet(\log\mathrm{D}_\bullet))\). It follows that the Hodge and weight filtrations on \(\mathrm{H}^\bullet(\mathrm{Y}_\bullet\operatorname{mod}\mathrm{X}_\bullet,\mathbf{C})\) are obtained from the filtrations on the cone \(\mathrm{C}(\Gamma v)\), and (8.3.9) follows formally.

\textbf{(8.3.10)} \textit{Proof of} (8.2.4) (ii).

The proof proceeds by induction on \(\mathrm{N}=\dim(\mathrm{X})\). We may assume \(\mathrm{X}\) reduced. Let \(\mathrm{U}\) be a smooth dense Zariski open subset of \(\mathrm{X}\) and \(\mathrm{Y}=\mathrm{X}-\mathrm{U}\). By resolution of singularities, there exists a proper birational morphism \(p:\mathrm{X}'\to\mathrm{X}\) with smooth source which induces an isomorphism of \(p^{-1}(\mathrm{U})\) with \(\mathrm{U}\). Let \(\mathrm{Y}'=p^{-1}(\mathrm{Y})\). For \(n\geq\mathrm{N}\), consider the diagram of mixed Hodge structures
\[
\begin{tikzcd}[column sep=small,row sep=huge]
 & \mathrm{H}^n(\mathrm{X}\operatorname{mod}\mathrm{Y},\mathbf{Z})
   \arrow[r] \arrow[d,"p^*"]
 & \mathrm{H}^n(\mathrm{X},\mathbf{Z}) \arrow[r] \arrow[d]
 & \mathrm{H}^n(\mathrm{Y},\mathbf{Z}) \arrow[d]\\
\mathrm{H}^{n-1}(\mathrm{Y}',\mathbf{Z}) \arrow[r]
 & \mathrm{H}^n(\mathrm{X}'\operatorname{mod}\mathrm{Y}',\mathbf{Z})
   \arrow[r]
 & \mathrm{H}^n(\mathrm{X}',\mathbf{Z}) \arrow[r]
 & \mathrm{H}^n(\mathrm{Y}',\mathbf{Z})
\end{tikzcd}
\]

Let \(\Phi\) be the family of supports on \(\mathrm{U}\) consisting of those \(\mathrm{A}\subset\mathrm{U}\) such that, in \(\mathrm{X}\) (or \(\mathrm{X}'\), which amounts to the same thing), \(\overline{\mathrm{A}}\subset\mathrm{U}\). We have
\[
\mathrm{H}^n(\mathrm{X}\operatorname{mod}\mathrm{Y},\mathbf{Z})
=\mathrm{H}^n_\Phi(\mathrm{U},\mathbf{Z})
=\mathrm{H}^n(\mathrm{X}'\operatorname{mod}\mathrm{Y}',\mathbf{Z})
\]
so that the arrow marked \(p^*\) is an isomorphism.

For \(h^{pq}\) a nonzero Hodge number of one of the mixed Hodge structures involved, we have:

% ---- 045-052 ----
\Page{45}
\noindent\textit{a)} for \(\mathrm{H}^{n}(\mathrm{X}',\mathbf{Z})\) : \((p,q)\in[n-\mathrm{N},\mathrm{N}]\times[n-\mathrm{N},\mathrm{N}]\)

\noindent\textit{b)} for \(\mathrm{H}^{n-1}(\mathrm{Y}',\mathbf{Z})\) : \((p,q)\in[n-\mathrm{N},\mathrm{N}-1]\times[n-\mathrm{N},\mathrm{N}-1]\)

\noindent\textit{c)} for \(\mathrm{H}^{n}(\mathrm{X}'\text{ mod }\mathrm{Y}',\mathbf{Z})=\mathrm{H}^{n}(\mathrm{X}\text{ mod }\mathrm{Y},\mathbf{Z})\) : \((p,q)\in[n-\mathrm{N},\mathrm{N}]\times[n-\mathrm{N},\mathrm{N}]\)

\noindent\textit{d)} for \(\mathrm{H}^{n}(\mathrm{Y},\mathbf{Z})\) : \((p,q)\in[n-\mathrm{N}+1,\mathrm{N}-1]\times[n-\mathrm{N}+1,\mathrm{N}-1]\)

\noindent\textit{e)} for \(\mathrm{H}^{n}(\mathrm{X},\mathbf{Z})\) : \((p,q)\in[n-\mathrm{N},\mathrm{N}]\times[n-\mathrm{N},\mathrm{N}]\).

\textit{a)} follows from the fact that \(\mathrm{X}'\) is smooth, and that the assertion holds for the terms \({}_{\mathrm{W}}\mathrm{E}_{1}\) of the spectral sequence (3.2.4.1). \textit{b)} and \textit{d)} follow from the induction hypothesis : we have \(\dim\mathrm{Y},\dim\mathrm{Y}'\leq\mathrm{N}-1\). \textit{c)} follows from \textit{a)} and \textit{b)}, and \textit{e)} follows from \textit{c)} and \textit{d)}. This completes the proof.

\section*{9. EXAMPLES AND APPLICATIONS}

\section*{9.1. Cohomology of groups and classifying spaces.}

Let \(\mathrm{G}\) be a connected complex linear algebraic group. We propose to compute the mixed Hodge structure on \(\mathrm{H}^{\bullet}(\mathrm{G})\). We shall begin by computing that on \(\mathrm{H}^{\bullet}(\mathrm{B}_{\bullet\mathrm{G}})\), by reduction to the case of tori; we shall then compute that on \(\mathrm{H}^{\bullet}(\mathrm{G})\) \textit{via} the Eilenberg-Moore spectral sequence (6.1.5). As a bonus, we obtain the strange theorem (9.1.7).

\textit{Theorem} \textbf{(9.1.1)}. — \textit{Let \(\mathrm{G}\) be a linear algebraic group. We have \(\mathrm{H}^{2n+1}(\mathrm{B}_{\bullet\mathrm{G}},\mathbf{Q})=0\), and \(\mathrm{H}^{2n}(\mathrm{B}_{\bullet\mathrm{G}},\mathbf{Q})\otimes\mathbf{C}\) is purely of type \((n,n)\) \textup{(III.0.5)}.}

Let \(\mathrm{G}^{0}\) be the identity component of \(\mathrm{G}\), and let \(\mathrm{T}\subset\mathrm{G}^{0}\) be a maximal torus. Since \(\mathrm{H}^{\bullet}(\mathrm{B}_{\bullet\mathrm{G}},\mathbf{Q})=\mathrm{H}^{\bullet}(\mathrm{B}_{\bullet\mathrm{G}^{0}},\mathbf{Q})^{\mathrm{G}/\mathrm{G}^{0}}\), we have, by the \textit{splitting principle} (see (6.1.6)),
\[
\mathrm{H}^{\bullet}(\mathrm{B}_{\bullet\mathrm{G}},\mathbf{Q})
\hookrightarrow
\mathrm{H}^{\bullet}(\mathrm{B}_{\bullet\mathrm{T}},\mathbf{Q}).
\]
This inclusion is a morphism of mixed Hodge structures : it suffices to treat the case of a torus, or even, by the Künneth formula, the case where \(\mathrm{G}=\mathbf{G}_{m}\). Here are two proofs of (9.1.1) in this case.

\textit{1\textsuperscript{st} proof.} — The cohomology of \(\mathbf{G}_{m}\) (\(=\)the projective line minus two points) is : \(\mathrm{H}^{0}=\mathbf{Z}\), \(\mathrm{H}^{1}=\mathbf{Z}\). The only Hodge numbers \(h^{pq}\) of \(\mathrm{H}^{1}(\mathbf{G}_{m})\) that could be nonzero are \(h^{01}=h^{10}\) and \(h^{11}\) ((8.2.4) (i) (iv)). Since \(h^{01}+h^{10}+h^{11}=2h^{01}+h^{11}=1\), we have \(h^{01}=0\), and \(\mathrm{H}^{1}(\mathbf{G}_{m})\) is purely of type \((1,1)\) (for another proof, see (10.3.8)). The term \(\mathrm{E}_{2}^{p,q}\) of the spectral sequence (6.1.5), zero for odd \(p+q\), is therefore purely of type \(\left(\dfrac{p+q}{2},\dfrac{p+q}{2}\right)\). Its abutment has the same homogeneity.

\textit{2\textsuperscript{nd} proof.} — We know that \(\mathbf{P}^{r}(\mathbf{C})\), equipped with \(\mathscr{O}(1)\), is an “approximation” to \(\mathrm{B}_{\mathbf{G}_{m}}\) : the homomorphism \([\mathscr{O}(1)]:\mathrm{H}^{i}(\mathrm{B}_{\bullet\mathbf{G}_{m}})\to\mathrm{H}^{i}(\mathbf{P}^{r}(\mathbf{C}))\) is bijective for \(i\leq2r\). The cohomology group \(\mathrm{H}^{2i}(\mathbf{P}^{r}(\mathbf{C}))\) is generated by the cohomology class of an algebraic cycle (namely, a linear subspace), and hence is of type \((i,i)\). The following proposition allows us to deduce that \(\mathrm{H}^{2i}(\mathrm{B}_{\mathbf{G}_{m}})\) is of type \((i,i)\).

\Page{46}
\textit{Proposition} \textbf{(9.1.2)}. — \textit{Let \(\mathrm{G}\) be an algebraic group (over \(\mathbf{C}\)), and let \(\mathrm{P}\) be an (algebraic) principal homogeneous space with group \(\mathrm{G}\) over a scheme \(\mathrm{X}\). The morphism}
\[
[\mathrm{P}]:\mathrm{H}^{\bullet}(\mathrm{B}_{\mathrm{G}})
=\mathrm{H}^{\bullet}(\mathrm{B}_{\bullet\mathrm{G}})
\to\mathrm{H}^{\bullet}(\mathrm{X})
\]
\textit{is a morphism of mixed Hodge structures.}

This follows from the definition (6.1.3.1) of \([\mathrm{P}]\) and from functoriality (8.3.4).

\textit{Corollary} \textbf{(9.1.3)}. — \textit{The characteristic classes, in \(\mathrm{H}^{2n}(\mathrm{S})\), of an (algebraic) principal homogeneous space \(\mathrm{P}\) over \(\mathrm{S}\), with a linear group \(\mathrm{G}\) as structure group, are purely of type \((n,n)\) : they arise from morphisms \(\mathbf{Q}(-n)\to\mathrm{H}^{2n}(\mathrm{S},\mathbf{Q})\).}

In the special case where \(\mathrm{S}\) is quasi-projective and smooth, we know that the characteristic classes of \(\mathrm{P}\) are even algebraic, i.e. in the image of the Chow ring \((\otimes\mathbf{Q})\).

\textbf{(9.1.4)} Let \(\mathrm{G}\) be a linear algebraic group over a finite field \(\mathbf{F}_{q}\), and let \(\overline{\mathrm{G}}/\overline{\mathbf{F}}_{q}\) be the group obtained from it by extension of scalars to an algebraic closure of \(\mathbf{F}_{q}\). A proof parallel to (9.1.1) shows that, in \(\ell\)-adic cohomology, the eigenvalues of the (geometric) Frobenius acting on \(\mathrm{H}^{2n}(\mathrm{B}_{\bullet\overline{\mathrm{G}}},\mathbf{Q}_{\ell})\) are of the form \(q^{n}.\zeta\) (\(\zeta\) a root of unity). This agrees with the philosophy (I.3). Similarly, the statement (9.1.5) below is inspired by the formulas giving the numbers of points of reductive groups over finite fields.

\textit{Theorem} \textbf{(9.1.5)}. — \textit{Let \(\mathrm{G}\) be a connected linear algebraic group. Let \(\mathrm{P}^{\bullet}(\mathrm{G})\subset\mathrm{H}^{\bullet}(\mathrm{G},\mathbf{Q})\) be the primitive part of the Hopf algebra \(\mathrm{H}^{\bullet}(\mathrm{G},\mathbf{Q})\). It is a mixed Hodge substructure. We have \(\mathrm{P}^{2i}(\mathrm{G})=0\), and \(\mathrm{P}^{2i-1}(\mathrm{G})\) is purely of type \((i,i)\). Finally, the isomorphism}
\[
\bigwedge\mathrm{P}^{\bullet}(\mathrm{G})
\to\mathrm{H}^{\bullet}(\mathrm{G},\mathbf{Q})
\tag*{\textbf{(9.1.5.1)}}
\]
\textit{is an isomorphism of mixed Hodge structures.}

In the spectral sequence (6.1.5), \(\mathrm{E}^{0,q}=0\) for \(q>0\). For \(n>0\), the \textit{edge homomorphism}
\[
\mathrm{H}^{2n}(\mathrm{B}_{\mathrm{G}},\mathbf{Q})
\to\mathrm{E}_{1}^{1,2n}
=\mathrm{H}^{2n-1}(\mathrm{G},\mathbf{Q})
\]
is a morphism of mixed Hodge structures (8.3.5); its image \(\mathrm{P}^{2n-1}(\mathrm{G})\) is therefore a mixed Hodge substructure, and is purely of type \((n,n)\) by (9.1.1). The isomorphism (9.1.5.1) is given by the cup-product, and we conclude by (8.2.11).

\textit{Corollary} \textbf{(9.1.6)}. — \textit{Let \(\mathrm{G}\) be a linear algebraic group. For \(n>0\), we have}
\[
\mathrm{W}_{n}(\mathrm{H}^{n}(\mathrm{G},\mathbf{Q}))=0.
\]
We reduce to the case where \(\mathrm{G}\) is connected, and apply (9.1.5).

\textit{Theorem} \textbf{(9.1.7)}. — \textit{Let \(\mathrm{X}\) be a proper smooth algebraic space on which a connected linear algebraic group \(\mathrm{G}\) acts. Let \(\rho:\mathrm{G}\times\mathrm{X}\to\mathrm{X}\). The diagram}
\Page{47}
\[
\begin{tikzcd}[ampersand replacement=\&, column sep=large, row sep=3.5em]
\mathrm{H}^{\bullet}(\mathrm{X},\mathbf{Q})
  \arrow[r, "\rho^{*}"] \arrow[d, equal]
\&
\mathrm{H}^{\bullet}(\mathrm{X},\mathbf{Q})\otimes
\mathrm{H}^{\bullet}(\mathrm{G},\mathbf{Q})
\\
\mathrm{H}^{\bullet}(\mathrm{X},\mathbf{Q})\otimes_{\mathbf{Q}}\mathbf{Q}
  \arrow[r, equal]
\&
\mathrm{H}^{\bullet}(\mathrm{X},\mathbf{Q})\otimes
\mathrm{H}^{0}(\mathrm{G},\mathbf{Q})
  \arrow[u]
\end{tikzcd}
\]
\textit{is commutative.}

Since the composite map
\[
\mathrm{H}^{\bullet}(\mathrm{X},\mathbf{Q})
\xrightarrow{\rho^{*}}
\mathrm{H}^{\bullet}(\mathrm{X},\mathbf{Q})\otimes\mathrm{H}^{\bullet}(\mathrm{G},\mathbf{Q})
\to
\mathrm{H}^{\bullet}(\mathrm{X},\mathbf{Q})\otimes\mathrm{H}^{0}(\mathrm{G},\mathbf{Q})
=\mathrm{H}^{\bullet}(\mathrm{X},\mathbf{Q})
\]
is the identity (the composite \(\mathrm{X}=\{e\}\times\mathrm{X}\to\mathrm{G}\times\mathrm{X}\to\mathrm{X}\) is the identity), it suffices to prove that for \(i>0\), the \(i\)-th Künneth component of \(\rho^{*}\)
\[
\rho_{i}^{*}:\mathrm{H}^{n}(\mathrm{X},\mathbf{Q})
\to
\mathrm{H}^{n-i}(\mathrm{X},\mathbf{Q})\otimes\mathrm{H}^{i}(\mathrm{G},\mathbf{Q})
\]
is zero. Since \(\mathrm{H}^{n}(\mathrm{X},\mathbf{Q})\) is pure of weight \(n\) and \(\mathrm{H}^{n-i}(\mathrm{X},\mathbf{Q})\) is pure of weight \(n-i\), we have
\[
\operatorname{Im}(\mathrm{H}^{n}(\mathrm{X},\mathbf{Q}))
\subset
\mathrm{W}_{n}(\mathrm{H}^{n-i}(\mathrm{X},\mathbf{Q})\otimes\mathrm{H}^{i}(\mathrm{G},\mathbf{Q}))
=
\mathrm{H}^{n-i}(\mathrm{X},\mathbf{Q})\otimes\mathrm{W}_{i}\mathrm{H}_{i}(\mathrm{G},\mathbf{Q}).
\]
By (9.1.6), this image is therefore zero.

\textit{Variant} \textbf{(9.1.8)}. — \textit{Let \(\mathrm{X}\) be a separated algebraic space on which a connected linear algebraic group \(\mathrm{G}\) acts.}

(i) \textit{If \(\mathrm{X}\) is smooth and \(\overline{\mathrm{X}}\) is a compactification of \(\mathrm{X}\), the diagram}
\[
\begin{tikzcd}[ampersand replacement=\&, column sep=1.8em, row sep=3.5em]
\mathrm{H}^{\bullet}(\overline{\mathrm{X}},\mathbf{Q})
  \arrow[r] \arrow[d]
\&
\mathrm{H}^{\bullet}(\mathrm{X},\mathbf{Q})
  \arrow[r, "\rho^{*}"]
\&
\mathrm{H}^{\bullet}(\mathrm{X},\mathbf{Q})\otimes
\mathrm{H}^{\bullet}(\mathrm{G},\mathbf{Q})
\\
\mathrm{H}^{\bullet}(\mathrm{X},\mathbf{Q})
  \arrow[rr, equal]
\& \&
\mathrm{H}^{\bullet}(\mathrm{X},\mathbf{Q})\otimes
\mathrm{H}^{0}(\mathrm{G},\mathbf{Q})
  \arrow[u, hook']
\end{tikzcd}
\]
\textit{is commutative.}

(ii) \textit{If \(\mathrm{X}\) is proper and \(\widetilde{\mathrm{X}}\) is a resolution of singularities of \(\mathrm{X}\), the diagram}
\[
\begin{tikzcd}[ampersand replacement=\&, column sep=1.5em, row sep=3.5em]
\mathrm{H}^{\bullet}(\mathrm{X},\mathbf{Q})
  \arrow[r, "\rho^{*}"] \arrow[d]
\&
\mathrm{H}^{\bullet}(\mathrm{X},\mathbf{Q})\otimes
\mathrm{H}^{\bullet}(\mathrm{G},\mathbf{Q})
  \arrow[r]
\&
\mathrm{H}^{\bullet}(\widetilde{\mathrm{X}},\mathbf{Q})\otimes
\mathrm{H}^{\bullet}(\mathrm{G},\mathbf{Q})
\\
\mathrm{H}^{\bullet}(\widetilde{\mathrm{X}},\mathbf{Q})
  \arrow[rr, equal]
\& \&
\mathrm{H}^{\bullet}(\widetilde{\mathrm{X}},\mathbf{Q})\otimes
\mathrm{H}^{0}(\mathrm{G},\mathbf{Q})
  \arrow[u, hook']
\end{tikzcd}
\]
\textit{is commutative.}

\Page{48}
It suffices to prove that for \(i>0\), the Künneth components
\[
\begin{aligned}
&\mathrm{H}^{n}(\overline{\mathrm{X}},\mathbf{Q})
\to\mathrm{H}^{n}(\mathrm{X},\mathbf{Q})
\to\mathrm{H}^{n-i}(\mathrm{X},\mathbf{Q})\otimes\mathrm{H}^{i}(\mathrm{G},\mathbf{Q})
\\
&\mathrm{H}^{n}(\mathrm{X},\mathbf{Q})
\to\mathrm{H}^{n-i}(\mathrm{X},\mathbf{Q})\otimes\mathrm{H}^{i}(\mathrm{G},\mathbf{Q})
\to\mathrm{H}^{n-i}(\widetilde{\mathrm{X}},\mathbf{Q})\otimes\mathrm{H}^{i}(\mathrm{G},\mathbf{Q})
\end{aligned}
\]
are respectively zero. In case (i) (resp. (ii)), we have (8.2.4)
\[
\mathrm{W}_{n}\mathrm{H}^{n}(\overline{\mathrm{X}},\mathbf{Q})
=\mathrm{H}^{n}(\overline{\mathrm{X}},\mathbf{Q})
\quad
(\text{resp. }\mathrm{W}_{n}\mathrm{H}^{n}(\mathrm{X},\mathbf{Q})
=\mathrm{H}^{n}(\mathrm{X},\mathbf{Q})),
\]
whereas, by (8.2.4) and (9.1.6),
\[
\mathrm{W}_{n}(\mathrm{H}^{n-i}(\mathrm{X},\mathbf{Q})\otimes\mathrm{H}^{i}(\mathrm{G},\mathbf{Q}))=0
\quad
(\text{resp. }\mathrm{W}_{n}(\mathrm{H}^{n-i}(\widetilde{\mathrm{X}},\mathbf{Q})\otimes\mathrm{H}^{i}(\mathrm{G},\mathbf{Q}))=0)
\]
whence the assertion.

\section*{9.2. Hodge theory of smooth hypersurfaces, after Griffiths.}

This section depends only on §§ 1 to 3. Here we explain Theorem (8.6) of [7].

\textbf{(9.2.1)} Let \(\mathrm{X}\) be a proper smooth algebraic variety, and let \(\mathrm{Y}\) be a smooth divisor in \(\mathrm{X}\). Let \(j\) be the inclusion of \(\mathrm{U}=\mathrm{X}-\mathrm{Y}\) in \(\mathrm{X}\). We have \(\mathrm{R}^{i}j_{*}\mathbf{Z}=0\) for \(i\neq0,\ 1\), and the Leray spectral sequence for \(j\) is identified with the long exact cohomology sequence
\begin{align*}
\cdots\to\mathrm{H}_{\mathrm{Y}}^{n}(\mathrm{X})
&\to\mathrm{H}^{n}(\mathrm{X})
\to\mathrm{H}^{n}(\mathrm{U})
\xrightarrow{\partial}\cdots
\tag*{\textbf{(9.2.1.1)}}\\
\mathrm{H}_{\mathrm{Y}}^{n}(\mathrm{X})
&=\mathrm{H}^{n-2}(\mathrm{Y},\mathscr{H}_{\mathrm{Y}}^{2}).
\end{align*}
For a reason made clear in (3.1.9), we shall identify \(\mathscr{H}_{\mathrm{Y}}^{2}\) with the constant sheaf \(\mathbf{Z}(-1)_{\mathrm{Z}}\). Since the Leray spectral sequence of \(j\) converges by definition to the mixed Hodge structure on \(\mathrm{H}^{n}(\mathrm{U})\), we find that the exact sequence
\[
\cdots\to
\mathrm{H}^{n-2}(\mathrm{Y},\mathbf{Z})(-1)
\to\mathrm{H}^{n}(\mathrm{X},\mathbf{Z})
\to\mathrm{H}^{n}(\mathrm{U},\mathbf{Z})
\xrightarrow{\partial}\cdots
\tag*{\textbf{(9.2.1.2)}}
\]
is an exact sequence of mixed Hodge structures.

The Leray spectral sequence of \(j\) is also the spectral sequence of the filtered complex \((\Omega_{\mathrm{X}}^{\bullet}(\log\mathrm{Y}),\mathrm{W})\); the preceding exact sequence is therefore also identified with the one deduced from the short exact sequence
\[
0\to\Omega_{\mathrm{X}}^{\bullet}
\to\Omega_{\mathrm{X}}^{\bullet}(\log\mathrm{Y})
\xrightarrow{\operatorname{Res}}
\Omega_{\mathrm{Y}}^{\bullet}\to0.
\tag*{\textbf{(9.2.1.3)}}
\]
For this reason, the homomorphism \(\partial\) in (9.2.1.2) is also called the \textit{residue operator}.

\textbf{(9.2.2)} The residue operator
\[
\operatorname{Res}:\mathrm{H}^{n+1}(\mathrm{U},\mathbf{Z})
\to\mathrm{H}^{n}(\mathrm{Y},\mathbf{Z})(-1)
\tag*{\textbf{(9.2.2.1)}}
\]
is strictly compatible with the Hodge filtration. The Hodge filtration on \(\mathrm{H}^{n}(\mathrm{U},\mathbf{C})\) therefore determines that on \(\operatorname{Ker}(\mathrm{H}^{n}(\mathrm{Y},\mathbf{C})\to\mathrm{H}^{n+2}(\mathrm{X},\mathbf{C}))\).

\Page{49}
Suppose that \(\mathrm{Y}\) is ample. If \(\mathrm{X}\) is of pure dimension \(\mathrm{N}+1\), we then have \(\mathrm{H}^{i}(\mathrm{X},\mathbf{Z})\xrightarrow{\sim}\mathrm{H}^{i}(\mathrm{Y},\mathbf{Z})\) for \(i<\mathrm{N}\) (Lefschetz). By Poincaré duality, for \(i\neq\mathrm{N}\), \(\mathrm{H}^{i}(\mathrm{Y},\mathbf{C})\) can therefore be computed from \(\mathrm{H}^{\bullet}(\mathrm{X},\mathbf{C})\). Put
\[
\begin{aligned}
\mathrm{H}_{\text{év}}^{\mathrm{N}}(\mathrm{Y},\mathbf{C})
&=\operatorname{Ker}(\mathrm{H}^{\mathrm{N}}(\mathrm{Y},\mathbf{C})\to\mathrm{H}^{\mathrm{N}+2}(\mathrm{X},\mathbf{C}))\\
&=\text{orthogonal complement of }\mathrm{H}^{\mathrm{N}}(\mathrm{X},\mathbf{C})
\hookrightarrow\mathrm{H}^{\mathrm{N}}(\mathrm{Y},\mathbf{C})
\end{aligned}
\]
(vanishing part of the cohomology). We have
\[
\mathrm{H}^{\mathrm{N}}(\mathrm{Y},\mathbf{C})
=\mathrm{H}_{\text{év}}^{\mathrm{N}}(\mathrm{Y},\mathbf{C})
\oplus\mathrm{H}^{\mathrm{N}}(\mathrm{X},\mathbf{C})
\]
and we thus obtain a great deal of information about the Hodge structure on \(\mathrm{H}^{\bullet}(\mathrm{Y},\mathbf{C})\) when we know that on \(\mathrm{H}^{\bullet}(\mathrm{X},\mathbf{C})\) and the Hodge filtration on \(\mathrm{H}^{\mathrm{N}+1}(\mathrm{U},\mathbf{C})\).

\textbf{(9.2.3)} By definition, the Hodge filtration on \(\mathrm{H}^{\bullet}(\mathrm{U},\mathbf{C})\) is the abutment of the spectral sequence \textit{degenerating} at \(\mathrm{E}_{1}\)
\[
\mathrm{E}_{1}^{pq}
=\mathrm{H}^{q}(\mathrm{X},\Omega_{\mathrm{X}}^{p}(\log\mathrm{Y}))
\Rightarrow\mathrm{H}^{p+q}(\mathrm{U},\mathbf{C}).
\]
By (3.1.11), this spectral sequence coincides with the spectral sequence of the complex \(j_{*}\Omega_{\mathrm{U}}^{\bullet}\), equipped with the pole order filtration \(\mathrm{P}\).

In particular, we have
\[
\mathrm{F}^{p}(\mathrm{H}^{\bullet}(\mathrm{U},\mathbf{C}))
=\mathbf{H}^{\bullet}(\mathrm{X},\mathrm{P}^{p}(j_{*}\Omega_{\mathrm{U}}^{\bullet})).
\]
Recall that \(\mathrm{P}^{p}(j_{*}\Omega_{\mathrm{U}}^{\bullet})\) is the following complex, whose first nonzero component is in degree \(p\)
\[
\Omega_{\mathrm{X}}^{p}(\mathrm{Y})
\to\Omega_{\mathrm{X}}^{p+1}(2\mathrm{Y})
\to\cdots\to
\Omega_{\mathrm{X}}^{i}((i-p+1)\mathrm{Y})
\to\cdots
\]
If \(\mathrm{H}^{i}(\mathrm{X},\Omega_{\mathrm{X}}^{p}(n\mathrm{Y}))=0\) for \(i>0\), \(n>0\) and \(p\geq0\), we simply have
\[
\mathbf{H}^{\bullet}(\mathrm{X},\mathrm{P}^{p}(j_{*}\Omega_{\mathrm{U}}^{\bullet}))
=\mathrm{H}^{\bullet}(\Gamma(\mathrm{X},\mathrm{P}^{p}(j_{*}\Omega_{\mathrm{U}}^{\bullet}))),
\]
whence the following result :

\textit{Proposition} \textbf{(9.2.4)}. — \textit{Suppose that \(\mathrm{H}^{i}(\mathrm{X},\Omega_{\mathrm{X}}^{p}(n\mathrm{Y}))=0\) for \(i>0\), \(n>0\) and \(p\geq0\) (this is the case if \(\mathrm{Y}\) is sufficiently ample). Then a cohomology class \(c\in\mathrm{H}^{d}(\mathrm{U},\mathbf{C})\) is of Hodge filtration \(\geq p\) if and only if it can be represented by a closed \(d\)-form \(\alpha\) on \(\mathrm{U}\), with a pole of order \(\leq d-p+1\) along \(\mathrm{Y}\). If \(\alpha\) is a closed form on \(\mathrm{U}\), with a pole of order \(k\) along \(\mathrm{Y}\), and if the class of \(\alpha\) in \(\mathrm{H}^{d}(\mathrm{U},\mathbf{C})\) is zero, then \(\alpha=d\beta\), with \(\beta\) having a pole of order \(\leq k-1\) along \(\mathrm{Y}\); for \(k\leq1\), we even have \(\alpha=0\).}

\textbf{(9.2.5)} It is well known (Bott’s theorem) that the hypotheses of (9.2.4) are satisfied when \(\mathrm{Y}\) is a hypersurface in projective space \(\mathbf{P}^{n}(\mathbf{C})\) \((n\geq1)\), i.e. that
\[
\mathrm{H}^{i}(\mathbf{P}^{n}(\mathbf{C}),\Omega^{j}(m))=0
\qquad\text{for}\qquad i>0,\ m>0
\]
(a proof can be found, for example, in SGA 7 XI). Moreover, for \(n>1\), the residue map identifies \(\mathrm{H}^{n}(\mathrm{U},\mathbf{C})\) with the primitive part of the cohomology of \(\mathrm{H}^{n-1}(\mathrm{Y},\mathbf{C})\). Hence ([7], (8.6)) :

\Page{50}
\textit{Proposition} \textbf{(9.2.6)}. — \textit{Let \(\mathrm{Y}\) be a smooth hypersurface in \(\mathbf{P}^{n}(\mathbf{C})\) \((n>1)\), and \(\mathrm{U}=\mathbf{P}^{n}(\mathbf{C})-\mathrm{Y}\).}

(i) \textit{For a primitive cohomology class \(c\in\mathrm{H}^{n-1}(\mathrm{Y},\mathbf{C})\) to be of Hodge filtration \(\leq p\), it is necessary and sufficient that it be the residue of a regular differential \(n\)-form on \(\mathrm{U}\), with a pole of order \(\leq n-p\) along \(\mathrm{Y}\).}

(ii) \textit{Let \(\alpha\) be a regular differential \(n\)-form on \(\mathrm{U}\), with a pole of order \(k\geq1\) along \(\mathrm{Y}\). For \(\operatorname{Res}(\alpha)\in\mathrm{H}^{n-1}(\mathrm{Y},\mathbf{C})\) to be zero, it is necessary and sufficient that \(\alpha=d\beta\), for an \((n-1)\)-form on \(\mathrm{U}\), with a pole of order \(\leq k-1\) along \(\mathrm{Y}\).}

\section*{9.3. Construction of complexes of 1\textsuperscript{st}-order differential operators.}

This section depends only on §§ 5 and 6. Here we prove the following theorem.

\textit{Proposition} \textbf{(9.3.1)}. — \textit{Let \(\mathrm{X}\) be a quasi-projective scheme (over \(\mathbf{C}\)). There exists on \(\mathrm{X}\) a complex of sheaves \(\mathrm{K}\) with the following properties.}

\textit{a) The \(\mathrm{K}^{n}\) are coherent algebraic sheaves and \(\mathrm{K}^{n}=0\) for \(n<0\).}

\textit{b) The differentials \(d:\mathrm{K}^{n}\to\mathrm{K}^{n+1}\) are 1\textsuperscript{st}-order algebraic differential operators.}

\textit{c) The complex of coherent analytic sheaves \(\mathrm{K}^{\mathrm{an}}\) on \(\mathrm{X}^{\mathrm{an}}\) is a resolution of the constant sheaf \(\underline{\mathbf{C}}\). More precisely, \(\mathscr{H}^{i}(\mathrm{K})=0\) for \(i>0\), and there exists}
\[
f:\Omega_{\mathrm{X}}^{\bullet}\to\mathrm{K}
\]
\textit{with \(f^{n}\) algebraic and linear, such that the composite map \(\underline{\mathbf{C}}\to\mathscr{O}_{\mathrm{X}}\xrightarrow{f^{0}}\mathrm{K}^{0}\) identifies \(\underline{\mathbf{C}}\) with \(\mathscr{H}^{0}(\mathrm{K})\).}

\textit{d) The natural map}
\[
\mathbf{H}^{i}(\mathrm{X},\mathrm{K})
\to\mathbf{H}^{i}(\mathrm{X}^{\mathrm{an}},\mathrm{K}^{\mathrm{an}})
\xleftarrow[(c)]{\sim}
\mathbf{H}^{i}(\mathrm{X}^{\mathrm{an}},\mathbf{C})
\]
\textit{is an isomorphism.}

As a corollary, we recover a special case of a theorem of Bloom and Herrera [2].

\textit{Corollary} \textbf{(9.3.2)}. — \textit{The map induced by \(\underline{\mathbf{C}}\to\Omega_{\mathrm{X}}^{\bullet\mathrm{an}}\) :}
\[
\mathrm{H}^{\bullet}(\mathrm{X}^{\mathrm{an}},\mathbf{C})
\to\mathbf{H}^{\bullet}(\mathrm{X}^{\mathrm{an}},\Omega_{\mathrm{X}}^{\bullet\mathrm{an}})
\]
\textit{identifies \(\mathrm{H}^{\bullet}(\mathrm{X}^{\mathrm{an}},\mathbf{C})\) with a direct summand of \(\mathbf{H}^{\bullet}(\mathrm{X}^{\mathrm{an}},\Omega_{\mathrm{X}}^{\bullet\mathrm{an}})\).}

We shall first prove the corollary. The proposition will be obtained by repeating the proof at the level of complexes. Let \(\varepsilon:\mathrm{X}_{\bullet}\to\mathrm{X}\) be a simplicial scheme augmented over \(\mathrm{X}\). We assume that \(\varepsilon\) is a proper hypercovering and that the \(\mathrm{X}_{n}\) are smooth. We have a commutative diagram
\[
\begin{tikzcd}[ampersand replacement=\&, column sep=large, row sep=3em]
\mathrm{H}^{\bullet}(\mathrm{X}^{\mathrm{an}},\mathbf{C})
  \arrow[r, "\sim", "\text{\textcircled{1}}"']
  \arrow[d]
\&
\mathrm{H}^{\bullet}(\mathrm{X}_{\bullet}^{\mathrm{an}},\mathbf{C})
  \arrow[d, "\wr", "\text{\textcircled{2}}"']
\\
\mathbf{H}^{\bullet}(\mathrm{X}^{\mathrm{an}},\Omega_{\mathrm{X}}^{\bullet})
  \arrow[r]
\&
\mathbf{H}^{\bullet}(\mathrm{X}_{\bullet}^{\mathrm{an}},\Omega_{\mathrm{X}_{\bullet}}^{\bullet})
\end{tikzcd}
\tag*{\textbf{(9.3.2.1)}}
\]

\Page{51}
Arrow 1 is bijective by (5.3.5) (II) (IV). Arrow 2 is bijective, because \(\Omega_{\mathrm{X}_{\bullet}}^{\bullet}\) is a resolution of \(\mathbf{C}\) (Poincaré lemma on the \(\mathrm{X}_{n}\)). Diagram (9.3.2.1) provides a retraction, compatible with the cup-product, of \(\mathbf{H}^{\bullet}(\mathrm{X}^{\mathrm{an}},\Omega_{\mathrm{X}}^{\bullet})\) onto \(\mathrm{H}^{\bullet}(\mathrm{X}^{\mathrm{an}},\mathbf{C})\), and this proves (9.3.2).

\textbf{(9.3.3)} The complex \(\mathrm{K}\) that we wish to construct will be none other, in the derived category, than \(\mathbf{s}\mathrm{R}\varepsilon_{\bullet *}\Omega_{\mathrm{X}_{\bullet}}^{\bullet}\). We shall treat only the case where the augmented simplicial space \(\mathrm{X}_{\bullet}\) is \(s\)-split, defined as in (6.2.5) by proper surjective maps whose sources are quasi-projective and smooth
\[
f'_{k}:\mathrm{N}_{k}\to
(\operatorname{cosq}\operatorname{sq}_{k-1}(\mathrm{X}_{\bullet}))_{k}
\qquad(k\geq0).
\]
We shall construct a “compatible system” of affine open coverings of the \(\mathrm{X}_{n}\). For this construction, it is useful to regard an open covering of a space \(\mathrm{Y}\) as a diagram
\[
\begin{tikzcd}[ampersand replacement=\&, column sep=1.4em]
\mathrm{Y}
\&
\mathrm{U}\arrow[l, "u"']\arrow[r, "\varphi"]
\&
\mathrm{I}
\end{tikzcd}
\]
with \(\mathrm{I}\) discrete, \(u\) surjective, and \(u_{i}:\mathrm{U}_{i}\underset{\mathrm{dfn}}{=}\varphi^{-1}(i)\hookrightarrow\mathrm{Y}\) an open immersion. Proceeding as in (6.2.5), we construct a diagram
\[
\begin{tikzcd}[ampersand replacement=\&, column sep=1.4em]
\mathrm{X}_{\bullet}
\&
\mathrm{U}_{\bullet}\arrow[l, "u"']\arrow[r, "\varphi"]
\&
\mathrm{I}_{\bullet}
\end{tikzcd}
\tag*{\textbf{(9.3.3.1)}}
\]
of the following type.

\textit{a)} \(\mathrm{I}_{\bullet}\) is a discrete simplicial set. The \(\mathrm{I}_{n}\) are finite.

\textit{b)} For every \(k\geq0\), (9.3.3.1) induces a diagram
\[
\begin{tikzcd}[ampersand replacement=\&, column sep=1.2em, row sep=3em]
(\operatorname{cosq}\operatorname{sq}_{k-1}(\mathrm{X}_{\bullet}))_{k}
\&
(\operatorname{cosq}\operatorname{sq}_{k-1}(\mathrm{U}_{\bullet}))_{k}
  \arrow[l]\arrow[r, "\varphi"]
\&
(\operatorname{cosq}\operatorname{sq}_{k-1}(\mathrm{I}_{\bullet}))_{k}
\\
\mathrm{N}_{k}\arrow[u, "f'_{k}"']
\&
\mathrm{N}_{k}(\mathrm{U}_{\bullet})
  \arrow[l]\arrow[u]\arrow[r, "\varphi"]
\&
\mathrm{N}_{k}(\mathrm{I}_{\bullet})\arrow[u, "a'_{k}"']
\end{tikzcd}
\]

\textit{c)} In diagram \textit{b)}, for \(i\in\operatorname{cosq}\operatorname{sq}_{k-1}(\mathrm{I}_{\bullet})_{k}\), the \(\varphi^{-1}(j)\) for \(a'_{k}(j)=i\) form an affine open covering of
\[
\mathrm{N}_{k}
\times_{(\operatorname{cosq}\operatorname{sq}_{k-1}(\mathrm{X}_{\bullet}))_{k}}
\varphi^{-1}(i).
\]

\textit{d)} (follows by induction from \textit{c)}). The first (and hence the second) row of diagram \textit{b)} defines an open covering. The \(\mathrm{X}_{k}\leftarrow\mathrm{U}_{k}\to\mathrm{I}_{k}\) therefore also form open coverings.

\textit{e)} (follows from a suitable choice of the affine coverings in \textit{c)}). \(\mathrm{U}_{n_i}=\varphi_{n}^{-1}(i)\) is the complement of a subscheme \(\mathrm{D}_{n,i}\) of \(\mathrm{X}_{n}\), the union of some connected components of \(\mathrm{X}_{n}\) and a divisor. For any coherent sheaf \(\mathscr{F}\) on \(\mathrm{X}_{n}\) and integers \((a_{i})_{i\in\mathrm{I}_{n}}\) \((a_{i}\geq0)\), put
\[
\mathscr{F}\Bigl(\sum_{i}a_{i}\mathrm{D}_{n,i}\Bigr)
=
\begin{cases}
0\text{ on the components of }\mathrm{X}_{n}
 \text{ contained in a }\mathrm{D}_{n,i}\text{ with }a_{i}>0\\[3pt]
\displaystyle
\mathscr{F}\otimes\mathscr{O}\Bigl(\sum_{a_{i}>0}a_{i}\mathrm{D}_{n,i}\Bigr)
 \text{ on the others.}
\end{cases}
\]

\Page{52}
In the inductive construction of (9.3.3.1), the open coverings (\textit{c}) are chosen arbitrarily. This allows us to assume

\textit{f)} For any \(a_{i}>0\) \((i\in\mathrm{I}_{k})\)
\[
\mathrm{R}^{\ell}\varepsilon_{k*}
\Bigl(\Omega_{\mathrm{X}_{k}}^{m}\Bigl(\sum_{i}a_{i}\mathrm{D}_{k,i}\Bigr)\Bigr)=0
\qquad\text{for}\qquad\ell>0.
\]

\textit{g)} For any \(\alpha:\Delta_{m}\to\Delta_{n}\) and \(i\in\mathrm{I}_{n}\), we have
\[
\mathrm{X}_{\bullet}(\alpha)^{-1}(\mathrm{D}_{m,\alpha(i)})
\subset\mathrm{D}_{n,i}
\qquad\text{(as a scheme).}
\]
For each \(k\), the complex of sheaves on \(\mathrm{X}\)
\[
\mathrm{R}\varepsilon_{k*}(\Omega_{\mathrm{X}_{k}}^{\bullet})
\]
can be computed by Čech methods. If, for \(\mathrm{P}\subset\mathrm{I}_{k}\), \(j_{\mathrm{P}}\) is the inclusion of \(\mathrm{U}_{\mathrm{P}}=\bigcap_{i\in\mathrm{P}}\mathrm{U}_{k_i}\) in \(\mathrm{X}_{k}\), it is the simple complex associated with the double complex with components
\[
\check{\mathrm{K}}^{k,p,q}
=\sum_{\#\mathrm{P}=p+1}
(\varepsilon_{k}j_{\mathrm{P}})_{*}j_{\mathrm{P}}^{*}\Omega_{\mathrm{X}_{k}}^{q}.
\]
This double complex has as a double subcomplex the complex whose components are the coherent sheaves
\[
\mathrm{K}^{k,p,q}
=\sum_{\#\mathrm{P}=p+1}
\varepsilon_{k*}\Omega_{\mathrm{X}_{k}}^{q}
\Bigl((q+1)\sum_{i\in\mathrm{P}}\mathrm{D}_{k,i}\Bigr).
\]
For fixed \(k\) and \(q\), it follows from \textit{f)} that the inclusion of \(\mathrm{K}^{k,\bullet,q}\) in \(\check{\mathrm{K}}^{k,\bullet,q}\) is a quasi-isomorphism. For fixed \(k\), the inclusion of \(\mathrm{K}^{k,\bullet,\bullet}\) in \(\check{\mathrm{K}}^{k,\bullet,\bullet}\) is also a quasi-isomorphism. Finally, for variable \(k\), the \(\mathrm{K}^{k,\bullet,\bullet}\) form, by \textit{g)}, a simplicial double complex. The associated simple complex \(\mathrm{K}\) coincides with \(\mathrm{R}\varepsilon_{\bullet *}\Omega_{\mathrm{X}_{\bullet}}^{\bullet}\), in the derived category.

Let us prove that \(\mathrm{K}\) satisfies (9.3.1). Properties \textit{a)} and \textit{b)} of \(\mathrm{K}\) are clear; we define \(f:\Omega_{\mathrm{X}}^{\bullet}\to\mathrm{K}\) by the obvious maps from the \(\Omega_{\mathrm{X}}^{q}\) to the \(\mathrm{K}^{0,0,q}\).

In the analytic category, the complex of coherent analytic sheaves \(\mathrm{K}^{\mathrm{an}}\) is given by the same formulas as \(\mathrm{K}\). We have a commutative diagram
\[
\begin{tikzcd}[ampersand replacement=\&, column sep=large, row sep=3em]
\mathbf{H}^{i}(\mathrm{X},\mathrm{K})
  \arrow[r, dash, "\sim"]\arrow[d]
\&
\mathbf{H}^{i}(\mathrm{X}_{\bullet},\Omega_{\mathrm{X}_{\bullet}}^{\bullet})
  \arrow[d, "\text{\textcircled{1}}"]
\\
\mathbf{H}^{i}(\mathrm{X}^{\mathrm{an}},\mathrm{K}^{\mathrm{an}})
  \arrow[r, dash, "\sim"]
\&
\mathbf{H}^{i}(\mathrm{X}_{\bullet}^{\mathrm{an}},
              \Omega_{\mathrm{X}_{\bullet}^{\mathrm{an}}}^{\bullet})
\end{tikzcd}
\]
and 1 is the abutment of a morphism of spectral sequences
\[
\begin{tikzcd}[ampersand replacement=\&, column sep=large, row sep=3em]
\mathbf{H}^{a}(\mathrm{X}_{b},\Omega_{\mathrm{X}_{\bullet}}^{\bullet})
  \arrow[r, Rightarrow]\arrow[d]
\&
\mathbf{H}^{a+b}(\mathrm{X}_{\bullet},\Omega_{\mathrm{X}_{\bullet}}^{\bullet})
  \arrow[d]
\\
\mathbf{H}^{a}(\mathrm{X}_{b}^{\mathrm{an}},
              \Omega_{\mathrm{X}_{\bullet}^{\mathrm{an}}}^{\bullet})
  \arrow[r, Rightarrow]
\&
\mathbf{H}^{a+b}(\mathrm{X}_{\bullet}^{\mathrm{an}},
                \Omega_{\mathrm{X}_{\bullet}^{\mathrm{an}}}^{\bullet}).
\end{tikzcd}
\]
This morphism is an isomorphism by [9], whence \textit{d)}.

% ---- 053-060 ----
\Page{53}
To prove \textit{c)}, we consider the commutative diagram
\[
\begin{tikzcd}[column sep=2.4em, row sep=3.5em]
\underline{\mathbf{C}}
  \arrow[rr, dash, "\approx", "{\text{\textcircled{1}}}"'{pos=.2}]
  \arrow[d]
& &
\mathrm{R}\varepsilon_{\bullet *}\underline{\mathbf{C}}
  \arrow[d, "\approx", "{\text{\textcircled{2}}}"'{pos=.2}]
\\
\Omega_{\mathrm{X}_{\bullet}}^{\bullet\mathrm{an}}
  \arrow[r, "f"]
&
\mathrm{K}^{\mathrm{an}}
  \arrow[r, dash, "\approx"]
&
\mathrm{R}\varepsilon_{\bullet *}
  \Omega_{\mathrm{X}_{\bullet}}^{\bullet\mathrm{an}}
\end{tikzcd}
\]
That 1 is a quasi-isomorphism follows from (5.3.5).\par
That 2 is one follows from the Poincaré lemma on the \(\mathrm{X}_n\).
Assertion \textit{c)} follows.

\textbf{(9.3.4)} The same construction yields resolutions satisfying
\textit{a)}, \textit{b)} for sheaves \(\mathscr{F}\) on
\(\mathrm{X}^{\mathrm{an}}\) other than the constant sheaf \(\mathbf{C}\)
(one can probably take \(\mathscr{F}\) to be any
“algebraically constructible” sheaf). Unfortunately, I cannot
prove that the resolutions obtained are essentially unique
(in a suitable derived category). In the absence of such
uniqueness, (9.3.1) is of little interest.

\section*{10. HODGE THEORY OF LEVEL \(\leq\mathbf{1}\)}

\section*{10.1. 1-motives}

\textit{Definition} \textbf{(10.1.1)}. — (i)
\textit{A mixed Hodge structure \(\mathrm{H}\) is
\textup{torsion-free} if \(\mathrm{H}_{\mathbf{Z}}\) is torsion-free.}

(ii) \textit{If \(\mathrm{H}\) is torsion-free, we again denote by
\(\mathrm{W}\) the filtration of
\(\mathrm{H}_{\mathbf{Z}}\subset\mathrm{H}_{\mathbf{Q}}\) induced by
\(\mathrm{W}\). We denote by
\(\operatorname{Gr}^{\mathrm{W}}_n(\mathrm{H})\) the
torsion-free \(\mathbf{Z}\)-module
\(\operatorname{Gr}^{\mathrm{W}}_n(\mathrm{H}_{\mathbf{Z}})\), equipped with
its Hodge structure of weight \(n\).}

\textit{Definition} \textbf{(10.1.2)}. —
\textit{A \textup{1-motive} over an algebraically closed field \(k\)
consists of:}

\noindent a) \textit{a finitely generated free \(\mathbf{Z}\)-module
\(\mathrm{X}\), an abelian variety \(\mathrm{A}\) and a torus
\(\mathrm{T}\);}\par
\noindent b) \textit{an extension \(\mathrm{G}\) of \(\mathrm{A}\)
by \(\mathrm{T}\);}\par
\noindent c) \textit{a homomorphism
\(u:\mathrm{X}\to\mathrm{G}(k)\).}

It is often useful to regard a 1-motive \(\mathrm{M}\) as
a complex of group schemes, concentrated in degrees 0 and 1
\[
\mathrm{M}=[\mathrm{X}\xrightarrow{u}\mathrm{G}].
\]

\textit{Construction} \textbf{(10.1.3)}. —
\textit{We shall construct an equivalence
\(\mathrm{M}\mapsto\mathrm{T}(\mathrm{M})\) between the category of
1-motives over \(\mathbf{C}\) and the category of torsion-free mixed Hodge
structures \(\mathrm{H}\) of type}
\[
\{(0,0),(0,-1),(-1,0),(-1,-1)\}
\]
\Page{54}
\textit{such that \(\operatorname{Gr}^{\mathrm{W}}_{-1}(\mathrm{H})\)
is polarizable. Moreover, we shall construct isomorphisms natural in
\(\mathrm{M}=(\mathrm{X},\mathrm{A},\mathrm{T},\mathrm{G},u)\):}
\[
\begin{aligned}
\mathrm{H}_1(\mathrm{T},\mathbf{Z})
  &\simeq \operatorname{Gr}^{\mathrm{W}}_{-2}
          \mathrm{T}(\mathrm{M})_{\mathbf{Z}}
\\
\mathrm{H}_1(\mathrm{A},\mathbf{Z})
  &\simeq \operatorname{Gr}^{\mathrm{W}}_{-1}\mathrm{T}(\mathrm{M})
  \qquad \text{\textit{(isomorphism of Hodge structures)}}
\\
\mathrm{X}
  &\simeq \operatorname{Gr}^{\mathrm{W}}_0
          \mathrm{T}(\mathrm{M})_{\mathbf{Z}}
\end{aligned}
\]

Let \(\mathrm{M}\) be a 1-motive. The exponential map
\(\exp:\operatorname{Lie}(\mathrm{G})\to\mathrm{G}\) has kernel
\(\mathrm{H}_1(\mathrm{G})\). We define
\(\mathrm{T}_{\mathbf{Z}}(\mathrm{M})\) (which we shall denote instead by
\(\mathrm{T}_{\mathrm{Z}}(\mathrm{M})\)) to be the fiber product
of \(\operatorname{Lie}(\mathrm{G})\) and \(\mathrm{X}\) over
\(\mathrm{G}\)
\[
\tag*{\textbf{(10.1.3.1)}}
\begin{tikzcd}[column sep=1.5em, row sep=3.5em]
0 \arrow[r]
& \mathrm{H}_1(\mathrm{G}) \arrow[r]
& \operatorname{Lie}(\mathrm{G}) \arrow[r, "\exp"]
& \mathrm{G} \arrow[r]
& 0
\\
0 \arrow[r]
& \mathrm{H}_1(\mathrm{G}) \arrow[r] \arrow[u, equal]
& \mathrm{T}_{\mathrm{Z}}(\mathrm{M})
    \arrow[r, "\beta"] \arrow[u, "\alpha"']
& \mathrm{X} \arrow[r] \arrow[u, "u"']
& 0
\end{tikzcd}
\]
We set
\[
\begin{aligned}
\mathrm{W}_{-1}(\mathrm{T}_{\mathrm{Z}}(\mathrm{M}))
  &=\operatorname{Ker}(\beta)=\mathrm{H}_1(\mathrm{G})
\\
\mathrm{W}_{-2}(\mathrm{T}_{\mathrm{Z}}(\mathrm{M}))
  &=\mathrm{H}_1(\mathrm{T})
    =\operatorname{Ker}(\mathrm{H}_1(\mathrm{G})\to\mathrm{H}_1(\mathrm{A})).
\end{aligned}
\]
This defines the \textit{weight filtration}. The morphism
\(\alpha\) extends to
\[
\alpha_{\mathbf{C}}:
\mathrm{T}_{\mathrm{Z}}(\mathrm{M})\otimes\mathbf{C}
\to\operatorname{Lie}(\mathrm{G}).
\]
We set
\(\mathrm{F}^0(\mathrm{T}_{\mathrm{Z}}(\mathrm{M})\otimes\mathbf{C})
=\operatorname{Ker}(\alpha_{\mathbf{C}})\), and this defines the
\textit{Hodge filtration} of
\[
\mathrm{T}_{\mathbf{C}}(\mathrm{M})
=\mathrm{T}_{\mathrm{Z}}(\mathrm{M})\otimes\mathbf{C}.
\]

\textit{Lemma} \textbf{(10.1.3.2)}. —
\textit{The triple
\(\mathrm{T}(\mathrm{M})=(\mathrm{T}_{\mathrm{Z}}(\mathrm{M}),
\mathrm{W},\mathrm{F})\) is a torsion-free mixed Hodge structure
of type \(\{(0,0)(0,-1)(-1,0)(-1,-1)\}\) and
\(\operatorname{Gr}^{\mathrm{W}}_{-1}(\mathrm{T}(\mathrm{M}))\)
is polarizable.}
\[
\begin{tikzcd}[column sep=1.1em, row sep=2.3em]
& & 0 \arrow[d] & 0 \arrow[d] &
\\
& &
\mathrm{W}_{-1}(\mathrm{T}_{\mathbf{C}}(\mathrm{M}))\cap\mathrm{F}^0
  \arrow[r, "\sim", "{\text{\textcircled{2}}}"']
  \arrow[d]
&
\operatorname{Ker}(\alpha_{\mathrm{A},\mathbf{C}})
  \arrow[d]
&
\\
0 \arrow[r]
&
\mathrm{H}_1(\mathrm{T},\mathbf{Z})\otimes\mathbf{C}
  \arrow[r]
  \arrow[d, "\approx", "{\text{\textcircled{1}}}"']
&
\mathrm{H}_1(\mathrm{G},\mathbf{Z})\otimes\mathbf{C}
  \arrow[r]
  \arrow[d, "{\alpha_{\mathrm{G},\mathbf{C}}}"]
&
\mathrm{H}_1(\mathrm{A},\mathbf{Z})\otimes\mathbf{C}
  \arrow[r]
  \arrow[d, "{\alpha_{\mathrm{A},\mathbf{C}}}"]
&
0
\\
0 \arrow[r]
&
\operatorname{Lie}(\mathrm{T}) \arrow[r]
&
\operatorname{Lie}(\mathrm{G}) \arrow[r] \arrow[d]
&
\operatorname{Lie}(\mathrm{A}) \arrow[r] \arrow[d]
&
0
\\
& & 0 & 0 &
\end{tikzcd}
\]

\Page{55}
\textit{a)} We have
\(\mathrm{W}_{-2}(\mathrm{T}_{\mathbf{C}}(\mathrm{M}))
\cap\mathrm{F}^0=0\), since 1 is injective: hence,
\((\mathrm{W}_{-2}(\mathrm{T}_{\mathbf{C}}(\mathrm{M})),\mathrm{F})\)
is a Hodge structure of type \((-1,-1)\).

\textit{b)} Since 2 is surjective, \(\mathrm{F}\) induces on
\(\operatorname{Gr}^{\mathrm{W}}_{-1}
(\mathrm{T}_{\mathrm{Z}}(\mathrm{M}))\otimes\mathbf{C}
\simeq\mathrm{H}_1(\mathrm{A},\mathbf{Z})\otimes\mathbf{C}\)
the Hodge filtration of
\(\mathrm{H}_1(\mathrm{A},\mathbf{Z})\otimes\mathbf{C}\):
this is a polarizable Hodge structure of type
\((-1,0)+(0,-1)\) (cf. (4.4.3)).

\textit{c)}
\[
\begin{tikzcd}[column sep=1.2em, row sep=2.8em]
& 0 \arrow[d] & 0 \arrow[d] & &
\\
0 \arrow[r]
&
\mathrm{W}_{-1}(\mathrm{T}_{\mathbf{C}}(\mathrm{M}))\cap\mathrm{F}^0
  \arrow[r] \arrow[d]
&
\mathrm{F}^0 \arrow[r] \arrow[d]
&
\mathrm{V} \arrow[r] \arrow[d, "\wr"]
&
0
\\
0 \arrow[r]
&
\mathrm{H}_1(\mathrm{G},\mathbf{Z})\otimes\mathbf{C}
  \arrow[r] \arrow[d]
&
\mathrm{T}_{\mathrm{Z}}(\mathrm{M})\otimes\mathbf{C}
  \arrow[r] \arrow[d]
&
\mathrm{X}\otimes\mathbf{C} \arrow[r]
&
0
\\
&
\operatorname{Lie}(\mathrm{G}) \arrow[r, equal] \arrow[d]
&
\operatorname{Lie}(\mathrm{G}) \arrow[d]
& &
\\
& 0 & 0 & &
\end{tikzcd}
\]
\(\mathrm{F}^0\) maps onto
\(\operatorname{Gr}^{\mathrm{W}}_0(\mathrm{T}_{\mathbf{C}}(\mathrm{M}))\),
which is therefore of type \((0,0)\).\par
This completes the construction of \(\mathrm{T}(\mathrm{M})\).
It is clear that \(\mathrm{T}(\mathrm{M})\) is functorial in
\(\mathrm{M}\).

Let \(\mathrm{H}\) be a torsion-free mixed Hodge structure of type
\[
\{(0,0)\ (0,-1)\ (-1,0)\ (-1,-1)\}.
\]
Suppose that \(\operatorname{Gr}^{\mathrm{W}}_{-1}(\mathrm{H})\)
is polarizable; the complex torus
\[
\mathrm{A}
=\mathrm{H}_{\mathbf{Z}}\backslash
\operatorname{Gr}^{\mathrm{W}}_{-1}(\mathrm{H}_{\mathbf{C}})
/
\mathrm{F}^0\operatorname{Gr}^{\mathrm{W}}_{-1}
(\mathrm{H}_{\mathbf{C}})
\]
is then an abelian variety (cf. (4.4.3)). Let \(\mathrm{T}\) be the
torus whose character group is the dual of
\(\operatorname{Gr}^{\mathrm{W}}_{-2}(\mathrm{H}_{\mathbf{Z}})\):
\[
\mathrm{H}_1(\mathrm{T})
=\operatorname{Gr}^{\mathrm{W}}_{-2}(\mathrm{H}_{\mathbf{Z}}).
\]
The complex analytic group
\[
\mathrm{G}
=\mathrm{W}_{-1}(\mathrm{H}_{\mathbf{Z}})\backslash
\mathrm{W}_{-1}(\mathrm{H}_{\mathbf{C}})
/
\mathrm{F}^0\cap\mathrm{W}_{-1}(\mathrm{H}_{\mathbf{C}})
\]
is an extension of \(\mathrm{A}\) by \(\mathrm{T}\).

\textit{Lemma} \textbf{(10.1.3.3)}. —
\textit{The functor \(\mathrm{E}\mapsto\mathrm{E}^{\mathrm{an}}\)
is an equivalence between the category of algebraic groups that are extensions
of \(\mathrm{A}\) by \(\mathrm{T}\) and that of analytic groups
that are extensions of \(\mathrm{A}\) by \(\mathrm{T}\).}

\Page{56}
One reduces to the case where \(\mathrm{T}=\mathbf{G}_m\).
By G.A.G.A., \(\mathrm{E}\mapsto\mathrm{E}^{\mathrm{an}}\) is
then, more generally, an equivalence between the category of principal homogeneous spaces
under \(\mathbf{G}_m\) over \(\mathrm{A}\) and that of principal homogeneous spaces under
\(\mathbf{G}_m\) over \(\mathrm{A}^{\mathrm{an}}\).

We thus have \(\mathrm{G}\), an extension of \(\mathrm{A}\) by
\(\mathrm{T}\). Let
\(\mathrm{X}=\operatorname{Gr}^{\mathrm{W}}_0(\mathrm{H}_{\mathbf{Z}})\).
We define \(u:\mathrm{X}\to\mathrm{G}\) as follows (cf. (2.2.1)):
\[
\begin{tikzcd}[column sep=1.1em, row sep=1.5em]
0 \arrow[r]
&
\mathrm{W}_{-1}(\mathrm{H}_{\mathbf{Z}}) \arrow[r]
&
\mathrm{W}_{-1}(\mathrm{H}_{\mathbf{C}})
/
\mathrm{W}_{-1}(\mathrm{H}_{\mathbf{C}})\cap\mathrm{F}^0
  \arrow[r] \arrow[d, "\wr"]
&
\mathrm{G} \arrow[r]
&
0
\\
& &
\mathrm{H}_{\mathbf{C}}/\mathrm{F}^0
& &
\\
0 \arrow[r]
&
\mathrm{W}_{-1}(\mathrm{H}_{\mathbf{Z}})
  \arrow[r] \arrow[uu, equal]
&
\mathrm{H}_{\mathbf{Z}} \arrow[r] \arrow[u]
&
\mathrm{X} \arrow[r] \arrow[uu]
&
0
\end{tikzcd}
\]
It is clear that the preceding construction, which associates to \(\mathrm{H}\)
a torsion-free 1-motive, is inverse to \(\mathrm{T}\), and is
functorial.

If \(\mathrm{G}_i\) \((i=1,2)\) is an extension of an abelian variety
\(\mathrm{A}_i\) by a torus \(\mathrm{T}_i\), then for every morphism
\(u:\mathrm{G}_1\to\mathrm{G}_2\),
\(u(\mathrm{T}_1)\subset\mathrm{T}_2\) and
\(\operatorname{Ker}(u:\mathrm{A}_1\to\mathrm{A}_2)^0
=\operatorname{Im}(\operatorname{Ker}(u:\mathrm{G}_1\to\mathrm{G}_2)
\to\mathrm{A}_1)^0\).
This corresponds to the fact that a morphism of mixed \(\mathbf{Q}\)-Hodge
structures is strictly compatible with the weight filtration.

\textbf{(10.1.4)} Let \(\mathrm{H}\) be a torsion-free mixed Hodge
structure of type
\[
\{(0,0)\ (0,-1)\ (-1,0)\ (-1,-1)\}.
\]
The filtration \(\mathrm{W}\) of \(\mathrm{H}_{\mathbf{Z}}\) then defines
a filtration of \(\mathrm{H}\) by mixed Hodge
substructures. Similarly, if
\(\mathrm{M}=(\mathrm{X},\mathrm{A},\mathrm{T},\mathrm{G},u)\)
is a 1-motive, we denote by \(\mathrm{W}\) the following increasing
filtration of \(\mathrm{M}\):
\[
\begin{aligned}
\mathrm{W}_i(\mathrm{M})
  &=0 \quad \text{for }i<-2
    \quad\text{and}\quad
    \mathrm{W}_i(\mathrm{M})=\mathrm{M}
    \quad\text{for }i\geq0;
\\
\mathrm{W}_{-1}(\mathrm{M})
  &=\mathrm{G}
    \quad\text{(i.e. }(\{0\},\mathrm{A},\mathrm{T},\mathrm{G},0)\text{)};
\\
\mathrm{W}_{-2}(\mathrm{M})
  &=\mathrm{T}
    \quad\text{(i.e. }(\{0\},0,\mathrm{T},\mathrm{T},0)\text{)}.
\end{aligned}
\]
In an obvious sense, the successive
\(\operatorname{Gr}^{\mathrm{W}}_i(\mathrm{M})\) for
\(i=0,-1,-2\) are \(\mathrm{X}\), \(\mathrm{A}\) and \(\mathrm{T}\).

\textbf{(10.1.5)} The construction (10.1.3) which associates to a 1-motive
\(\mathrm{M}\) over \(\mathbf{C}\)
\(\mathrm{T}_{\mathrm{Z}}(\mathrm{M})\) is transcendental.
We shall show that
\[
\widehat{\mathrm{T}}(\mathrm{M})
=\mathrm{T}_{\mathrm{Z}}(\mathrm{M})\otimes\widehat{\mathbf{Z}}
=\prod_{\ell}\mathrm{T}_{\mathrm{Z}}(\mathrm{M})\otimes\mathbf{Z}_{\ell}
\]
can be defined purely algebraically.

Let \(\mathrm{M}\) be a 1-motive over an algebraically closed field \(k\)
of characteristic 0. We identify \(\mathrm{M}\) with a complex
\([\mathrm{X}\xrightarrow{u}\mathrm{G}]\) (degrees 0 and 1).
For every integer \(n>0\), consider the complex
\([\mathbf{Z}\xrightarrow{n}\mathbf{Z}]\) (degrees \(-1\) and 0),
and let \(\mathrm{T}_{\mathbf{Z}/n\mathbf{Z}}(\mathrm{M})\) be the
\(\mathrm{H}^0\) of the complex
\([\mathrm{X}\xrightarrow{u}\mathrm{G}]
\otimes[\mathbf{Z}\xrightarrow{n}\mathbf{Z}]\):
\Page{57}
\[
\begin{tikzcd}[column sep=2em, row sep=2.5em]
\mathrm{X} \arrow[r, "u"] & \mathrm{G}
\\
\mathrm{X} \arrow[r, "u"] \arrow[u, "n"']
& \mathrm{G} \arrow[u, "-n"']
\end{tikzcd}
\]
We have
\[
\mathrm{T}_{\mathbf{Z}/n\mathbf{Z}}(\mathrm{M})
=\{(x,g)\mid u(x)=ng\}/\{(nx,u(x))\mid x\in\mathrm{X}\}.
\]
In the derived category, we have
\[
\mathrm{T}_{\mathbf{Z}/n\mathbf{Z}}(\mathrm{M})
=\mathrm{H}^0
\bigl(\mathrm{M}\overset{\mathrm{L}}{\otimes}\mathbf{Z}/n\mathbf{Z}\bigr).
\]
For \(n=md\), we define transition morphisms
\[
\varphi_{m,n}:
\mathrm{T}_{\mathbf{Z}/n\mathbf{Z}}(\mathrm{M})
\to\mathrm{T}_{\mathbf{Z}/m\mathbf{Z}}(\mathrm{M})
\qquad\text{by}\qquad
\varphi_{m,n}((x,g))=(x,dg).
\]
We set
\[
\widehat{\mathrm{T}}(\mathrm{M})
=\varprojlim\mathrm{T}_{\mathbf{Z}/n\mathbf{Z}}(\mathrm{M}).
\]
The filtration \(\mathrm{W}\) of \(\mathrm{M}\) induces a filtration
\(\mathrm{W}\) on the
\(\mathrm{T}_{\mathbf{Z}/n\mathbf{Z}}(\mathrm{M})\) and on
\(\widehat{\mathrm{T}}(\mathrm{M})\). We have
\[
\begin{aligned}
&\operatorname{Gr}^{\mathrm{W}}_0(\widehat{\mathrm{T}}(\mathrm{M}))
  =\mathrm{X}\otimes\widehat{\mathbf{Z}}
\\
&\operatorname{Gr}^{\mathrm{W}}_{-1}(\widehat{\mathrm{T}}(\mathrm{M}))
  =\varprojlim\mathrm{A}_n=\widehat{\mathrm{T}}(\mathrm{A})
\\
&\operatorname{Gr}^{\mathrm{W}}_{-2}(\widehat{\mathrm{T}}(\mathrm{M}))
  =\varprojlim\mathrm{T}_n=\mathrm{Y}\otimes\widehat{\mathbf{Z}}(1),
\end{aligned}
\]
where \(\mathrm{Y}\) is the dual of the character group of
\(\mathrm{T}=\mathrm{W}_{-2}(\mathrm{M})\).

\textit{Construction} \textbf{(10.1.6)}. —
\textit{Let \(\mathrm{M}=[\mathrm{X}\xrightarrow{u}\mathrm{G}]\)
be a 1-motive over \(\mathbf{C}\). We have}
\[
\tag*{\textbf{(10.1.6.1)}}
\widehat{\mathrm{T}}(\mathrm{M})
\simeq\mathrm{T}_{\mathrm{Z}}(\mathrm{M})\otimes\widehat{\mathbf{Z}}.
\]
The natural morphism
\([\mathrm{T}_{\mathrm{Z}}(\mathrm{M})\to\operatorname{Lie}(\mathrm{G})]
\to[\mathrm{X}\to\mathrm{G}]\) is a quasi-isomorphism.
The quasi-isomorphisms
\[
{
\left[
\begin{tikzcd}[ampersand replacement=\&, column sep=1.5em, row sep=2.6em]
\mathrm{X} \arrow[r, "u"] \& \mathrm{G}
\\
\mathrm{X} \arrow[r, "u"] \arrow[u, "n"']
\& \mathrm{G} \arrow[u, "-n"']
\end{tikzcd}
\right]
}
\longleftarrow
{
\left[
\begin{tikzcd}[ampersand replacement=\&, column sep=1.5em, row sep=2.6em]
\mathrm{T}_{\mathrm{Z}} \arrow[r, "u"]
\& \operatorname{Lie}(\mathrm{G})
\\
\mathrm{T}_{\mathrm{Z}} \arrow[r, "u"] \arrow[u, "n"']
\& \operatorname{Lie}(\mathrm{G}) \arrow[u, "\wr", "-n"']
\end{tikzcd}
\right]
}
\longrightarrow
{
\left[
\begin{tikzcd}[ampersand replacement=\&, column sep=1.5em, row sep=2.6em]
\mathrm{T}_{\mathrm{Z}} \arrow[r] \& 0
\\
\mathrm{T}_{\mathrm{Z}} \arrow[r, "c"] \arrow[u, "n"']
\& 0 \arrow[u]
\end{tikzcd}
\right]
}
\]
yield isomorphisms
\[
\tag*{\textbf{(10.1.6.2)}}
\mathrm{T}_{\mathbf{Z}/n\mathbf{Z}}(\mathrm{M})
\simeq\mathrm{T}_{\mathrm{Z}}/n\mathrm{T}_{\mathrm{Z}}.
\]
Under this isomorphism, \(t\in\mathrm{W}_{-1}(\mathrm{T}_{\mathrm{Z}})\)
is sent to \(\exp(t/n)\in\mathrm{G}_n\). Passing to the limit,
we deduce (10.1.6.1) from (10.1.6.2).

Techniques of Grothendieck allow us to carry over to de Rham
cohomology what we have just done in \(\ell\)-adic cohomology.
Let \(\mathrm{M}\) be a 1-motive over an algebraically closed field \(k\).

\Page{58}
\textit{Construction} \textbf{(10.1.7)}. —
\textit{Let us construct a vector space
\(\mathrm{T}_{\mathrm{DR}}(\mathrm{M})\) over \(k\), equipped with filtrations
\(\mathrm{W}\) and \(\mathrm{F}\).}

Let \(\mathrm{M}=(\mathrm{X},\mathrm{A},\mathrm{T},\mathrm{G},u)\),
and consider \(\mathrm{M}\) as a complex concentrated in degrees
\(-1\) and 0. An \textit{extension} of \(\mathrm{M}\) by
\(\mathbf{G}_a\) is an extension of \(\mathrm{M}\) by the complex
consisting only of \(\mathbf{G}_a\) placed in degree 0. We have
\(\operatorname{Hom}(\mathrm{G},\mathbf{G}_a)
=\operatorname{Ext}^1(\mathrm{X},\mathbf{G}_a)=0\).
Hence:

\noindent \textit{a)} extensions of \(\mathrm{M}\) by
\(\mathbf{G}_a\) have no automorphisms;\par
\noindent \textit{b)} the sequence
\[
0\to\operatorname{Hom}(\mathrm{X},\mathbf{G}_a)
\to\operatorname{Ext}^1(\mathrm{M},\mathbf{G}_a)
\to\operatorname{Ext}^1(\mathrm{G},\mathbf{G}_a)\to0
\]
is exact.

We have
\(\operatorname{Hom}(\mathrm{T},\mathbf{G}_a)
=\operatorname{Ext}^1(\mathrm{T},\mathbf{G}_a)=0\), so that
\[
\tag*{\textit{c)}}
\operatorname{Ext}^1(\mathrm{A},\mathbf{G}_a)
\xrightarrow{\sim}
\operatorname{Ext}^1(\mathrm{G},\mathbf{G}_a).
\]
The vector space \(\operatorname{Ext}^1(\mathrm{M},\mathbf{G}_a)\)
is therefore finite-dimensional and, by \textit{a)}, there thus exists a
universal extension of \(\mathrm{M}\) by a vector group;
we denote it by \(\mathrm{M}^{\natural}:\mathrm{X}\to\mathrm{G}^{\natural}\).
\[
\begin{tikzcd}[column sep=1.8em, row sep=3.5em]
&
0 \arrow[r] \arrow[d]
&
\mathrm{X} \arrow[r, equal] \arrow[d]
&
\mathrm{X} \arrow[d]
&
\\
0 \arrow[r]
&
\operatorname{Ext}^1(\mathrm{M},\mathbf{G}_a)^* \arrow[r]
&
\mathrm{G}^{\natural} \arrow[r]
&
\mathrm{G} \arrow[r]
&
0
\end{tikzcd}
\]
We set
\[
\begin{aligned}
\mathrm{T}_{\mathrm{DR}}(\mathrm{M})
  &=\operatorname{Lie}(\mathrm{G}^{\natural})
\\
\mathrm{F}^0\mathrm{T}_{\mathrm{DR}}(\mathrm{M})
  &=\operatorname{Ker}
    (\operatorname{Lie}\mathrm{G}^{\natural}\to\operatorname{Lie}\mathrm{G})
    \simeq(\operatorname{Ext}^1(\mathrm{M},\mathbf{G}_a))^*.
\end{aligned}
\]
\(\mathrm{T}_{\mathrm{DR}}\) is functorial in \(\mathrm{M}\)
(as is \(\mathrm{M}^{\natural}\)), and we define the filtration
\(\mathrm{W}\) on \(\mathrm{T}_{\mathrm{DR}}(\mathrm{M})\) to be the one
arising from the filtration \(\mathrm{W}\) on \(\mathrm{M}\).

\textit{Construction} \textbf{(10.1.8)}. —
\textit{Let \(\mathrm{M}\) be a 1-motive over \(\mathbf{C}\). We have}
\[
(\mathrm{T}_{\mathrm{DR}}(\mathrm{M}),\mathrm{F},\mathrm{W})
\simeq
(\mathrm{T}_{\mathbf{C}}(\mathrm{M}),\mathrm{F},\mathrm{W}).
\]
The \(\operatorname{Ext}^i(\mathrm{X},\mathbf{G}_a)\),
\(\operatorname{Ext}^i(\mathrm{A},\mathbf{G}_a)\) and
\(\operatorname{Ext}^i(\mathrm{T},\mathbf{G}_a)\) \((i=0,1)\)
are the same in the algebraic and analytic categories.
Consider the map
\(\mathrm{T}_{\mathbf{C}}(\mathrm{M})
=\mathrm{T}_{\mathrm{Z}}(\mathrm{M})\otimes\mathbf{C}
\to\operatorname{Lie}(\mathrm{G})\to\mathrm{G}\).
The diagram
\[
\begin{tikzcd}[column sep=2em, row sep=3em]
\mathrm{X} \arrow[r] \arrow[d, equal]
&
\mathrm{T}_{\mathbf{C}}(\mathrm{M})/\mathrm{H}_1(\mathrm{G})
  \arrow[d]
\\
\mathrm{X} \arrow[r]
&
\mathrm{G}
\end{tikzcd}
\]
\Page{59}
defines an extension of \(\mathrm{M}^{\mathrm{an}}\) by the vector
group \(\mathrm{F}^0\mathrm{T}_{\mathbf{C}}\mathrm{M}\), hence an
extension of \(\mathrm{M}\) by this group. We must prove that
this is the universal extension. For this, it suffices to check that
the category of extensions of
\([\mathrm{X}\to\mathrm{T}_{\mathbf{C}}(\mathrm{M})/\mathrm{H}_1(\mathrm{G})]\)
by \(\mathbf{G}_a\) is trivial (a single object, no automorphisms).
It is also the category of extensions of
\(\mathrm{T}_{\mathbf{C}}(\mathrm{M})/\mathrm{T}_{\mathrm{Z}}(\mathrm{M})\)
by \(\mathbf{G}_a\), and indeed we have
\[
\operatorname{Ext}^i
(\mathrm{T}_{\mathbf{C}}(\mathrm{M})/\mathrm{T}_{\mathrm{Z}}(\mathrm{M}),
\mathbf{G}_a)=0
\qquad\text{for}\qquad i=0,1.
\]

\textit{Remark} \textbf{(10.1.9)}. —
\(\mathrm{M}^{\natural}\) is characterized by the commutative diagram
\[
\begin{tikzcd}[column sep=2em, row sep=3.2em]
\mathrm{X} \arrow[r]
&
\mathrm{G}^{\natural} \arrow[r]
&
\mathrm{G}
\\
\mathrm{T}_{\mathrm{Z}}(\mathrm{M}) \arrow[r] \arrow[u]
&
\operatorname{Lie}(\mathrm{G}^{\natural}) \arrow[r] \arrow[u]
&
\operatorname{Lie}(\mathrm{G}) \arrow[u]
\end{tikzcd}
\]
whose outer square is (10.1.3.1), and which induces an isomorphism
\(\mathrm{T}_{\mathbf{C}}(\mathrm{M})
\xrightarrow{\sim}\operatorname{Lie}(\mathrm{G}^{\natural})\).

\textit{Variant} \textbf{(10.1.10)}. —
A \textit{smooth 1-motive} \(\mathrm{M}\) over a scheme \(\mathrm{S}\)
consists of:

\noindent \textit{a)} a group scheme \(\mathrm{X}\) over
\(\mathrm{S}\), which locally for the étale topology is a constant
group scheme defined by a finitely generated free \(\mathbf{Z}\)-module;
an abelian scheme \(\mathrm{A}\) over \(\mathrm{S}\), and a torus
\(\mathrm{T}\) over \(\mathrm{S}\);\par
\noindent \textit{b)} an extension \(\mathrm{G}\) of \(\mathrm{A}\)
by \(\mathrm{T}\);\par
\noindent \textit{c)} a morphism \(u:\mathrm{X}\to\mathrm{G}\).

For every integer \(n\), a construction analogous to (10.1.5) associates
to a smooth 1-motive \(\mathrm{M}\) over \(\mathrm{S}\) a group
scheme \(\mathrm{T}_{\mathbf{Z}/n\mathbf{Z}}(\mathrm{M})\) finite and
flat over \(\mathrm{S}\). For \(\ell\) invertible on \(\mathrm{S}\),
the inverse system
\[
\mathrm{T}_{\ell}(\mathrm{M})
=\varprojlim\mathrm{T}_{\mathbf{Z}/\ell^n\mathbf{Z}}(\mathrm{M})
\]
is a \(\mathbf{Z}_{\ell}\)-sheaf on \(\mathrm{S}\). In general,
\(\mathrm{T}_{\ell}(\mathrm{M})\) corresponds to an
\(\ell\)-divisible (\(=\) Barsotti-Tate) group over \(\mathrm{S}\).

The filtration \(\mathrm{W}\) of \(\mathrm{M}\), defined as in
(10.1.4), defines a filtration \(\mathrm{W}\) on
\(\mathrm{T}_{\ell}(\mathrm{M})\). We again have
\[
\begin{aligned}
&\operatorname{Gr}^{\mathrm{W}}_0\mathrm{T}_{\ell}(\mathrm{M})
  =\mathrm{X}\otimes\widehat{\mathbf{Z}}
\\
&\operatorname{Gr}^{\mathrm{W}}_{-1}\mathrm{T}_{\ell}(\mathrm{M})
  =\mathrm{T}_{\ell}(\mathrm{A})
\\
&\operatorname{Gr}^{\mathrm{W}}_{-2}\mathrm{T}_{\ell}(\mathrm{M})
  =\mathrm{Y}\otimes\mathbf{Z}_{\ell}(1)
\end{aligned}
\]
where \(\mathrm{Y}=\mathit{Hom}(\mathrm{T},\mathbf{G}_m)^{\vee}\).

Similarly, a construction analogous to (10.1.7) associates to \(\mathrm{M}\)
a vector bundle \(\mathrm{T}_{\mathrm{DR}}(\mathrm{M})\).

\textit{Terminology} \textbf{(10.1.11)}. —
Let \(\mathrm{M}\) be a 1-motive. We call \(\mathrm{T}(\mathrm{M})\),
\(\mathrm{T}_{\ell}(\mathrm{M})\) and
\(\mathrm{T}_{\mathrm{DR}}(\mathrm{M})\) the
\textit{Hodge, \(\ell\)-adic} and
\textit{de Rham realizations} of the 1-motive \(\mathrm{M}\).

\Page{60}
\section*{10.2. 1-motives and biextensions}

The results of this subsection will not be used in the remainder of this
section. We shall use the notion of biextension
(S.G.A. 7, VII (2.1)) and the following generalization.

\textbf{(10.2.1)} Let
\(\mathrm{K}_i:[\mathrm{A}_i\to\mathrm{B}_i]\) \((i=1,2)\)
be two complexes of sheaves, concentrated in degrees 0 and \(-1\)
(on an arbitrary topos). A \textit{biextension} of
\(\mathrm{K}_1\) and \(\mathrm{K}_2\) by an abelian sheaf
\(\mathrm{H}\) consists of:

\noindent \textit{a)} a biextension \(\mathrm{P}\) of
\(\mathrm{B}_1\) and \(\mathrm{B}_2\) by \(\mathrm{H}\), i.e. an
\(\mathrm{H}\)-principal homogeneous space \(\mathrm{P}\) over
\(\mathrm{B}_1\times\mathrm{B}_2\), with
\(\mathrm{P}_{b_1b_2}\) depending biadditively on \(b_1\) and \(b_2\);\par
\noindent \textit{b)} a trivialization (\(=\) biadditive section)
of the biextension of \(\mathrm{B}_1\) and \(\mathrm{A}_2\) by
\(\mathrm{H}\), the pullback of \(\mathrm{P}\);\par
\noindent \textit{c)} a trivialization of the biextension of
\(\mathrm{A}_1\) and \(\mathrm{B}_2\) by \(\mathrm{H}\), the
pullback of \(\mathrm{P}\); we require that the trivializations
\textit{b)} and \textit{c)} coincide on
\(\mathrm{A}_1\times\mathrm{A}_2\).

We denote by
\(\symbfit{Biext}(\mathrm{K}_1,\mathrm{K}_2;\mathrm{H})\)
the category of biextensions of \(\mathrm{K}_1\) and \(\mathrm{K}_2\)
by \(\mathrm{H}\). It is a Picard category (the objects can be
added). The group
\(\operatorname{Biext}^0(\mathrm{K}_1,\mathrm{K}_2;\mathrm{H})\)
of automorphisms of any biextension of \(\mathrm{K}_1\)
and \(\mathrm{K}_2\) by \(\mathrm{H}\) is
\[
\operatorname{Biext}^0(\mathrm{K}_1,\mathrm{K}_2;\mathrm{H})
=\operatorname{Hom}
(\mathrm{H}^0(\mathrm{K}_1)\otimes\mathrm{H}^0(\mathrm{K}_2),\mathrm{H}).
\]
We denote by
\(\operatorname{Biext}^1(\mathrm{K}_1,\mathrm{K}_2;\mathrm{H})\)
the group of isomorphism classes of biextensions. One checks as
in (S.G.A. 7, VII (3.6.5)) that
\[
\operatorname{Biext}^i(\mathrm{K}_1,\mathrm{K}_2;\mathrm{H})
=\operatorname{Ext}^i
(\mathrm{K}_1\overset{\mathrm{L}}{\otimes}\mathrm{K}_2;\mathrm{H})
\qquad(i=0,1).
\]
Using this formula, or by elementary means, one checks that if
\(\mathrm{K}'_i\to\mathrm{K}_i\) is a quasi-isomorphism
\((i=1,2)\), giving a biextension of
\(\mathrm{K}_1\) and \(\mathrm{K}_2\) by \(\mathrm{H}\) is equivalent to giving one of
\(\mathrm{K}'_1\) and \(\mathrm{K}'_2\) by \(\mathrm{H}\).

This still applies on replacing “sheaves” by “algebraic
groups” (resp. “complex analytic groups”): these can
be interpreted as sheaves on a suitable site.

\textbf{(10.2.2)} In this section, we shall identify a 1-motive
\(\mathrm{M}=(\mathrm{X},\mathrm{T},\mathrm{A},\mathrm{G},u)\)
with a complex concentrated in \textit{degrees \(0\) and \(-1\)}
\[
\mathrm{M}:[\mathrm{X}\to\mathrm{G}].
\]
Over \(\mathbf{C}\), we shall denote by \(\mathrm{M}^{\mathrm{an}}\)
the complex of complex analytic groups
\([\mathrm{X}\to\mathrm{G}^{\mathrm{an}}]\).
We have the following rigidity result:

\textit{Lemma} \textbf{(10.2.2.1)}. —
\textit{Let
\(\mathrm{M}_i=(\mathrm{X}_i,\mathrm{A}_i,\mathrm{T}_i,\mathrm{G}_i,u_i)\)
\((i=1,2)\) be two 1-motives. We have}
\[
\operatorname{Biext}^0(\mathrm{M}_1,\mathrm{M}_2;\mathbf{G}_m)=0.
\]
\textit{Over \(\mathbf{C}\), we likewise have
\(\operatorname{Biext}^0
(\mathrm{M}_1^{\mathrm{an}},\mathrm{M}_2^{\mathrm{an}};\mathbf{G}_m)=0\).}

Every biadditive morphism
\(\mathrm{G}_1\times\mathrm{G}_2\to\mathbf{G}_m\)
is indeed trivial.

% ---- 061-068 ----
\Page{61}
\textit{Construction} \textbf{(10.2.3)}. — \textit{Let
\(\mathrm{M}_i=(\mathrm{X}_i,\mathrm{A}_i,\mathrm{T}_i,\mathrm{G}_i,u_i)\)
\((i=1,2)\) be two \(1\)-motives over \(\mathbf{C}\). We have}
\[
\operatorname{Biext}^{1}(\mathrm{M}_1,\mathrm{M}_2;\mathbf{G}_m)
=\operatorname{Hom}(\mathrm{T}(\mathrm{M}_1)\otimes
\mathrm{T}(\mathrm{M}_2),\mathbf{Z}(1)).
\]

We shall first compute
\(\operatorname{Biext}^{1}(\mathrm{M}_1^{\mathrm{an}},
\mathrm{M}_2^{\mathrm{an}};\mathbf{G}_m)\).

If \(\mathrm{V}\) and \(\mathrm{W}\) are two vector spaces, we have
\textit{in the analytic category}
\[
\operatorname{Biext}^{0}(\mathrm{V},\mathrm{W};\mathbf{G}_m)
\xleftarrow[\exp]{\sim}
\operatorname{Hom}(\mathrm{V}\otimes\mathrm{W},\mathbf{C})
\tag*{\textbf{(10.2.3.1)}\textsubscript{an}}
\]
\[
\operatorname{Biext}^{1}(\mathrm{V},\mathrm{W};\mathbf{G}_m)=0.
\tag*{\textbf{(10.2.3.2)}\textsubscript{an}}
\]
Indeed, it follows formally from the fact that an extension of
\(\mathrm{W}\) by \(\mathbf{G}_m\), over any base, is always
locally trivial that
\[
\begin{aligned}
\operatorname{Biext}^{i}(\mathrm{V},\mathrm{W};\mathbf{G}_m)
&=\operatorname{Ext}^{i}(\mathrm{V},
  \mathit{Hom}(\mathrm{W},\mathbf{G}_m))\\
&=\operatorname{Ext}^{i}(\mathrm{V},\mathrm{W}^{*})
\qquad(i=1,2).
\end{aligned}
\]

Let \(\mathrm{V}_{\mathbf{Z}}\) and \(\mathrm{W}_{\mathbf{Z}}\) be two
free \(\mathbf{Z}\)-modules of finite rank. Since
\[
[\mathrm{V}_{\mathbf{Z}}\to\mathrm{V}_{\mathbf{C}}]
\to[0\to\mathrm{V}_{\mathbf{C}}/\mathrm{V}_{\mathbf{Z}}]
\]
is a quasi-isomorphism, the “pullback” functor is an
equivalence.
\[
\symbfit{Biext}(\mathrm{V}_{\mathbf{C}}/\mathrm{V}_{\mathbf{Z}},
\mathrm{W}_{\mathbf{C}}/\mathrm{W}_{\mathbf{Z}};\mathbf{G}_m)
\xrightarrow{\approx}
\symbfit{Biext}([\mathrm{V}_{\mathbf{Z}}\to\mathrm{V}_{\mathbf{C}}],
[\mathrm{W}_{\mathbf{Z}}\to\mathrm{W}_{\mathbf{C}}];\mathbf{G}_m).
\tag*{\textbf{(10.2.3.3)}\textsubscript{an}}
\]
Every biadditive morphism from
\(\mathrm{V}_{\mathbf{C}}/\mathrm{V}_{\mathbf{Z}}\times
\mathrm{W}_{\mathbf{C}}/\mathrm{W}_{\mathbf{Z}}\)
to \(\mathbf{G}_m\) is trivial, whence \textit{a)} below:

\textit{Lemma} \textbf{(10.2.3.4)}
\[
\begin{aligned}
\text{\textit{a)}}\quad&
\operatorname{Biext}^{0}(\mathrm{V}_{\mathbf{C}}/\mathrm{V}_{\mathbf{Z}},
\mathrm{W}_{\mathbf{C}}/\mathrm{W}_{\mathbf{Z}};\mathbf{G}_m)=0.\\
\text{\textit{b)}}\quad&
\operatorname{Biext}^{1}(\mathrm{V}_{\mathbf{C}}/\mathrm{V}_{\mathbf{Z}},
\mathrm{W}_{\mathbf{C}}/\mathrm{W}_{\mathbf{Z}};\mathbf{G}_m)
=\operatorname{Hom}(\mathrm{V}_{\mathbf{Z}}\otimes
\mathrm{W}_{\mathbf{Z}},\mathbf{Z}(1)).
\end{aligned}
\]

Let \(\psi_i:\mathrm{V}_{\mathbf{C}}\otimes
\mathrm{W}_{\mathbf{C}}\to\mathbf{C}\) be two bilinear maps,
with \(\psi=\psi_1-\psi_2\) taking values in
\(\mathbf{Z}(1)=2\pi i\mathbf{Z}\) on
\(\mathrm{V}_{\mathbf{Z}}\otimes\mathrm{W}_{\mathbf{Z}}\).
Let \(\mathrm{P}(\psi_1,\psi_2)\) be the following biextension of
\([\mathrm{V}_{\mathbf{Z}}\to\mathrm{V}_{\mathbf{C}}]\) and
\([\mathrm{W}_{\mathbf{Z}}\to\mathrm{W}_{\mathbf{C}}]\) by
\(\mathbf{G}_m\): the trivial biextension of
\(\mathrm{V}_{\mathbf{C}}\) and \(\mathrm{W}_{\mathbf{C}}\) by
\(\mathbf{G}_m\), equipped with the trivializations
\(\varphi_1=\exp\psi_1:\mathrm{V}_{\mathbf{Z}}\otimes
\mathrm{W}_{\mathbf{C}}\to\mathbf{G}_m\) and
\(\varphi_2=\exp\psi_2:\mathrm{V}_{\mathbf{C}}\otimes
\mathrm{W}_{\mathbf{Z}}\to\mathbf{G}_m\).
By (10.2.3.2), every biextension is of this form; by
(10.2.3.1)
\(\mathrm{P}(\psi_1,\psi_2)\simeq\mathrm{P}(\psi_1',\psi_2')\)
if and only if \(\psi_1'-\psi_2'=\psi_1-\psi_2\).
The \(\mathrm{P}(\psi,0)\) therefore represent each biextension
once and only once.

\textbf{(10.2.3.5)} Under the general hypotheses of (10.2.1), let
\(\mathrm{K}_i=[\mathrm{A}_i+\mathrm{F}_i\to\mathrm{B}_i]\)
\((i=1,2)\) be complexes. Let \(\mathrm{K}_i'\) be the subcomplex
\([\mathrm{A}_i\to\mathrm{B}_i]\). One checks that giving
a biextension of \(\mathrm{K}_1\) and
\(\mathrm{K}_2\) by \(\mathrm{H}\) is equivalent to giving:

\noindent\textit{a)} a biextension \(\mathrm{P}\) of
\(\mathrm{K}_1'\) and \(\mathrm{K}_2'\) by \(\mathrm{H}\); and

\noindent\textit{b)} trivializations of the biextensions of
\([0\to\mathrm{F}_1]\) and \(\mathrm{K}_2'\)
(resp. \(\mathrm{K}_1'\) and \([0\to\mathrm{F}_2]\)) by
\(\mathrm{H}\) obtained by pulling back \(\mathrm{P}\); these
trivializations must agree on
\(\mathrm{F}_1\times\mathrm{F}_2\).

In the particular case where
\begin{align*}
&\operatorname{Biext}^{i}(\mathrm{F}_1,\mathrm{K}_2';\mathrm{H})=0
\qquad(i=1,2)\\
\text{and}\qquad&
\operatorname{Biext}^{i}(\mathrm{K}_1',\mathrm{F}_2;\mathrm{H})=0
\qquad(i=1,2),
\end{align*}
\Page{62}
there is one and only one trivialization \(\varphi_1\), a section of
\(\mathrm{P}\) over \(\mathrm{F}_1\times\mathrm{B}_2\)
(resp. \(\varphi_2\), a section of \(\mathrm{P}\) over
\(\mathrm{B}_1\times\mathrm{F}_2\)) as in \textit{b)}.
The category
\(\symbfit{Biext}(\mathrm{K}_1,\mathrm{K}_2;\mathrm{H})\)
is then identified with the subcategory of
\(\symbfit{Biext}(\mathrm{K}_1',\mathrm{K}_2';\mathrm{H})\)
consisting of those \(\mathrm{P}\) for which \(\varphi_1=\varphi_2\)
on \(\mathrm{F}_1\times\mathrm{F}_2\).

\textbf{(10.2.3.6)} Let \(\mathrm{V}_{\mathbf{Z}}\) and
\(\mathrm{W}_{\mathbf{Z}}\) be two finitely generated free
\(\mathbf{Z}\)-modules, and let \(\mathrm{F},\mathrm{G}\) be vector subspaces
of \(\mathrm{V}_{\mathbf{C}}\) and \(\mathrm{W}_{\mathbf{C}}\). We have
\[
\operatorname{Biext}^{i}(\mathrm{F},
[\mathrm{W}_{\mathbf{Z}}\to\mathrm{W}_{\mathbf{C}}];\mathbf{G}_m)=0.
\]
For \(i=0\), this follows from (10.2.3.1) and from the fact that a
bilinear map
\[
\mathrm{B}:\mathrm{F}\times\mathrm{W}_{\mathbf{C}}\to\mathbf{C}
\qquad\text{with}\qquad
\mathrm{B}(\mathrm{F},\mathrm{W}_{\mathbf{Z}})\subset\mathbf{Z}(1)
\]
is zero. For \(i=1\), this follows from the fact that every biadditive
morphism
\(\mathrm{F}\times\mathrm{W}_{\mathbf{Z}}\to\mathbf{G}_m\)
comes from a biadditive morphism
\(\mathrm{F}\times\mathrm{W}_{\mathbf{C}}\to\mathbf{G}_m\).
In other words, we use the long exact sequence
\[
\begin{aligned}
0&\to\operatorname{Biext}^{0}(\mathrm{F},
[\mathrm{W}_{\mathbf{Z}}\to\mathrm{W}_{\mathbf{C}}];\mathbf{G}_m)
\to\operatorname{Biext}^{0}(\mathrm{F},\mathrm{W}_{\mathbf{C}};\mathbf{G}_m)
\xrightarrow{\sim}
\operatorname{Biext}^{0}(\mathrm{F},\mathrm{W}_{\mathbf{Z}};\mathbf{G}_m)\\
&\to\operatorname{Biext}^{1}(\mathrm{F},
[\mathrm{W}_{\mathbf{Z}}\to\mathrm{W}_{\mathbf{C}}];\mathbf{G}_m)
\to\operatorname{Biext}^{1}(\mathrm{F},\mathrm{W}_{\mathbf{C}};\mathbf{G}_m).
\end{aligned}
\]

We can therefore apply (10.2.3.5) to the complexes
\(\mathrm{K}_1=[\mathrm{V}_{\mathbf{Z}}\oplus\mathrm{F}
\to\mathrm{V}_{\mathbf{C}}]\) and
\(\mathrm{K}_2=[\mathrm{W}_{\mathbf{Z}}\oplus\mathrm{G}
\to\mathrm{W}_{\mathbf{C}}]\).
We find that the biextensions of \(\mathrm{K}_1\) and
\(\mathrm{K}_2\) by \(\mathbf{G}_m\) are identified with certain
biextensions of
\(\mathrm{V}_{\mathbf{C}}/\mathrm{V}_{\mathbf{Z}}\) and
\(\mathrm{W}_{\mathbf{C}}/\mathrm{W}_{\mathbf{Z}}\) by
\(\mathbf{G}_m\), i.e. with certain forms
\(\psi:\mathrm{V}_{\mathbf{Z}}\otimes\mathrm{W}_{\mathbf{Z}}
\to\mathbf{Z}(1)\) (10.2.3.4).
For \(\mathrm{P}(\psi_1,\psi_2)\) (10.2.3.4) to come from a
biextension of \(\mathrm{K}_1\) and \(\mathrm{K}_2\) by
\(\mathbf{G}_m\), it is necessary and sufficient that
\(\exp\psi_2:\mathrm{F}\times\mathrm{W}_{\mathbf{C}}\to\mathbf{G}_m\)
and
\(\exp\psi_1:\mathrm{V}_{\mathbf{C}}\times\mathrm{G}\to\mathbf{G}_m\)
agree on \(\mathrm{F}\times\mathrm{G}\), that is, that
\(\psi(\mathrm{F},\mathrm{G})=0\).

\textit{Lemma} \textbf{(10.2.3.7)}. — \textit{With the preceding
notation, we have
\(\operatorname{Biext}^{0}(\mathrm{K}_1,\mathrm{K}_2;\mathbf{G}_m)=0\)
and
\(\operatorname{Biext}^{1}(\mathrm{K}_1,\mathrm{K}_2;\mathbf{G}_m)\)
is identified with the set of forms
\(\psi:\mathrm{V}_{\mathbf{Z}}\otimes\mathrm{W}_{\mathbf{Z}}
\to\mathbf{Z}(1)\) whose complexifications satisfy
\(\psi(\mathrm{F},\mathrm{G})=0\).}

For \(\psi\) as in (10.2.3.7), with \(\psi=\psi_1-\psi_2\),
the biextension corresponding to \(\psi\) is given by the
following trivializations of the trivial biextension of
\(\mathrm{V}_{\mathbf{C}}\) and \(\mathrm{W}_{\mathbf{C}}\) by
\(\mathbf{G}_m\):
\[
\begin{aligned}
&\text{on }(\mathrm{V}_{\mathbf{Z}}\oplus\mathrm{F})
\times\mathrm{W}_{\mathbf{C}}
&&:\quad \exp\psi_1(v_{\mathbf{Z}},w_{\mathbf{C}})
.\exp\psi_2(f,w_{\mathbf{C}})\\
&\text{on }\mathrm{V}_{\mathbf{C}}\times
(\mathrm{W}_{\mathbf{Z}}\oplus\mathrm{G})
&&:\quad \exp\psi_2(v_{\mathbf{C}},w_{\mathbf{Z}})
.\exp\psi_1(v_{\mathbf{C}},g).
\end{aligned}
\]

\textbf{(10.2.3.8)} For the 1-motives \(\mathrm{M}_i\), we have
quasi-isomorphisms
\[
\begin{tikzcd}[row sep=3.5em, column sep=2em]
\mathrm{T}_{\mathbf{Z}}(\mathrm{M}_i)\oplus
\mathrm{F}^{0}\mathrm{T}_{\mathbf{C}}(\mathrm{M}_i)
  \arrow[r] \arrow[d]
&
\mathrm{T}_{\mathbf{Z}}(\mathrm{M}_i)
  \arrow[r] \arrow[d]
&
\mathrm{X}_i \arrow[d]
\\
\mathrm{T}_{\mathbf{C}}(\mathrm{M}_i) \arrow[r]
&
\operatorname{Lie}(\mathrm{G}_i) \arrow[r]
&
\mathrm{G}_i
\end{tikzcd}
\]
\Page{63}
The group
\(\operatorname{Biext}^{1}(\mathrm{M}_1^{\mathrm{an}},
\mathrm{M}_2^{\mathrm{an}};\mathbf{G}_m)\)
is therefore identified with the group of maps
\[
\psi:\mathrm{T}_{\mathbf{Z}}(\mathrm{M}_1)\otimes
\mathrm{T}_{\mathbf{Z}}(\mathrm{M}_2)\to\mathbf{Z}(1)
\]
whose complexifications are compatible with the Hodge filtration.

\textit{Lemma} \textbf{(10.2.3.9)}. — \textit{The map
\(\operatorname{Biext}^{1}(\mathrm{M}_1,\mathrm{M}_2;\mathbf{G}_m)
\to\operatorname{Biext}^{1}(\mathrm{M}_1^{\mathrm{an}},
\mathrm{M}_2^{\mathrm{an}};\mathbf{G}_m)\)
is injective. Its image consists of the biextensions \(\mathrm{P}\)
whose restriction to \(\mathrm{G}_1\times\mathrm{G}_2\) is the
pullback of a biextension of \(\mathrm{A}_1\) and
\(\mathrm{A}_2\) by \(\mathbf{G}_m\).}

We have
\begin{align*}
&\operatorname{Biext}^{1}(\mathrm{A}_1,\mathrm{A}_2;\mathbf{G}_m)
\xrightarrow{\sim}
\operatorname{Biext}^{1}(\mathrm{G}_1,\mathrm{G}_2;\mathbf{G}_m)
\qquad\text{(S.G.A. 7, VIII (3.5))}\\
\text{and}\qquad&
\operatorname{Biext}^{1}(\mathrm{A}_1,\mathrm{A}_2;\mathbf{G}_m)
\xrightarrow{\sim}
\operatorname{Biext}^{1}(\mathrm{A}_1^{\mathrm{an}},
\mathrm{A}_2^{\mathrm{an}};\mathbf{G}_m)
\qquad\text{(G.A.G.A.).}
\end{align*}

Let \(\mathrm{P}\) be a biextension of \(\mathrm{M}_1\) and
\(\mathrm{M}_2\) by \(\mathbf{G}_m\). If
\(\mathrm{P}^{\mathrm{an}}\) is trivial, then:

\textit{a)} the restriction of \(\mathrm{P}\) to
\(\mathrm{G}_1\times\mathrm{G}_2\) is the image of a biextension
\(\mathrm{P}_0\) of \(\mathrm{A}_1\) and \(\mathrm{A}_2\) by
\(\mathbf{G}_m\); the latter is trivial. Otherwise it would be
analytically nontrivial, would give rise to a nonzero bilinear
form \(\psi_0\), and the form \(\psi\) corresponding to
\(\mathrm{P}\) would be nonzero \textit{a fortiori}. The restriction
of \(\mathrm{P}\) to \(\mathrm{G}_1\times\mathrm{G}_2\) is therefore
trivial (and uniquely trivializable, even analytically).

\textit{b)} \(\mathrm{P}\) is therefore defined by biadditive
maps \(\mathrm{X}_1\times\mathrm{G}_2\to\mathbf{G}_m\) and
\(\mathrm{G}_1\times\mathrm{X}_2\to\mathbf{G}_m\). These are
zero in the analytic category, hence are zero, and
\(\mathrm{P}\) is trivial.

Let \(\mathrm{P}\) be a biextension of
\(\mathrm{M}_1^{\mathrm{an}}\) and \(\mathrm{M}_2^{\mathrm{an}}\)
by \(\mathbf{G}_m\) satisfying condition (10.2.3.9).
It remains to prove that \(\mathrm{P}\) is algebraizable. By
hypothesis and G.A.G.A., its restriction \(\mathrm{P}_1\) to
\(\mathrm{G}_1\times\mathrm{G}_2\) is algebraizable. We conclude by showing
that a trivialization of \(\mathrm{P}_1\) on
\(\mathrm{X}_1\times\mathrm{G}_2\)
(resp. \(\mathrm{G}_1\times\mathrm{X}_2\)) is automatically
algebraic: \(\mathrm{P}_1\) can be interpreted as an extension of
\(\mathrm{X}_1\otimes\mathrm{G}_2\) by \(\mathbf{G}_m\), which
reduces matters to showing that
\[
\operatorname{Ext}^{i}(\mathrm{G}_1,\mathbf{G}_m)
\xrightarrow{\sim}
\operatorname{Ext}^{i}(\mathrm{G}_1^{\mathrm{an}},\mathbf{G}_m)
\qquad(i=0,1).
\]
This follows from the same assertion for an abelian variety
(G.A.G.A.) and for \(\mathbf{G}_m\).

By (10.2.3.9),
\(\operatorname{Biext}^{1}(\mathrm{M}_1,\mathrm{M}_2;\mathbf{G}_m)\)
is identified with the set of
\[
\psi:\mathrm{T}_{\mathbf{Z}}(\mathrm{M}_1)\otimes
\mathrm{T}_{\mathbf{Z}}(\mathrm{M}_2)\to\mathbf{Z}(1)
\]
such that:

\noindent\textit{a)} \(\psi\) is compatible with the Hodge
filtration (10.2.3.8);

\noindent\textit{b)} the restriction of \(\psi\) to
\(\mathrm{W}_{-1}\mathrm{T}_{\mathbf{Z}}(\mathrm{M}_1)\otimes
\mathrm{W}_{-1}\mathrm{T}_{\mathbf{Z}}(\mathrm{M}_2)\)
comes from a form on
\[
\operatorname{Gr}^{\mathrm{W}}_{-1}
\mathrm{T}_{\mathbf{Z}}(\mathrm{M}_1)\otimes
\operatorname{Gr}^{\mathrm{W}}_{-1}
\mathrm{T}_{\mathbf{Z}}(\mathrm{M}_2).
\]
Condition \textit{b)} means that \(\psi\) is compatible with
\(\mathrm{W}\), and this completes construction (10.2.3).

\textit{Remark} \textbf{(10.2.4)}. — A biextension
\(\mathrm{P}\) of \(\mathrm{M}_1\) and \(\mathrm{M}_2\) by
\(\mathbf{G}_m\) defines in an obvious way a biextension
\(\mathrm{P}^{0}\) of \(\mathrm{M}_2\) and \(\mathrm{M}_1\) by
\(\mathbf{G}_m\). If \(\psi\) (resp. \(\psi^{0}\)) is the bilinear
form
\Page{64}
defined by \(\mathrm{P}\) (resp. \(\mathrm{P}^{0}\)), we have
\(\psi(x,y)=-\psi^{0}(y,x)\). This is clear from the formulas in
(10.2.3.7).

Construction (10.2.3) has purely algebraic analogues in
\(\ell\)-adic or de Rham cohomology.

\textit{Construction} \textbf{(10.2.5)}. — \textit{Let
\(\mathrm{M}_i=(\mathrm{X}_i,\mathrm{A}_i,\mathrm{T}_i,\mathrm{G}_i,u_i)\)
\((i=1,2)\) be two \(1\)-motives over an algebraically closed
field \(k\) of characteristic \(0\). Let \(\mathrm{P}\) be a biextension
of \(\mathrm{M}_1\) and \(\mathrm{M}_2\) by \(\mathbf{G}_m\).
We associate to it a morphism}
\[
\mathrm{T}_{\mathbf{Z}/n\mathbf{Z}}(\mathrm{M}_1)\otimes
\mathrm{T}_{\mathbf{Z}/n\mathbf{Z}}(\mathrm{M}_2)
\to\mathbf{Z}/n\mathbf{Z}(1).
\]

Let \(m_i\in\mathrm{T}_{\mathbf{Z}/n\mathbf{Z}}(\mathrm{M}_1)\)
be represented by \((x_i,g_i)\) with \(u_i(x_i)=ng_i\).
Let us construct two isomorphisms \(a_1\) and \(a_2\) from
\(\mathrm{P}_{g_1g_2}^{\otimes n}\) to the trivial principal homogeneous space
\(\mathbf{G}_m\):
\[
\begin{aligned}
a_1:\quad
\mathrm{P}_{g_1,g_2}^{\otimes n}
&\xrightarrow{\sim}\mathrm{P}_{ng_1,g_2}
=\mathrm{P}_{u(x_1),g_2}\xrightarrow{\sim}\mathbf{G}_m\\
a_2:\quad
\mathrm{P}_{g_1,g_2}^{\otimes n}
&\xrightarrow{\sim}\mathrm{P}_{g_1,ng_2}
=\mathrm{P}_{g_1,u(x_2)}\xrightarrow{\sim}\mathbf{G}_m.
\end{aligned}
\]
Set \(a_2=\Phi(m_1,m_2)a_1\). One checks that
\(\Phi(m_1,m_2)\) is independent of the choice of the \((x_i,g_i)\) and
is an \(n\)-th root of unity. This is the required form.

\textit{Proposition} \textbf{(10.2.6)}. — \textit{Over
\(\mathbf{C}\), the form \((10.2.5)\) is obtained from that constructed
in \((10.2.3)\) by reduction modulo \(n\).}

This is verified using the formulas given in (10.2.3.7).

\textit{Construction} \textbf{(10.2.7)}. — \textit{Let
\(\mathrm{M}_i=(\mathrm{X}_i,\mathrm{A}_i,\mathrm{T}_i,\mathrm{G}_i,u_i)\)
be two \(1\)-motives over an algebraically closed field \(k\).
Let \(\mathrm{P}\) be a biextension of \(\mathrm{M}_1\) and
\(\mathrm{M}_2\) by \(\mathbf{G}_m\).
We associate to it a morphism}
\[
\mathrm{T}_{\mathrm{DR}}(\mathrm{M}_1)\otimes_k
\mathrm{T}_{\mathrm{DR}}(\mathrm{M}_2)\to k.
\]

Let us recall the definitions (due to Grothendieck) of
\(\natural\)-extensions and biextensions.

\textbf{(10.2.7.1)} A \(\natural\)-\textit{extension} of a smooth
commutative group \(\mathrm{G}\) by a smooth commutative group
\(\mathrm{H}\) consists of:

\noindent\textit{a)} an extension \(\mathrm{E}\) of \(\mathrm{G}\)
by \(\mathrm{H}\); if
\(\mu:\mathrm{G}\times\mathrm{G}\to\mathrm{G}\) is addition,
we regard \(\mathrm{E}\) as a principal homogeneous \(\mathrm{H}\)-space over
\(\mathrm{G}\) equipped with a morphism of principal homogeneous \(\mathrm{H}\)-spaces over
\(\mathrm{G}\times\mathrm{G}\) (addition)
\[
\nu:\operatorname{pr}_1^{*}\mathrm{E}
+\operatorname{pr}_2^{*}\mathrm{E}\to\mu^{*}\mathrm{E};
\]

\noindent\textit{b)} a connection on the principal homogeneous \(\mathrm{H}\)-space
\(\mathrm{E}\), such that the map \(\nu\) is horizontal.
This connection automatically has zero curvature.

If \(\mathrm{G}\) is an extension of an abelian variety by a
torus, every extension of \(\mathrm{G}\) by \(\mathbf{G}_m\)
admits a \(\natural\)-structure; two \(\natural\)-structures
differ by an invariant form on \(\mathrm{G}\) (with values
in the Lie algebra of \(\mathbf{G}_m\)). The Baer sum of two
\(\natural\)-extensions is defined in the obvious way.

\Page{65}
\textbf{(10.2.7.2)} Let \(\mathrm{P}\) be a biextension of
\(\mathrm{G}_1\) and \(\mathrm{G}_2\) by \(\mathrm{H}\).
Regard \(\mathrm{P}\) as a principal homogeneous \(\mathrm{H}\)-space
\(\mathrm{P}\) over \(\mathrm{G}_1\times\mathrm{G}_2\), equipped with
morphisms
\[
\begin{aligned}
\nu_1:\quad
\operatorname{pr}_{13}^{*}\mathrm{P}
+\operatorname{pr}_{23}^{*}\mathrm{P}
&\to(\mu_1\times\mathrm{Id})^{*}\mathrm{P}
&&\quad\text{over}\quad
\mathrm{G}_1\times\mathrm{G}_1\times\mathrm{G}_2\\
\nu_2:\quad
\operatorname{pr}_{12}^{*}\mathrm{P}
+\operatorname{pr}_{13}^{*}\mathrm{P}
&\to(\mathrm{Id}\times\mu_2)^{*}\mathrm{P}
&&\quad\text{over}\quad
\mathrm{G}_1\times\mathrm{G}_2\times\mathrm{G}_2.
\end{aligned}
\]
A \(\natural\)-1-\textit{structure} on \(\mathrm{P}\) is a
connection on the principal homogeneous \(\mathrm{H}\)-space \(\mathrm{P}\), relative to
\(\mathrm{G}_1\times\mathrm{G}_2\to\mathrm{G}_2\) (the covariant
derivative is defined only for vector fields parallel
to \(\mathrm{G}_1\)), such that \(\nu_1\) and \(\nu_2\) are
horizontal. The
\(\natural\)-2-\textit{structures} are defined similarly. A
\(\natural\)-\textit{structure} is the datum of a
\(\natural\)-1-structure and a \(\natural\)-2-structure.
It is also a connection on the principal homogeneous \(\mathrm{H}\)-space
\(\mathrm{P}\), such that \(\nu_1\) and \(\nu_2\) are horizontal.

The \textit{curvature} \(\mathrm{R}\) of a \(\natural\)-biextension
\(\mathrm{P}\) is the curvature of the connection on \(\mathrm{P}\).
It is a translation-invariant 2-form on
\(\mathrm{G}_1\times\mathrm{G}_2\), with values in
\(\operatorname{Lie}(\mathrm{H})\), i.e. it is an alternating form
on
\(\operatorname{Lie}(\mathrm{G}_1)\times
\operatorname{Lie}(\mathrm{G}_2)\), with values in
\(\operatorname{Lie}(\mathrm{H})\).
Its restriction to \(\operatorname{Lie}(\mathrm{G}_1)\) and to
\(\operatorname{Lie}(\mathrm{G}_2)\) is zero, so that
\(\mathrm{R}\) defines a pairing (also called the
\textit{curvature} of \(\mathrm{P}\))
\[
\Phi:\operatorname{Lie}(\mathrm{G}_1)\otimes
\operatorname{Lie}(\mathrm{G}_2)\to\operatorname{Lie}(\mathrm{H}),
\tag*{\textbf{(10.2.7.3)}}
\]
with
\[
\mathrm{R}(g_1+g_2,g_1'+g_2')
=\Phi(g_1,g_2')-\Phi(g_1',g_2).
\]

For complexes of smooth groups
\(\mathrm{K}_i:[\mathrm{A}_i\to\mathrm{B}_i]\), a
\(\natural\)-\textit{biextension}
(resp. \(\natural\)-\(i\)-\textit{biextension}) of
\(\mathrm{K}_1\) and \(\mathrm{K}_2\) by \(\mathrm{H}\) is a
\(\natural\)-biextension (resp. \(\natural\)-\(i\)-biextension)
of \(\mathrm{B}_1\) and \(\mathrm{B}_2\) by \(\mathrm{H}\),
equipped with trivializations (as a \(\natural\)-biextension,
resp. \(\natural\)-\(i\)-biextension) as in (10.2.1).

\textit{Proposition} \textbf{(10.2.7.4)}. — \textit{Let
\(\mathrm{P}^{\natural}\) be the biextension of
\(\mathrm{M}_1^{\natural}\) and \(\mathrm{M}_2^{\natural}\) by
\(\mathbf{G}_m\) obtained by pulling back \(\mathrm{P}\). There is
one and only one \(\natural\)-structure
on \(\mathrm{P}^{\natural}\).}

Let us prove uniqueness: we must check that every
\(\natural\)-structure on the trivial biextension of
\(\mathrm{M}_1^{\natural}\) and
\(\mathrm{M}_2^{\natural}\) by \(\mathbf{G}_m\) is trivial. A \(\natural\)-structure
is first of all a connection on the trivial principal homogeneous \(\mathbf{G}_m\)-space
over \(\mathrm{G}_1^{\natural}\times\mathrm{G}_2^{\natural}\),
i.e. a field of forms on
\(\mathrm{G}_1^{\natural}\times\mathrm{G}_2^{\natural}\),
say \(\Gamma_{g_1,g_2}(t_1+t_2)\).
The axiom for \(\natural\)-structures implies that
\[
\Gamma_{g_1g_2}(t_1+t_2)
=\Phi_1(g_1,t_2)+\Phi_2(t_1,g_2)
\]
with \(\Phi_1\) and \(\Phi_2\) biadditive.
\(\mathrm{G}_i^{\natural}\) is an extension of
\(\mathrm{X}_i\otimes k\) by the universal extension
\(\mathrm{G}_i'\) of \(\mathrm{G}_i\) by an additive group.
Every morphism \(\mathrm{G}_i'\to\mathbf{G}_a\) is trivial.
Consequently, \(\Phi_1\) factors through
\((\mathrm{X}_1\otimes k)\otimes
\operatorname{Lie}(\mathrm{G}_2^{\natural})\).
Finally, the restriction of \(\Gamma\) to
\(\mathrm{X}_1\times\mathrm{G}_2^{\natural}\) must be trivial,
hence so must the restriction of \(\Phi_1\) to
\(\mathrm{X}_1\times\operatorname{Lie}(\mathrm{G}_2^{\natural})\).
We therefore have \(\Phi_1=0\). Similarly, \(\Phi_2=0\), and \(\Gamma=0\).

To prove existence, it suffices to construct a
\(\natural\)-2-structure on the biextension of
\(\mathrm{M}_1^{\natural}\) and \(\mathrm{M}_2\) by
\(\mathbf{G}_m\) obtained by pulling back \(\mathrm{P}\): by
pullback, we shall obtain a \(\natural\)-2-structure on
\(\mathrm{P}^{\natural}\) and, symmetrically, a
\(\natural\)-1-structure.

For \(g\in\mathrm{G}_1\), \(\mathrm{P}_g\) is an extension
of \(\mathrm{G}_2\) by \(\mathbf{G}_m\).
Let \(\mathrm{C}_g\) be the set of \(\natural\)-structures on
\(\mathrm{P}_g\). It is a principal homogeneous space under
\(\operatorname{Lie}(\mathrm{G}_2)^{*}\).
For \(g\in\mathrm{G}_1\), the \(\mathrm{C}_g\) are the fibers
of a principal homogeneous space \(\mathrm{C}\) over \(\mathrm{G}_1\); Baer
addition of \(\natural\)-extensions even makes \(\mathrm{C}\)
an extension of \(\mathrm{G}_1\) by
\Page{66}
\(\operatorname{Lie}(\mathrm{G}_2)^{*}\).
Lift \(u_1:\mathrm{X}\to\mathrm{G}_1\) to
\(u_1':\mathrm{X}\to\mathrm{C}\) by associating to
\(x\in\mathrm{X}\) the trivial connection on the trivialized
extension \(\mathrm{P}_{u_1(x)}\).
We thus obtain a \(\natural\)-2-structure on the biextension
of \([\mathrm{X}\to\mathrm{C}]\) and \(\mathrm{M}_2\) by
\(\mathbf{G}_m\) obtained by pulling back \(\mathrm{P}\).

In view of the universal property of \(\mathrm{M}_1^{\natural}\),
there exists a unique commutative diagram
\[
\begin{tikzcd}[row sep=3em, column sep=1.5em]
& {[\mathrm{X}\longrightarrow\mathrm{G}]} &\\
\mathrm{M}_1^{\natural}
  \arrow[ur] \arrow[rr, "v"]
&& {[\mathrm{X}\longrightarrow\mathrm{C}]} \arrow[ul]
\end{tikzcd}
\]
taking a pullback along \(v\), we find the required
\(\natural\)-2-structure.

By definition, the pairing (10.2.7) is the negative of the
curvature of the \(\natural\)-biextension (10.2.7.4).

\textit{Proposition} \textbf{(10.2.8)}. — \textit{Over
\(\mathbf{C}\), the pairing \((10.2.7)\) is the complexification
of the pairing \((10.2.3)\).}

The \(\natural\)-biextension \(\mathrm{P}^{\natural}\) defines,
in the analytic category, a \(\natural\)-biextension of
\([\mathrm{T}_{\mathbf{Z}}(\mathrm{M}_1)\to
\mathrm{T}_{\mathbf{C}}(\mathrm{M}_1)]\) and
\([\mathrm{T}_{\mathbf{Z}}(\mathrm{M}_2)\to
\mathrm{T}_{\mathbf{C}}(\mathrm{M}_2)]\) by \(\mathbf{G}_m\).
Proposition (10.2.8) then follows from the

\textit{Lemma} \textbf{(10.2.9)}. — \textit{Let
\(\mathrm{V}_{\mathbf{Z}}\) and \(\mathrm{W}_{\mathbf{Z}}\)
be two free \(\mathbf{Z}\)-modules of finite rank, and \(\mathrm{P}\)
a \(\natural\)-biextension of
\([\mathrm{V}_{\mathbf{Z}}\to\mathrm{V}_{\mathbf{C}}]\) and
\([\mathrm{W}_{\mathbf{Z}}\to\mathrm{W}_{\mathbf{C}}]\) by
\(\mathbf{G}_m\), with underlying biextension \(\mathrm{P}\).}

\textit{Then, the curvature
\(\psi:\mathrm{V}_{\mathbf{C}}\otimes\mathrm{W}_{\mathbf{C}}
\to\mathbf{C}\) of \(\mathrm{P}^{\natural}\) coincides with
the negative of the complexification of the form corresponding to
\(\mathrm{P}\) under \((10.2.3.4)\).}

Take \(\mathrm{P}=\mathrm{P}(\psi_1,\psi_2)\) (10.2.3.4).
A \(\natural\)-structure on \(\mathrm{P}\) is first of all a
connection on the trivial principal homogeneous \(\mathbf{G}_m\)-space over
\(\mathrm{V}_{\mathbf{C}}\times\mathrm{W}_{\mathbf{C}}\),
i.e. a field of forms \(\Gamma_{v,w}(v'+w')\) on
\(\mathrm{V}_{\mathbf{C}}\times\mathrm{W}_{\mathbf{C}}\).
This connection defines a \(\natural\)-structure on this
principal homogeneous space if and only if
\[
\Gamma_{v,w}(v'+w')=\Phi_1(v,w')+\Phi_2(v',w)
\]
with \(\Phi_1\) and \(\Phi_2\) bilinear.
The trivializations \(\exp\psi_1(v_{\mathbf{Z}},w_{\mathbf{C}})\)
and \(\exp\psi_2(v_{\mathbf{C}},w_{\mathbf{Z}})\) are
horizontal if and only if \(\Phi_i=-\psi_i\).
The curvature of the connection is the field of 2-forms \(d\Gamma\):
\[
\mathrm{R}_{v,w}(v'+w',v''+w'')
=\Phi_1(v',w'')+\Phi_2(v'',w')
-\Phi_1(v'',w')-\Phi_2(v',w'').
\]
The curvature of \(\mathrm{P}\) is therefore the form
\[
\Phi_1(v',w'')-\Phi_2(v',w'')
=-(\psi_1(v',w'')-\psi_2(v',w''))
=-\psi(v',w'').
\]

\textbf{(10.2.10)} Let \(\mathrm{M}\) be a 1-motive over
\(\mathbf{C}\). To \(\mathrm{M}\) corresponds a mixed Hodge
structure \(\mathrm{T}(\mathrm{M})\) of type
\(\{(0,0)\ (0,-1)\ (-1,0)\ (-1,-1)\}\).
The mixed Hodge structure
\(\operatorname{Hom}(\mathrm{T}(\mathrm{M}),\mathbf{Z}(1))\)
is again of this type, hence defines a 1-motive \(\mathrm{M}^{*}\).
The obvious map
\(\mathrm{T}(\mathrm{M})\otimes\mathrm{T}(\mathrm{M}^{*})
\to\mathbf{Z}(1)\) defines a biextension \(\mathrm{P}\) of
\(\mathrm{M}\) and \(\mathrm{M}^{*}\) by \(\mathbf{G}_m\).
The 1-motive \(\mathrm{M}^{*}\), equipped with the biextension
\(\mathrm{P}\), is the \textit{Cartier dual} of \(\mathrm{M}\).
We call \(\mathrm{P}\) the \textit{Poincaré biextension}.

\Page{67}
\textbf{(10.2.11)} Construction (10.2.10) can be made
purely algebraic. To a 1-motive
\(\mathrm{M}=(\mathrm{X},\mathrm{A},\mathrm{T},\mathrm{G},u)\)
over an algebraically closed field \(k\), we shall associate
a \textit{Cartier dual}
\(\mathrm{M}^{*}=(\mathrm{X}',\mathrm{A}',\mathrm{T}',
\mathrm{G}',u')\):

\textit{a)} \(\mathrm{X}'\) is the character group of
\(\mathrm{T}\), \(\mathrm{A}'\) the dual abelian variety of
\(\mathrm{A}\), and \(\mathrm{T}'\) the torus with character
group \(\mathrm{X}\).

\textit{b)} Let us first treat the case where \(\mathrm{T}=0\).
An extension of \(\mathrm{M}\) by \(\mathbf{G}_m\) then consists
of an extension \(\mathrm{E}\) of \(\mathrm{A}\) by
\(\mathbf{G}_m\), together with a trivialization \(\tilde{u}\) of
the extension \(u^{*}\mathrm{E}\) of \(\mathrm{X}\) by
\(\mathbf{G}_m\)
\[
\begin{tikzcd}[row sep=2.5em, column sep=2em]
&&& \mathrm{X} \arrow[dl, "\tilde{u}"'] \arrow[d, "u"] &\\
0 \arrow[r] & \mathbf{G}_m \arrow[r] &
\mathrm{E} \arrow[r] & \mathrm{A} \arrow[r] & 0
\end{tikzcd}
\]
We have an exact sequence
\[
0\to\operatorname{Hom}(\mathrm{X},\mathbf{G}_m)
\to\operatorname{Ext}^{1}(\mathrm{M},\mathbf{G}_m)
\to\operatorname{Ext}^{1}(\mathrm{A},\mathbf{G}_m)\to0.
\]
Extensions of \(\mathrm{M}\) by \(\mathbf{G}_m\) have
no nontrivial automorphism. The functor of isomorphism
classes of extensions is representable; we define
\(\mathrm{G}'\) to be the algebraic group representing
this functor; it is an extension of \(\mathrm{A}'\) by
\(\mathrm{T}'\). We have a biextension \(\mathrm{P}\)
of \(\mathrm{M}\) and \(\mathrm{G}'\) by \(\mathbf{G}_m\).

\textit{c)} In the general case, set
\(\mathrm{G}'=\operatorname{Ext}^{1}
(\mathrm{M}/\mathrm{W}_{-2}\mathrm{M},\mathbf{G}_m)\).
The natural biextension \(\mathrm{P}''\) of
\(\mathrm{M}/\mathrm{W}_{-2}\mathrm{M}\) and \(\mathrm{G}'\)
by \(\mathbf{G}_m\) induces a biextension \(\mathrm{P}'\)
of \(\mathrm{M}\) and \(\mathrm{G}'\) by \(\mathbf{G}_m\).
For each \(\chi\in\mathrm{X}'\), the extension \(\mathrm{M}\)
of \(\mathrm{M}/\mathrm{W}_{-2}\mathrm{M}\) by \(\mathrm{T}\)
defines an extension of
\(\mathrm{M}/\mathrm{W}_{-2}\mathrm{M}\) by \(\mathbf{G}_m\),
whence \(u'(\chi)\in\mathrm{G}'\).
\[
\begin{tikzcd}[row sep=3em, column sep=2em]
0 \arrow[r] &
\mathrm{T} \arrow[r] \arrow[d] &
\mathrm{M} \arrow[r] \arrow[d, "v"] &
\mathrm{M}/\mathrm{W}_{-2}\mathrm{M}
  \arrow[r] \arrow[d, equal] &
0\\
0 \arrow[r] &
\mathbf{G}_m \arrow[r] &
\mathrm{P}''_{u(\chi)} \arrow[r] &
\mathrm{M}/\mathrm{W}_{-2}\mathrm{M} \arrow[r] &
0
\end{tikzcd}
\]
The map \(v\) trivializes the extension
\(\mathrm{P}'_{u(\chi)}\) of \(\mathrm{M}\) by \(\mathbf{G}_m\)
\[
\begin{tikzcd}[row sep=3em, column sep=2em]
0 \arrow[r] &
\mathbf{G}_m \arrow[r] \arrow[d, equal] &
\mathrm{P}_{u(\chi)} \arrow[r] \arrow[d] &
\mathrm{M} \arrow[r] \arrow[d] \arrow[dl, "v"'] &
0\\
& \mathbf{G}_m \arrow[r] &
\mathrm{P}'_{u(\chi)} \arrow[r] &
\mathrm{M}/\mathrm{W}_{-2}\mathrm{M} &
\end{tikzcd}
\]
These trivializations make \(\mathrm{P}'\) a biextension of
\(\mathrm{M}\) and
\(\mathrm{M}^{*}=[\mathrm{X}'\xrightarrow{u'}\mathrm{G}']\)
by \(\mathbf{G}_m\). The 1-motive \(\mathrm{M}^{*}\), equipped with
\(\mathrm{P}'\), is the desired Cartier dual.

\Page{68}
\textbf{(10.2.12)} This construction allows us to give a
more symmetric description of 1-motives.

Let \(\mathrm{M}=(\mathrm{X},\mathrm{A},\mathrm{T},\mathrm{G},u)\)
be a 1-motive, with Cartier dual
\(\mathrm{M}^{*}=(\mathrm{X}',\mathrm{A}',\mathrm{T}',
\mathrm{G}',u')\), and let \(\mathrm{P}_0\) be the Poincaré
biextension. The morphisms \(u\) and \(u'\) define, by passing
to the quotient, morphisms
\[
\bar{u}:\mathrm{X}\to\mathrm{A}
\qquad\text{and}\qquad
\bar{u}':\mathrm{X}'\to\mathrm{A}'.
\tag*{\textbf{(10.2.12.1)}}
\]
The biextension of \(\mathrm{G}\) and \(\mathrm{G}'\) by
\(\mathbf{G}_m\) obtained from \(\mathrm{P}_0\) is the pullback
of the Poincaré biextension \(\mathrm{P}\) of \(\mathrm{A}\)
and \(\mathrm{A}'\) by \(\mathbf{G}_m\). By hypothesis, the
trivializations (of the pullbacks) of \(\mathrm{P}\) on
\(\mathrm{X}\times\mathrm{G}'\) and
\(\mathrm{G}\times\mathrm{X}'\) agree on
\(\mathrm{X}\times\mathrm{X}'\), whence
\[
\text{a trivialization }\psi\text{ of }
(\bar{u}\times\bar{u}')^{*}\mathrm{P}.
\tag*{\textbf{(10.2.12.2)}}
\]
A morphism \(f:\mathrm{M}_1\to\mathrm{M}_2\), with transpose
(in an obvious sense) \(\mathrm{M}_2^{*}\to\mathrm{M}_1^{*}\),
gives rise to commutative diagrams
\[
\begin{tikzcd}[row sep=2.5em, column sep=2em]
\mathrm{X}_1 \arrow[r] \arrow[d, "f_{\mathrm{X}}"] &
\mathrm{A}_1 \arrow[d, "f_{\mathrm{A}}"]\\
\mathrm{X}_2 \arrow[r] & \mathrm{A}_2
\end{tikzcd}
\qquad
\begin{tikzcd}[row sep=2.5em, column sep=2em]
\mathrm{X}_1' \arrow[r] & \mathrm{A}_1'\\
\mathrm{X}_2' \arrow[r] \arrow[u, "f_{\mathrm{X}}'"'] &
\mathrm{A}_2' \arrow[u, "f_{\mathrm{A}}'"']
\end{tikzcd}
\tag*{\textbf{(10.2.12.3)}}
\]
Moreover, \(f_{\mathrm{A}}\) and \(f_{\mathrm{A}}'\) are transposes
and, \textit{via} the isomorphism
\((1,f_{\mathrm{A}}')^{*}\mathrm{P}_1
\simeq(f_{\mathrm{A}},1)^{*}\mathrm{P}_2\), we have
\[
\psi_1(x_1,f_{\mathrm{X}}'(x_2'))
=\psi_2(f_{\mathrm{X}}(x_1),x_2')
\tag*{\textbf{(10.2.12.4)}}
\]
\[
\begin{tikzcd}[row sep=4em, column sep=2em]
\mathrm{X}_1\times\mathrm{X}_1' \arrow[r, "\psi_1"] &
(\mathrm{P}_1\text{ over }\mathrm{A}_1\times\mathrm{A}_1')
\\
\mathrm{X}_1\times\mathrm{X}_2'
  \arrow[u] \arrow[d] \arrow[r, dotted] &
(1,f_{\mathrm{A}}')^{*}\mathrm{P}_1
=(f_{\mathrm{A}},1)^{*}\mathrm{P}_2
\text{ over }\mathrm{A}_1\times\mathrm{A}_2'
  \arrow[u] \arrow[d]
\\
\mathrm{X}_2\times\mathrm{X}_2' \arrow[r, "\psi_2"] &
(\mathrm{P}_2\text{ over }\mathrm{A}_2\times\mathrm{A}_2').
\end{tikzcd}
\]

\textbf{(10.2.13)} We thus obtain a functor from the category
of 1-motives to the following category. The objects consist of:

\noindent\textit{a)} two abelian varieties in duality
\(\mathrm{A}\) and \(\mathrm{A}'\);

\noindent\textit{b)} two free \(\mathbf{Z}\)-modules of finite
rank \(\mathrm{X}\) and \(\mathrm{X}'\);

\noindent\textit{c)} morphisms
\(v:\mathrm{X}\to\mathrm{A}\) and
\(v':\mathrm{X}'\to\mathrm{A}'\);

\noindent\textit{d)} a trivialization \(\psi\) of the pullback
along \((v,v')\) of the Poincaré biextension of
\(\mathrm{A}\) and \(\mathrm{A}'\) by \(\mathbf{G}_m\).

% ---- 069-077 ----
\Page{69}
The arrows consist of the diagrams (10.2.12.3) satisfying the conditions of (10.2.12).

\textit{Proposition} \textbf{(10.2.14)}. — \textit{The functor \textup{(10.2.13)} is an equivalence of categories.}

The verification is left as an exercise to the reader.

\section*{10.3. Algebraic interpretation of mixed \(\mathbf{H}^{1}\) : the case of curves.}

\textbf{(10.3.1)} Let \(\mathrm{X}_0\) be an algebraic curve over an algebraically closed field \(k\). Let \(p:\mathrm{X}'\to\mathrm{X}_0\) be the normalization of \(\mathrm{X}_0\), and let \(\overline{\mathrm{X}}'\) be the nonsingular projective curve of which \(\mathrm{X}'\) is a dense open subset. We shall denote by \(\mathrm{X}\) (resp. by \(\overline{\mathrm{X}}\)) the curve obtained from \(\mathrm{X}'\) (resp. from \(\overline{\mathrm{X}}'\)) by contracting to a point each of the finite sets \(p^{-1}(s)\) for \(s\in\mathrm{X}_0\).
\[
\begin{tikzcd}[column sep=2em,row sep=1.5em]
\mathrm{X}' \arrow[r,hook,"j'"'] \arrow[d,"r"']
  & \overline{\mathrm{X}}' \arrow[d,"\bar r"'] \\
\mathrm{X} \arrow[r,hook,"j"'] \arrow[d,"p_0"']
  & \overline{\mathrm{X}} \\
\mathrm{X}_0 &
\end{tikzcd}
\]

We can characterize \(\mathrm{X}\) and \(\overline{\mathrm{X}}\) by the following properties.

\textit{a)} \(p_0\) is finite and radicial, hence induces isomorphisms
\[
\mathrm{H}^{\bullet}(\mathrm{X}_0,\mathbf{Z}_{\ell})
\longrightarrow
\mathrm{H}^{\bullet}(\mathrm{X},\mathbf{Z}_{\ell})
\qquad
(\ell\text{ prime to the characteristic}).
\]

\textit{b)} The singularities of \(\mathrm{X}\) are analytically isomorphic to those of a union of coordinate axes in affine space.

\textit{c)} \(\overline{\mathrm{X}}\) is projective and \(\overline{\mathrm{X}}-\mathrm{X}\) is a finite set of points at which \(\overline{\mathrm{X}}\) is smooth.

\textbf{(10.3.2)} The Picard group \(\operatorname{Pic}(\overline{\mathrm{X}})\) admits the following dévissage.

\textit{a)} Let \(\mathrm{I}\) be the set of irreducible components of \(\overline{\mathrm{X}}\) (= the set of irreducible components of \(\mathrm{X}_0\)). For each irreducible component \(\mathrm{Y}\) of \(\overline{\mathrm{X}}\), the degree of the restriction to \(\mathrm{Y}\) of an invertible sheaf is a homomorphism
\[
\deg_{\mathrm{Y}}:\operatorname{Pic}(\overline{\mathrm{X}})
\longrightarrow\mathbf{Z}:
\mathcal{L}\longrightarrow\deg_{\mathrm{Y}}(\mathcal{L}).
\]
This gives
\[
\mathit{deg}:\operatorname{Pic}(\overline{\mathrm{X}})
\longrightarrow\mathbf{Z}^{\mathrm{I}}.
\]
This homomorphism is surjective, with kernel \(\operatorname{Pic}^{0}(\overline{\mathrm{X}})\).

\textit{b)} The morphism from \(\operatorname{Pic}^{0}(\overline{\mathrm{X}})\) to the abelian variety \(\operatorname{Pic}^{0}(\overline{\mathrm{X}}')\) is surjective. Its kernel is a torus.

\Page{70}
\textbf{(10.3.3)} Let \(\mathrm{S}\) be the finite set \(\overline{\mathrm{X}}-\mathrm{X}\), \(\alpha_0:\mathrm{S}\to\mathrm{I}\) the map which assigns to \(s\in\mathrm{S}\) the unique irreducible component of \(\overline{\mathrm{X}}\) such that \(s\in\alpha_0(s)\) and let \(\alpha:\mathbf{Z}^{\mathrm{S}}\to\mathbf{Z}^{\mathrm{I}}\) be induced by \(\alpha_0\).

Each \(s\in\mathrm{S}\) defines an invertible sheaf \(\mathcal{O}(s)\), hence \(u_0:\mathbf{Z}^{\mathrm{S}}\to\operatorname{Pic}(\overline{\mathrm{X}})\). We have \(\mathit{deg}\,u_0=\alpha\).

\textit{Definition} \textbf{(10.3.4)}. — \textit{The \textup{motivic} \(\mathrm{H}^{1}\) of \(\mathrm{X}_0\), denoted by \(\mathrm{H}^{1}_{m}(\mathrm{X}_0)(1)\) is the \(1\)-motive}
\[
\operatorname{Ker}(\alpha)\longrightarrow
\operatorname{Pic}^{0}(\overline{\mathrm{X}}).
\]

By definition, we therefore have \(\mathrm{H}^{1}_{m}(\mathrm{X}_0)(1)=\mathrm{H}^{1}_{m}(\mathrm{X})(1)\).

\textit{Remarks} \textbf{(10.3.5)}. — (i) We have
\[
\begin{aligned}
\operatorname{Gr}^{\mathrm{W}}_0(\mathrm{H}^{1}_{m}(\mathrm{X})(1))
  &=\operatorname{Ker}(\alpha:\mathbf{Z}^{\mathrm{S}}\to\mathbf{Z}^{\mathrm{I}})\\
\operatorname{Gr}^{\mathrm{W}}_{-1}(\mathrm{H}^{1}_{m}(\mathrm{X})(1))
  &=\operatorname{Pic}^{0}(\overline{\mathrm{X}}').
\end{aligned}
\]

(ii) The cokernel of \(\alpha\) is torsion-free : it is identified with the free \(\mathbf{Z}\)-module with basis \(\mathrm{I}-\alpha_0(\mathrm{S})\). We have a short exact sequence of complexes and a quasi-isomorphism
\[
\begin{tikzcd}[column sep=small,row sep=3em]
0 \arrow[r]
 & {[\operatorname{Ker}(\alpha)\longrightarrow
       \operatorname{Pic}^{0}(\overline{\mathrm{X}}_0)]} \arrow[r]
 & {[\mathbf{Z}^{\mathrm{S}}\longrightarrow
       \operatorname{Pic}(\overline{\mathrm{X}}_0)]} \arrow[r]
 & {[\operatorname{Im}(\alpha)\longrightarrow\mathbf{Z}^{\mathrm{I}}]}
       \arrow[r] \arrow[d,"\wr"]
 & 0 \\
 & & & {[0\longrightarrow\operatorname{coker}(\alpha)]} &
\end{tikzcd}
\]

\textit{Construction} \textbf{(10.3.6)}. — \textit{For \(n\) an integer \(>0\) prime to the characteristic of \(k\), we construct an isomorphism}
\[
\mathrm{H}^{1}(\mathrm{X},\mathbf{Z}/n\mathbf{Z})(1)
=\mathrm{H}^{1}(\mathrm{X},\mu_n)
\simeq
\mathrm{T}_{\mathbf{Z}/n\mathbf{Z}}(\mathrm{H}^{1}_{m}(\mathrm{X})(1)).
\]

We know that \(\mathrm{H}^{1}(\mathrm{X},\mu_n)\) is identified with the set of isomorphism classes of invertible sheaves \(\mathcal{L}\) on \(\mathrm{X}\), equipped with an isomorphism \(\mathcal{L}^{\otimes n}\simeq\mathcal{O}\). Since \(p_0:\mathrm{X}_0\to\mathrm{X}\) is finite and radicial, we have \(\mathrm{H}^{1}(\mathrm{X},\mathbf{Z}/n)\xrightarrow{\sim}\mathrm{H}^{1}(\mathrm{X}_0,\mathbf{Z}/n)\), and \(\mathrm{H}^{1}(\mathrm{X},\mu_n)\) is again the set of isomorphism classes of pairs \((\mathcal{L},\alpha)\) where \(\mathcal{L}\) is an invertible sheaf on \(\mathrm{X}_0\), and \(\alpha\) an isomorphism \(\alpha:\mathcal{L}^{\otimes n}\simeq\mathcal{O}\).

Let \((\mathcal{L},\alpha)\) be such a pair. The invertible sheaf \(\mathcal{L}\) extends to \(\overline{\mathcal{L}}\) on \(\overline{\mathrm{X}}\), and there exists a divisor \(\mathrm{D}\) supported on \(\mathrm{S}\) such that \(\alpha\) extends to \(\bar\alpha:\overline{\mathcal{L}}^{\otimes n}\xrightarrow{\sim}\mathcal{O}(\mathrm{D})\). If there exists an isomorphism \(\beta\) between \(\overline{\mathcal{L}}^{\otimes n}\) and \(\mathcal{O}(\mathrm{E})\), this isomorphism is uniquely determined up to multiplication by an element of the divisible group \(\mathrm{H}^{0}(\overline{\mathrm{X}},\mathcal{O}^{*})\). We deduce that \((\overline{\mathcal{L}},\mathrm{D})\) determines \((\mathcal{L},\alpha)\) up to isomorphism. For a pair \((\overline{\mathcal{L}},\mathrm{D})\) to come from a suitable pair \((\mathcal{L},\alpha)\), it is necessary and sufficient that \(n[\mathcal{L}]=[\mathcal{O}(\mathrm{D})]\) in \(\operatorname{Pic}(\overline{\mathrm{X}})\). It comes from \((\mathcal{O}_{\mathrm{X}_0},0)\) if and only if it is of the form \((\mathcal{O}_{\mathrm{X}_0}(\mathrm{D}),n\mathrm{D})\).

\Page{71}
This identifies \(\mathrm{H}^{1}(\mathrm{X},\mu_n)\) with the \(\mathrm{H}^{0}\) of the tensor product complex of \([\mathbf{Z}^{\mathrm{S}}\to\operatorname{Pic}(\overline{\mathrm{X}})]\) (degrees \(0\) and \(1\)) and \([\mathbf{Z}\xrightarrow{n}\mathbf{Z}]\) (degrees \(-1\) and \(0\)) :
\[
\begin{tikzcd}[column sep=2em,row sep=2.5em]
\mathbf{Z}^{\mathrm{S}} \arrow[r,"u_0"]
 & \operatorname{Pic}(\overline{\mathrm{X}})\\
\mathbf{Z}^{\mathrm{S}} \arrow[r,"u_0"] \arrow[u,"n"']
 & \operatorname{Pic}(\overline{\mathrm{X}}) \arrow[u,"-n"']
\end{tikzcd}
\]
\[
\mathrm{H}^{1}(\mathrm{X},\mu_n)
\simeq
\mathrm{H}^{0}\bigl(
[\mathbf{Z}^{\mathrm{S}}\to\operatorname{Pic}(\overline{\mathrm{X}})]
\overset{\mathrm{L}}{\otimes}\mathbf{Z}/n\mathbf{Z}\bigr).
\]

Since \(\operatorname{coker}(\alpha)\) is torsion-free, we see from (10.3.5) (ii) that
\[
\begin{aligned}
\mathrm{H}^{1}(\mathrm{X},\mu_n)
&\simeq
\mathrm{H}^{0}\bigl(
[\mathbf{Z}^{\mathrm{S}}\to\operatorname{Pic}(\overline{\mathrm{X}})]
\overset{\mathrm{L}}{\otimes}\mathbf{Z}/n\mathbf{Z}\bigr)
\xleftarrow{\sim}
\mathrm{H}^{0}\bigl(
\mathrm{H}^{1}_{m}(\mathrm{X})
\overset{\mathrm{L}}{\otimes}\mathbf{Z}/n\mathbf{Z}\bigr)\\
&\simeq
\mathrm{T}_{\mathbf{Z}/n\mathbf{Z}}(\mathrm{H}^{1}_{m}(\mathrm{X})(1)).
\end{aligned}
\]

\textit{Variant} \textbf{(10.3.7)}. — \textit{For \(\ell\) prime to the characteristic \(p\) of \(k\), we have}
\[
\mathrm{H}^{1}(\mathrm{X}_0,\mathbf{Z}_{\ell}(1))
\simeq
\mathrm{T}_{\ell}(\mathrm{H}^{1}_{m}(\mathrm{X}_0)(1)).
\]

\textit{Construction} \textbf{(10.3.8)}. — \textit{Let \(\mathrm{X}_0\) be a curve over \(\mathbf{C}\). We shall construct an isomorphism of mixed Hodge structures}
\[
\mathrm{H}^{1}(\mathrm{X}_0)(1)
\simeq\mathrm{T}(\mathrm{H}^{1}_{m}(\mathrm{X}_0)(1)).
\]
We may and shall assume that \(\mathrm{X}=\mathrm{X}_0\).

\textit{Construction} \textbf{(10.3.9)}. — (i) \textit{Let \(j_*^{m}\mathcal{O}_{\mathrm{X}}^{*}\) be the subsheaf of meromorphic functions in \(j_*\mathcal{O}_{\mathrm{X}}^{*}\). We have}
\[
\mathrm{H}^{1}(\mathrm{X},\mathbf{Z}(1))
=\mathbf{H}^{1}\bigl(\overline{\mathrm{X}},
[\mathcal{O}_{\overline{\mathrm{X}}}\xrightarrow{\exp}
j_*^{m}\mathcal{O}_{\mathrm{X}}^{*}]\bigr).
\]
\[
\mathrm{H}^{1}(\mathrm{X},\mathbf{C})
=\mathbf{H}^{1}\bigl(\overline{\mathrm{X}},
[\mathcal{O}_{\overline{\mathrm{X}}}\xrightarrow{d}
r_*\Omega^{1}_{\overline{\mathrm{X}}'}(\log\mathrm{S})]\bigr).
\tag*{(ii)}
\]

(iii) \textit{The inclusion of \(\mathrm{H}^{1}(\mathrm{X},\mathbf{Z})\) in \(\mathrm{H}^{1}(\mathrm{X},\mathbf{C})\) is defined by the morphism of complexes}
\[
\begin{tikzcd}[column sep=2em,row sep=2.5em]
\mathcal{O}_{\overline{\mathrm{X}}}
  \arrow[r,"\exp"] \arrow[d,"\wr"']
 & j_*^{m}\mathcal{O}_{\mathrm{X}}^{*} \arrow[d,"df/f"]\\
\mathcal{O}_{\overline{\mathrm{X}}} \arrow[r]
 & r_*\Omega^{1}_{\overline{\mathrm{X}}'}(\log\mathrm{S})
\end{tikzcd}
\]

Consider the commutative diagram
\[
\begin{tikzcd}[column sep=0.7em,row sep=3.5em]
\mathrm{H}^{1}(\mathrm{X},\mathbf{Z}(1))
  \arrow[r,"\sim"] \arrow[d]
  \arrow[dr,phantom,"\text{\textcircled{1}}" description]
&
\mathbf{H}^{1}(\mathrm{X},
 [\mathcal{O}_{\mathrm{X}}\xrightarrow{\exp}\mathcal{O}_{\mathrm{X}}^{*}])
  \arrow[d]
  \arrow[dr,phantom,"\text{\textcircled{2}}" description]
&
\mathbf{H}^{1}(\overline{\mathrm{X}},
 [j_*\mathcal{O}_{\mathrm{X}}\xrightarrow{\exp}j_*\mathcal{O}_{\mathrm{X}}^{*}])
  \arrow[l,"\sim"'] \arrow[d]
  \arrow[dr,phantom,"\text{\textcircled{3}}" description]
&
\mathbf{H}^{1}(\overline{\mathrm{X}},
 [\mathcal{O}_{\overline{\mathrm{X}}}\xrightarrow{\exp}
 j_*^{m}\mathcal{O}_{\mathrm{X}}^{*}])
  \arrow[l,"\sim"'] \arrow[d]
\\
\mathrm{H}^{1}(\mathrm{X},\mathbf{C}) \arrow[r]
&
\mathbf{H}^{1}(\mathrm{X},
 [\mathcal{O}_{\mathrm{X}}\xrightarrow{d}
 r_*\Omega^{1}_{\overline{\mathrm{X}}'}])
&
\mathbf{H}^{1}(\overline{\mathrm{X}},
 [j_*\mathcal{O}_{\mathrm{X}}\longrightarrow
 j_*r_*\Omega^{1}_{\mathrm{X}'}])
 \arrow[l]
&
\mathbf{H}^{1}(\overline{\mathrm{X}},
 [\mathcal{O}_{\overline{\mathrm{X}}}\longrightarrow
 r_*\Omega^{1}_{\overline{\mathrm{X}}'}(\log\mathrm{S})])
 \arrow[l,"\sim"']
\end{tikzcd}
\]

\Page{72}
The horizontal arrows of square 1 are isomorphisms, since the sequences
\[
0\longrightarrow\mathbf{Z}(1)
\longrightarrow\mathcal{O}_{\mathrm{X}}
\xrightarrow{\exp}\mathcal{O}_{\mathrm{X}}^{*}
\longrightarrow0
\]
\[
0\longrightarrow\mathbf{C}
\longrightarrow\mathcal{O}_{\mathrm{X}}
\xrightarrow{d}r_*\Omega^{1}_{\mathrm{X}'}
\longrightarrow0
\]
are exact. Those of 2 are isomorphisms because \(\mathrm{R}^{1}j_*\mathcal{O}_{\mathrm{X}}=0\). Those of 3 are isomorphisms because the morphisms of complexes defining them are quasi-isomorphisms. This diagram defines (10.3.9).

\textbf{(10.3.10)} Recall that if \(d:\mathrm{F}\to\mathrm{G}\) is a morphism of abelian sheaves on a space \(\mathrm{Z}\), then \(\mathbf{H}^{1}(\mathrm{Z},[\mathrm{F}\to\mathrm{G}])\) is identified with the set of isomorphism classes of pairs \((\mathrm{P},\alpha)\), where \(\mathrm{P}\) is an \(\mathrm{F}\)-torsor (= a principal homogeneous space under \(\mathrm{F}\)) and where \(\alpha\) is a trivialization of the \(\mathrm{G}\)-torsor \(d\mathrm{P}\).

In particular :

\textit{a)} \(\mathbf{H}^{1}(\overline{\mathrm{X}},[\mathcal{O}_{\overline{\mathrm{X}}}^{*}\to r_*\Omega^{1}_{\overline{\mathrm{X}}'}(\log\mathrm{S})])\) is identified with the set of isomorphism classes of pairs \((\mathcal{L},\alpha)\), where \(\mathcal{L}\) is an invertible sheaf on \(\overline{\mathrm{X}}\), and \(\alpha\) a connection on \(\bar r^{*}\mathcal{L}\), holomorphic on \(\mathrm{X}'\) and possibly having a simple pole at the points of \(\mathrm{S}\). Let \(a\) be the projection of \(\overline{\mathrm{X}}\) to \(\operatorname{Spec}(\mathbf{C})\) (= a point). In the \textit{fppf} topology, the sheaf \(\mathrm{R}^{1}a_*([\mathcal{O}_{\overline{\mathrm{X}}}^{*}\to r_*\Omega^{1}_{\overline{\mathrm{X}}'}(\log\mathrm{S})])\) is representable by a group scheme \(\operatorname{Pic}^{\natural}(\mathrm{X})\), whose set of points is \(\mathbf{H}^{1}(\overline{\mathrm{X}},[\mathcal{O}_{\overline{\mathrm{X}}}^{*}\to r_*\Omega^{1}_{\overline{\mathrm{X}}'}(\log\mathrm{S})])\). \(\operatorname{Pic}^{\natural}(\mathrm{X})\) is an extension of a subgroup of \(\operatorname{Pic}(\overline{\mathrm{X}})\), containing \(\operatorname{Pic}^{0}(\overline{\mathrm{X}})\), by \(\mathrm{H}^{0}(\overline{\mathrm{X}}',\Omega^{1}_{\overline{\mathrm{X}}'}(\log\mathrm{S}))\).

We define a map
\[
\mathbf{Z}^{\mathrm{S}}\longrightarrow\operatorname{Pic}^{\natural}(\mathrm{X})
\tag*{\textbf{(10.3.10.1)}}
\]
by associating to each divisor \(\mathrm{D}\) on \(\overline{\mathrm{X}}\), supported on \(\mathrm{S}\), the invertible sheaf \(\mathcal{O}(\mathrm{D})\).

\textit{b)} \(\mathrm{H}^{1}(\mathrm{X},\mathbf{C})=\mathbf{H}^{1}(\overline{\mathrm{X}},[\mathcal{O}_{\overline{\mathrm{X}}}\to r_*\Omega^{1}_{\overline{\mathrm{X}}'}(\log\mathrm{S})])\) is the Lie algebra of \(\operatorname{Pic}^{\natural}(\mathrm{X})\). It is the set of isomorphism classes of pairs \((\mathrm{P},\alpha)\), where \(\mathrm{P}\) is an \(\mathcal{O}_{\overline{\mathrm{X}}}\)-torsor and \(\alpha\) a connection as in \textit{a)} on \(\mathrm{P}\).

\textit{c)} \(\mathrm{H}^{1}(\mathrm{X},\mathbf{Z}(1))=\mathbf{H}^{1}(\overline{\mathrm{X}},[\mathcal{O}_{\overline{\mathrm{X}}}\to j_*^{m}\mathcal{O}_{\mathrm{X}}^{*}])\) is the set of isomorphism classes of pairs \((\mathrm{P},\alpha)\), where \(\mathrm{P}\) is an \(\mathcal{O}_{\overline{\mathrm{X}}}\)-torsor, and \(\alpha\) an isomorphism between the invertible sheaf \(\exp(\mathrm{P})\) and an invertible sheaf \(\mathcal{O}_{\mathrm{X}}(\mathrm{D})\) (\(\mathrm{D}\) supported on \(\mathrm{S}\)). The map
\[
\operatorname{Aut}(\mathrm{P})=\mathbf{C}
\longrightarrow\operatorname{Aut}(\exp(\mathrm{P}))=\mathbf{C}^{*}
\]
is surjective. This also allows us to identify \(\mathrm{H}^{1}(\mathrm{X},\mathbf{Z}(1))\) with the set of pairs \((p,d)\), where \(p\) is an isomorphism class of \(\mathcal{O}_{\overline{\mathrm{X}}}\)-torsors, i.e. an element of \(\operatorname{Lie}(\operatorname{Pic}\overline{\mathrm{X}})\), and \(d\in\operatorname{Ker}(\alpha)\), defining a divisor supported on \(\mathrm{S}\) having \(\exp(p)\) as its image in \(\operatorname{Pic}^{0}(\overline{\mathrm{X}})\). In other words, we have defined an isomorphism
\[
\mathrm{H}^{1}(\mathrm{X},\mathbf{Z}(1))
\xrightarrow{\sim}
\mathrm{T}_{\mathbf{Z}}(\mathrm{H}^{1}_{m}(\mathrm{X})(1)).
\tag*{\textbf{(10.3.10.2)}}
\]

From part (i) of the following lemma, we deduce that this isomorphism is compatible with the weight filtration.

\Page{73}
\textit{Lemma} \textbf{(10.3.11)}. — (i) \textit{We have}
\[
\mathrm{W}^{1}(\mathrm{H}^{1}(\mathrm{X},\mathbf{Z}))
=\operatorname{Im}\bigl(
\mathrm{H}^{1}(\overline{\mathrm{X}},\mathbf{Z})
\to\mathrm{H}^{1}(\mathrm{X},\mathbf{Z})\bigr)
\]
\[
\mathrm{W}^{0}(\mathrm{H}^{1}(\mathrm{X},\mathbf{Z}))
=\operatorname{Ker}\bigl(
\mathrm{H}^{1}(\overline{\mathrm{X}},\mathbf{Z})
\to\mathrm{H}^{1}(\overline{\mathrm{X}}',\mathbf{Z})\bigr).
\]

(ii) \textit{The spectral sequence defined by the stupid filtration of
\([\mathcal{O}_{\overline{\mathrm{X}}}\to
\bar r_*\Omega^{1}_{\overline{\mathrm{X}}'}(\log\mathrm{S})]\)
degenerates and abuts to the Hodge filtration of
\(\mathbf{H}^{*}(\mathrm{X},\mathbf{C})\).}

Let us prove (i). It suffices to prove the first assertion : the second follows from (8.2.5). Let us compare the exact sequences
\[
\begin{tikzcd}[column sep=1.5em,row sep=3em]
\mathrm{H}^{1}(\overline{\mathrm{X}},\mathbf{Z})
 \arrow[r] \arrow[d]
&
\mathrm{H}^{1}(\mathrm{X},\mathbf{Z})
 \arrow[r,"\partial"] \arrow[d]
&
\mathrm{H}^{2}(\overline{\mathrm{X}}\bmod\mathrm{X},\mathbf{Z})
 \arrow[d,"\wr"]
&
\\
\mathrm{H}^{1}(\overline{\mathrm{X}}',\mathbf{Z})
 \arrow[r]
&
\mathrm{H}^{1}(\mathrm{X}',\mathbf{Z})
 \arrow[r,"\partial"]
&
\mathrm{H}^{2}(\overline{\mathrm{X}}'\bmod\mathrm{X}',\mathbf{Z})
 \arrow[r]
&
\mathrm{H}^{2}(\overline{\mathrm{X}}',\mathbf{Z})
\end{tikzcd}
\]

It follows from (8.3.7) and (3.2.17) that \(\mathrm{H}^{2}(\overline{\mathrm{X}}'\bmod\mathrm{X},\mathbf{Z})\), hence \(\mathrm{H}^{2}(\overline{\mathrm{X}}\bmod\mathrm{X},\mathbf{Z})\) is pure of weight 2. In view of (8.2.4), this proves (i).

To prove (ii), we shall have to return to definition (8.2.1) and use the fact that the simplicial scheme \(((\mathrm{X}'/\mathrm{X})^{\Delta_n})_{n\geq0}\) is smooth (the iterated \((n+1)\)-fold fiber product \((\mathrm{X}'/\mathrm{X})^{\Delta_n}\) is the disjoint union of \(\mathrm{X}'\), embedded diagonally, and a finite number of points) and that it admits \(((\overline{\mathrm{X}}'/\overline{\mathrm{X}})^{\Delta_n})_{n\geq0}\) as a smooth compactification. Let \(\varepsilon:((\overline{\mathrm{X}}'/\overline{\mathrm{X}})^{\Delta_n})_{n\geq0}\to\overline{\mathrm{X}}\) be the augmentation morphism. Let \((\mathrm{K},\mathrm{F})\) be the complex \(\mathcal{O}\to\Omega^{1}(\log\mathrm{S})\) on \(((\overline{\mathrm{X}}'/\overline{\mathrm{X}})^{\Delta_n})_{n\geq0}\), equipped with the stupid filtration. We have an augmentation morphism
\[
\alpha:
[\mathcal{O}_{\overline{\mathrm{X}}}\to
 \bar r_*\Omega^{1}_{\overline{\mathrm{X}}'}(\log\mathrm{S})]
\longrightarrow s\varepsilon_{\bullet *}\mathrm{K}.
\]

The diagram
\[
\begin{tikzcd}[column sep=0.7em,row sep=3em]
\mathrm{H}^{1}(\mathrm{X},\mathbf{C})
 \arrow[r,phantom,"{\textstyle=}" description]
 \arrow[d,equal]
&
\mathbf{H}^{1}\bigl(
 ((\overline{\mathrm{X}}'/\overline{\mathrm{X}})^{\Delta_n})_{\bullet},
 \mathrm{K}\bigr)
&
\mathbf{H}^{1}(\overline{\mathrm{X}},s\varepsilon_{\bullet *}\mathrm{K})
 \arrow[l,"\text{\textcircled{1}}"']
\\
\mathrm{H}^{1}(\mathrm{X},\mathbf{C}) \arrow[rr,equal]
&&
\mathbf{H}^{1}\bigl(\overline{\mathrm{X}},
 [\mathcal{O}_{\overline{\mathrm{X}}}\longrightarrow
 \bar r_*\Omega^{1}_{\overline{\mathrm{X}}'}(\log\mathrm{S})]\bigr)
 \arrow[u]
\end{tikzcd}
\]
is commutative, and 1 is an isomorphism because the \(\varepsilon_n\) are finite, so that \(\mathrm{R}^{i}\varepsilon_{n*}=0\) for \(i>0\). The assertion then follows from the fact that \(\alpha\) is a filtered quasi-isomorphism.

\Page{74}
\textbf{(10.3.12)} Let us complete construction (10.3.8). The diagram
\[
\begin{tikzcd}[column sep=1.5em,row sep=3em]
\mathrm{H}^{1}(\mathrm{X},\mathbf{Z}(1))
 \arrow[r,"\sim"] \arrow[d]
&
\mathrm{T}_{\mathbf{Z}}(\mathrm{H}^{1}_{m}(\mathrm{X})(1))
 \arrow[r]
&
\operatorname{Lie}\operatorname{Pic}(\overline{\mathrm{X}})
 =\mathrm{H}^{1}(\overline{\mathrm{X}},\mathcal{O}_{\overline{\mathrm{X}}})
 \arrow[d,equal]
\\
\mathrm{H}^{1}(\mathrm{X},\mathbf{C})
 \arrow[r,"\sim"]
&
\operatorname{Lie}\operatorname{Pic}^{\natural}(\mathrm{X})
 \arrow[r]
&
\operatorname{Lie}\operatorname{Pic}(\overline{\mathrm{X}})
 =\mathrm{H}_{1}(\overline{\mathrm{X}},\mathcal{O}_{\overline{\mathrm{X}}})
\end{tikzcd}
\]
is commutative. By (10.3.11) (ii) and the definition of the Hodge filtrations, (10.3.10.2) is compatible with the Hodge filtration.

\textit{Corollary} \textbf{(10.3.13)}. — \textit{\(\mathrm{H}^{1}_{m}(\mathrm{X})(1)^{\natural}\) is the extension}
\[
\operatorname{Ker}(\alpha)\longrightarrow\operatorname{Pic}^{\natural}(\mathrm{X})^{0}
\]
\textit{of \(\mathrm{H}^{1}_{m}(\mathrm{X})(1)\) obtained from \textup{(10.3.10.1)}.}

Indeed we have a commutative diagram
\[
\begin{tikzcd}[column sep=2em,row sep=3em]
\operatorname{Ker}(\alpha) \arrow[r]
&
\operatorname{Pic}^{\natural}(\mathrm{X})^{0} \arrow[r]
&
\operatorname{Pic}^{0}(\overline{\mathrm{X}})
\\
\mathrm{H}^{1}(\mathrm{X},\mathbf{Z}) \arrow[r] \arrow[u]
&
\operatorname{Lie}\operatorname{Pic}^{\natural}(\mathrm{X})
 \arrow[r] \arrow[u]
&
\operatorname{Lie}\operatorname{Pic}(\overline{\mathrm{X}}), \arrow[u]
\end{tikzcd}
\]
and we apply (10.1.9).

\textit{Exercise} \textbf{(10.3.14)}. — Verify that the diagram
\[
\begin{tikzcd}[column sep=2em,row sep=3em]
\mathrm{H}^{1}(\mathrm{X},\mathbf{Z}(1))
 \arrow[r,equal] \arrow[d]
&
\mathrm{T}_{\mathbf{Z}}(\mathrm{H}^{1}_{m}(\mathrm{X})(1))
 \arrow[d]
\\
\mathrm{H}^{1}(\mathrm{X},\mu_n) \arrow[r,equal]
&
\mathrm{T}_{\mathbf{Z}/n\mathbf{Z}}(\mathrm{H}^{1}_{m}(\mathrm{X})(1))
\end{tikzcd}
\]
is commutative.

\textit{Remark} \textbf{(10.3.15)}. — The isomorphism deduced from (10.3.8) and (10.3.9) (ii), between \(\mathbf{H}(\overline{\mathrm{X}},\mathcal{O}_{\overline{\mathrm{X}}}\to r_*\Omega^{1}_{\overline{\mathrm{X}}'}(\log\mathrm{S}))\) (which we may call \(\mathrm{H}^{1}_{\mathrm{DR}}(\mathrm{X})\)) and \(\mathrm{T}_{\mathrm{DR}}(\mathrm{H}^{1}_{m}(\mathrm{X})(1))\), equal to \(\operatorname{Lie}\operatorname{Pic}^{\natural}(\mathrm{X})\) by (10.3.13), has been defined in a purely algebraic way.

\Page{75}
\section*{10.4. Translation of a theorem of Picard.}

\textbf{(10.4.1)} Let \(\mathcal{H}\) be a mixed Hodge structure such that \(h^{pq}=0\) for \(p\) or \(q<0\) (resp. for \(p\) or \(q>n\)). I shall provisionally denote by \(\mathrm{I}(\mathcal{H})\) (resp. \(\mathrm{II}_n(\mathcal{H})\)) the 1-motive defined by the largest mixed Hodge substructure of \((\mathcal{H}_{\mathbf{Z}}/\mathrm{torsion})(1)\) (resp. by the largest quotient of \((\mathcal{H}_{\mathbf{Z}}/\mathrm{torsion})(n)\) which is purely of type
\[
\{(-1,-1),(-1,0),(0,-1),(0,0)\}.
\]
If \(\mathrm{X}\) is a complex algebraic variety of dimension \(\leq\mathrm{N}\), I conjecture that the 1-motives \(\mathrm{I}(\mathrm{H}^{n}(\mathrm{X},\mathbf{Z}))\), \(\mathrm{II}_n(\mathrm{H}^{n}(\mathrm{X},\mathbf{Z}))\) (for \(n\leq\mathrm{N}\)) and \(\mathrm{II}_{\mathrm{N}}(\mathrm{H}^{n}(\mathrm{X},\mathbf{Z}))\) (for \(n\geq\mathrm{N}\)) admit a purely algebraic definition. The morphisms
\[
\begin{aligned}
\mathrm{T}_{\ell}(\mathrm{I}(\mathrm{H}^{n}(\mathrm{X},\mathbf{Z})))
 &\longrightarrow
 (\mathrm{H}^{n}(\mathrm{X},\mathbf{Z}_{\ell})/\mathrm{torsion})(1)\\
\mathrm{T}_{\ell}(\mathrm{II}_n(\mathrm{H}^{n}(\mathrm{X},\mathbf{Z})))
 &\longleftarrow
 \mathrm{H}^{n}(\mathrm{X},\mathbf{Z}_{\ell})(n)
 &&(\text{for }n\leq\mathrm{N})\\
\mathrm{T}_{\ell}(\mathrm{II}_{\mathrm{N}}(\mathrm{H}^{n}(\mathrm{X},\mathbf{Z})))
 &\longleftarrow
 \mathrm{H}^{n}(\mathrm{X},\mathbf{Z}_{\ell})(\mathrm{N})
 &&(\text{for }n\geq\mathrm{N})
\end{aligned}
\]
and their analogues in de Rham cohomology should also admit a purely algebraic definition. This is what we have verified for \(\dim(\mathrm{X})=1\). Here is another example.

\textbf{(10.4.2)} Let \(\mathrm{S}\) be a smooth projective surface, \(\mathrm{U}\) a dense (Zariski) open subset of \(\mathrm{S}\) and \(\mathrm{C}=\mathrm{S}-\mathrm{U}\). Put
\[
\mathrm{H}^{\log}_{2}(\mathrm{U},\mathbf{Z})
=\operatorname{Ker}\bigl(
\mathrm{H}_{2}(\mathrm{U},\mathbf{Z})
\to\mathrm{H}_{2}(\mathrm{S},\mathbf{Z})\bigr).
\]
According to [8] \S~9, this group does not depend on the chosen smooth compactification \(\mathrm{S}\) of \(\mathrm{U}\). Modulo finite groups, it is just (3.2.17) \(\mathrm{W}_{-3}\mathrm{H}_{2}(\mathrm{U},\mathbf{Z})\) (\(\mathrm{H}_{2}(\mathrm{U},\mathbf{Q})\)), the dual of \(\mathrm{H}^{2}(\mathrm{U},\mathbf{Q})\), is equipped with the dual mixed Hodge structure). One can verify that the 1-motive defined by \((\mathrm{H}^{\log}_{2}(\mathrm{U},\mathbf{Z})/\mathrm{torsion})(-1)\) admits the following description :

\textit{a)} By Poincaré duality, we have \(\mathrm{H}_{2}(\mathrm{U},\mathbf{Z})=\mathrm{H}^{2}_{c}(\mathrm{U},\mathbf{Z})\). The long exact cohomology sequence of \((\mathrm{S},\mathrm{C})\) gives an exact sequence
\[
\mathrm{H}^{1}(\mathrm{S},\mathbf{Z})
\longrightarrow\mathrm{H}^{1}(\mathrm{C},\mathbf{Z})
\longrightarrow\mathrm{H}^{\log}_{2}(\mathrm{U},\mathbf{Z})
\longrightarrow0.
\]

\textit{b)} Let \(\mathrm{H}^{1}_{m}(\mathrm{C})\) be the quotient of \(\operatorname{Pic}^{0}(\mathrm{C})\) by its unipotent radical. We know that \(\mathrm{H}^{1}(\mathrm{C},\mathbf{Z})=\mathrm{T}_{\mathbf{Z}}(\mathrm{H}^{1}_{m}(\mathrm{C}))\) and \(\mathrm{H}^{1}(\mathrm{S},\mathbf{Z})=\mathrm{T}_{\mathbf{Z}}(\operatorname{Pic}^{0}(\mathrm{S}))\). We therefore have a map
\[
\mathrm{H}^{\log}_{2}(\mathrm{U},\mathbf{Z})
\longrightarrow
\mathrm{T}_{\mathbf{Z}}(\mathrm{H}^{1}_{m}(\mathrm{C})/\operatorname{Pic}^{0}(\mathrm{S})),
\]
which is an isomorphism modulo finite groups.

\textit{c)} The map obtained from \textit{b)}
\[
(\mathrm{H}^{\log}_{2}(\mathrm{U},\mathbf{Z})/\mathrm{torsion})(-1)
\longrightarrow
\mathrm{T}_{\mathbf{Z}}(\mathrm{H}^{1}_{m}(\mathrm{C})/\operatorname{Pic}^{0}(\mathrm{S}))
\tag*{\textbf{(10.4.2.1)}}
\]
is a morphism of mixed Hodge structures (and an isomorphism modulo finite groups).

\Page{76}
One can define in a purely algebraic way the isomorphisms (modulo finite groups) deduced from (10.4.2.1) between the analogues of \(\mathrm{H}^{\log}_{2}(\mathrm{U},\mathbf{Z})(-1)\) in \(\ell\)-adic or de Rham cohomology, and the \(\ell\)-adic or de Rham realizations of \(\mathrm{H}^{1}_{m}(\mathrm{C})/\operatorname{Pic}^{0}(\mathrm{S})\).

\textbf{(10.4.3)} We shall not prove the preceding constructions and compatibilities. The result is implicit in the following theorem of E. Picard ([11], volume I, p. 62). Here \(\mathrm{P}\) and \(\mathrm{A}\) denote polynomials.

\textit{Every residue of the double integral}
\[
\iint\frac{\mathrm{P}(x,y)\,dx\,dy}{\mathrm{A}(x,y)}
\]
\textit{can be regarded as a logarithmic or cyclic period of an abelian integral.}

I recall that :

\textit{a)} A \textit{\(k\)-fold integral} on a nonsingular projective variety \(\mathrm{V}\) of dimension \(d\) is a closed rational \(k\)-form on \(\mathrm{V}\) : it is an element \(\alpha\in\mathrm{H}^{0}(\mathrm{U},\Omega^{k})\), where \(\mathrm{U}\) is a dense open subset of \(\mathrm{V}\), satisfying \(d\alpha=0\). Here, \(k=2\) and \(\mathrm{V}=\mathbf{P}^{2}\), so that \(k=\dim\mathrm{V}\) and the condition \(d\alpha=0\) is vacuous.

\textit{b)} A \(k\)-fold integral \(\alpha\) defines an element of \(\mathrm{H}^{k}(\mathrm{U},\mathbf{C})\) again denoted by \(\alpha\). The \textit{periods} of \(\alpha\) are the numbers \(\langle\alpha,c\rangle\) for \(c\in\mathrm{H}_{k}(\mathrm{U},\mathbf{Z})\). If \(c\) is defined by a singular chain \(\mathrm{C}\), we have
\[
\langle\alpha,c\rangle=\int_{\mathrm{C}}\alpha.
\]

\textit{c)} The \textit{residues} of a \(k\)-fold integral \(\alpha\in\mathrm{H}^{0}(\mathrm{U},\Omega^{k})\) are the numbers \(\langle\alpha,c\rangle\) for
\[
c\in\operatorname{Ker}\bigl(
\mathrm{H}_{k}(\mathrm{U},\mathbf{Z})
\to\mathrm{H}_{k}(\mathrm{V},\mathbf{Z})\bigr).
\]

In the case considered by Picard, the map \(\mathrm{H}_{2}(\mathrm{U},\mathbf{Z})\to\mathrm{H}_{2}(\mathbf{P}^{2},\mathbf{Z})\) is zero and there is no need to distinguish between periods and residues. For the general case, cf. Lefschetz, [10], note I.

\textit{d)} By “logarithmic or cyclic period of an abelian integral”, Picard simply means “period of a single integral \(\beta\) on a curve \(\mathrm{Y}\)”. It is understood that \(\beta\) and \(\mathrm{Y}\) are rationally determined by \(\mathrm{P}\) and \(\mathrm{A}\).

\textbf{(10.4.4)} Let us translate this theorem. Thus let \(\mathrm{S}\) and \(\mathrm{U}\) be as in (10.4.2) (for example \(\mathrm{S}=\mathbf{P}^{2}(\mathbf{C})\)), let \(\beta\in\mathrm{H}^{0}(\mathrm{U},\Omega^{2}_{\mathrm{U}})\) be a double integral and \(c\in\mathrm{H}^{\log}_{2}(\mathrm{U},\mathbf{Z})\) a cycle. By definition
\[
\langle\beta,c\rangle=\int_c\beta
\]
is a residue of \(\beta\). Let \(\mathrm{A}=\mathrm{H}^{1}_{m}(\mathrm{C})(1)/\operatorname{Pic}^{0}(\mathrm{S})\) (an extension of an abelian variety by a torus). The 2-form \(\beta\) defines a cohomology class \(\operatorname{cl}(\beta)\in\mathrm{H}^{2}_{\mathrm{DR}}(\mathrm{U})\) and induces
\Page{77}
a linear form \(\beta'\) on \(\mathrm{H}^{\log}_{\mathrm{DR}2}(\mathrm{U})=\mathrm{T}_{\mathrm{DR}}(\mathrm{A})\). The correspondence \(\beta\mapsto\beta'\) is rational. The class \(c\) defines \(c'\in\mathrm{T}_{\mathbf{Z}}(\mathrm{A})\otimes\mathbf{Q}\), and we have
\[
\langle\beta,c\rangle=2\pi i\langle\beta',c'\rangle.
\tag*{\textbf{(10.4.4.1)}}
\]
We can represent \(c'\) by a cycle on \(\mathrm{C}\), and \(\beta'\) as a single integral on \(\mathrm{C}\).

{\small
\begin{center}
\textbf{BIBLIOGRAPHY}
\end{center}

\noindent [1] A. \textsc{Blanchard}, On complex analytic varieties, \textit{Ann. Sc. E.N.S.}, \textbf{73} (1956).

\noindent [2] T. \textsc{Bloom} and M. \textsc{Herrera}, De Rham cohomology of an analytic space, \textit{Inv. Math.}, \textbf{7} (1969), 275-296.

\noindent [3] P. \textsc{Deligne}, Lefschetz's theorem and criteria for degeneration of spectral sequences, \textit{Publ. Math. I.H.E.S.}, \textbf{35} (1968), 107-126.

\noindent [4] P. \textsc{Deligne}, Hodge theory I, \textit{Proceedings of the International Congress of Mathematicians (Nice, 1970)}, Gauthier-Villars, 1971, \textbf{1}, 425-430.

\noindent [5] P. \textsc{Deligne}, Hodge theory II, \textit{Publ. Math. I.H.E.S.}, \textbf{40} (1971), 5-58.

\noindent [6] P. A. \textsc{Griffiths}, Periods of integrals on algebraic manifolds. Summary of main results and discussion of open problems, \textit{Bull. Am. Math. Soc.}, \textbf{76} (1970), 228-296.

\noindent [7] P. A. \textsc{Griffiths}, On the periods of certain rational integrals : I, \textit{Annals of Math.}, \textbf{90} (1969), 460-495.

\noindent [8] A. \textsc{Grothendieck}, The Brauer group III : Examples and complements, in \textit{Ten lectures on the cohomology of schemes}, North Holland Publ. Co., 88-188.

\noindent [9] A. \textsc{Grothendieck}, On the De Rham cohomology of algebraic varieties, \textit{Publ. Math. I.H.E.S.}, \textbf{29} (1966), 96-103.

\noindent [10] S. \textsc{Lefschetz}, \textit{Analysis situs and algebraic geometry}, Gauthier-Villars, 1924 (reprinted in \textit{Selected Papers}, Chelsea Publ. Co., 1971).

\noindent [11] E. \textsc{Picard} and G. \textsc{Simart}, \textit{Theory of algebraic functions of two independent variables}, vol. I and II, Paris, Gauthier-Villars, 1895, 1906.

\noindent [12] D. G. \textsc{Quillen}, \textit{Notes on the homology of commutative rings}, M.I.T., 1968.

\noindent [13] B. \textsc{Saint-Donat} (based on notes by P. \textsc{Deligne}), Techniques of cohomological descent, S.G.A. \textbf{4} V \textit{bis} (to appear with S.G.A. \textbf{4}, in the \textit{Lecture Notes}, Springer-Verlag).

\noindent [14] G. \textsc{Segal}, Classifying spaces and spectral sequences, \textit{Publ. Math. I.H.E.S.}, \textbf{34} (1968), 105-112.

\noindent [15] J.-P. \textsc{Serre}, Algebraic geometry and analytic geometry, \textit{Ann. Inst. Fourier}, \textbf{6} (1966), cited as G.A.G.A.

\begin{center}
\textbf{ABBREVIATIONS}
\end{center}

\noindent S.G.A., Algebraic Geometry Seminar at Bois-Marie, \textbf{1} to \textbf{7}, distributed by the I.H.E.S.; \textbf{1}, \textbf{3}, \textbf{4}, \textbf{6}, \textbf{7} I have appeared in the \textit{Springer Lecture Notes}, nos. \textbf{224}; \textbf{151}, \textbf{152} and \textbf{153}; \textbf{269}, \textbf{270}, \textbf{305}; \textbf{225}, \textbf{288} — S.G.A. \textbf{2} was published by North Holland Publ. Co.
\par
}

\begin{flushright}
\textit{Manuscript received on 2 March 1972.}
\end{flushright}

\end{document}
